\documentclass[a4paper,12pt,oneside]{article}  % 13pt는 직접 설정함
\special{papersize=210mm,297mm}

% --- Language and Font Settings ---
\usepackage[english]{babel}
\usepackage{fontspec}
\usepackage{kotex}  % Korean typesetting

% --- Font Configuration (Main + Korean) ---
\usepackage{libertinus}
\setmainhangulfont[
  Path = ./,
  UprightFont = *Batang Light.ttf,
  BoldFont    = *Batang Medium.ttf
]{KoPubWorld}

% --- Page Geometry: A4 with wide margins for print readability ---
\usepackage[A4paper, margin=1in]{geometry}

% --- Line Spacing ---
\usepackage{setspace}
\setstretch{1.45}

% --- Section Formatting: Disable numbering ---
\usepackage{titlesec}
\setcounter{secnumdepth}{0}
\setcounter{tocdepth}{3}

\usepackage{chngcntr}
\counterwithin*{footnote}{section}

\titleformat{\chapter}[display]
  {\normalfont\Huge\bfseries}
  {}{0pt}{\Huge}\titleformat{\section}
  {\normalfont\Large\bfseries}
  {}{0pt}{\Large}
\titleformat{\subsection}
  {\normalfont\large\bfseries}
  {}{0pt}{\large}

\usepackage{etoolbox}
\pretocmd{\section}{\clearpage}{}{}

% --- Math and Tables ---
\usepackage{amsmath, amssymb, amsthm, mathtools}
\usepackage{graphicx}
\usepackage{enumitem}
\usepackage{tabularx, booktabs}
\usepackage{footmisc}

% --- Quotes, Hyperlink, Header/Footer ---
\usepackage{csquotes}
\usepackage[hidelinks]{hyperref}
\usepackage{fancyhdr}
\pagestyle{fancy}
\fancyhf{}
\fancyfoot[C]{\thepage}
\setlength{\headheight}{15pt}

% --- Bibliography (biblatex) ---
\usepackage[style=verbose-note, backend=biber, maxbibnames=99]{biblatex}
\addbibresource{references.bib}
\AtEveryBibitem{\clearfield{pages}}

% --- Title Info ---
\title{\Huge\textsc{사법심사에 반대하는 핵심 논거} \\[2ex] \Large The Core of the Case Against Judicial Review}
\author{\Large JEREMY WALDRON}
\date{}

\begin{document}

% --- Title Page ---
\begin{titlepage}
  \centering
  \vspace*{3cm}
  {\Huge\textsc{사법심사에 반대하는 핵심 논거}}\\[1.5ex]
  {\textsc{제레미 월드런 (JEREMY WALDRON)}\\[6ex]
  {\small 2025 소논문 특강 강의용 한국어판\\
  강좌 외의 사용을 불허함}
  \vfill
\end{titlepage}

% ==========================
% 📚 본문 시작
% ==========================
\begin{document}

\frontmatter
\maketitle

\mainmatter
\cleardoublepage
\markboth{}{}

\textbf{초록(ABSTRACT).} 이 에세이는 입법에 대한 사법심사(judicial
review of legislation)에 반대하는 일반적인 논거를 명확히 제시하며, 특정
판결이나 헌법 제도 내에서의 역사적 기원에 대한 논의로 복잡하게 얽히지
않도록 구성되어 있다. 본문은 두 가지 주요 이유를 들어 사법심사를
비판한다. 첫째, 사법심사가 민주적으로 구성된 입법부보다 권리를 더 잘
보호한다고 가정할 이유가 없다는 점을 논증한다. 둘째, 그 결과와는
무관하게, 사법심사는 민주적으로 정당하지 않다는 점을 논증한다. 두 번째
논증은 비교적 익숙한 반면, 첫 번째 논증은 덜 알려져 있다.

그러나 사법심사에 대한 반대는 절대적이거나 무조건적인 것이 아니다. 이
글에서 제시되는 주장은 몇 가지 조건들을 전제로 한다. 예컨대, 해당 사회가
잘 작동하는 민주적 제도를 갖추고 있으며, 대다수 시민들이 권리를 진지하게
받아들이고 있다는 전제이다(비록 자신들이 어떤 권리를 갖고 있는지에
대해서는 의견이 다를 수 있더라도). 이 글은 이러한 조건들이 충족되지 않을
경우 어떤 결과가 따를 수 있는지를 고찰하며 마무리된다.

\textbf{저자 소개(AUTHOR).} 컬럼비아 대학교 로스쿨 석좌교수. (2006년
7월부터 뉴욕대학교(NYU) 로스쿨 교수로 재직.) 이 에세이의 초기 버전은
런던대학교 유니버시티칼리지(UCL)의 법철학 콜로키움, 예루살렘
히브리대학교의 법학 워크숍, 그리고 하버드 로스쿨에서 열린 헌법학
컨퍼런스에서 발표되었다. 이 자리에서 형식적인 논평을 제공해 준 로널드
드워킨, 루스 가비슨, 시어나 쉬프린에게 특별히 감사를 전한다. 또한 제임스
앨런, 아하론 바라크, 리처드 벨라미, 아일린 캐버너, 아서 차스칼슨, 마이클
도르프, 리처드 팰런, 찰스 프라이드, 앤드루 게디스, 스티븐 게스트, 이언
헤이니-로페즈, 알론 하렐, 데이비드 헤이드, 샘 이샤로프, 엘레나 케이건,
케네스 키스, 마이클 클라만, 존 매닝, 안드레이 마르모르, 프랭크 미첼먼,
헨리 모나한, 베로니크 무뇨스-다르데, 존 몰리, 매튜 팔머, 리처드 필데스,
조셉 라즈, 캐롤 생어, 데이비드 위긴스, 조 울프 등 여러 학자들의 제안과
비판에 감사한다. 이 문제에 대해 수년에 걸쳐 나와 논쟁해 온 수백 명의
이들에게도, 이 글을 우정과 감사를 담아 바친다.

\tableofcontents
\clearpage

\newpage

\subsection{서론 (INTRODUCTION)}\label{uxc11cuxb860-introduction}

\textbf{(p.~1348)} 판사는, 특정 법률이 개인의 권리(individual rights)를
침해한다고 확신할 경우, 그 법률을 무효화(strike down)할 권한을 가져야
하는가? 많은 나라에서 판사는 그러한 권한을 갖고 있다. 가장 잘 알려진
사례는 미국이다. 2003년 11월, 매사추세츠주 대법원(Supreme Judicial Court
of Massachusetts)은 해당 주의 혼인 허가법(marriage licensing laws)이
실질적으로 혼인을 남성과 여성 간의 결합으로 한정함으로써, 절차적
적법성(due process)과 평등 보호(equal protection)에 대한 주(州) 헌법상
권리를 침해하였다고 판결하였다.\footnote{Goodridge v. Dep't of Pub.
  Health, 798 N.E.2d 941 (Mass. 2003).} 이 결정은, 자신들의 권리가
법적으로 인정받지 못했고, 동성애자(gay men and women)로서 기존의 혼인법
하에서 이등 시민(second-class citizens)처럼 취급받아 왔다고 느끼던 많은
사람들에게 용기를 주었다.\footnote{See \emph{Landmark Ruling: The
  Victors}, \emph{Boston Herald}, Nov.~19, 2003, at 5.} 비록 이 결정이
훗날 주 헌법 개정을 통해 번복되더라도, 원고(plaintiffs)와 그 지지자들은
적어도 이제 권리에 관한 문제가 직접적으로 논의되고 있다는 점에서 위안을
받을 수 있다. 합리적인 판단(good decision)과 권리 청구(claims of
rights)가 지속적이고 진지하게 고려되는 절차(process)는\footnote{이
  표현은 Ronald Dworkin의 표현을 변형한 것이다. 참조: \emph{Ronald
  Dworkin, A Matter of Principle} 9--32 (1985).}---많은 사람들에게
있어---사법심사(judicial review)라는 제도를 소중히 여기는 이유가 된다.
이들은 사법심사가 때때로 잘못된 판단, 예컨대 로크너 시대(Lochner era)
동안 주 및 연방 법원이 170건의 노동 관련 법률을 무효화한
것처럼,\footnote{\emph{Lochner v. New York}, 198 U.S. 45 (1905).
  1880--1935년 사이에 노동 관계 및 노동 조건에 관한 주 및 연방 법률이
  위헌으로 판결된 전체 사례 수는 다음 문헌에 수록된 목록을 기반으로
  계산되었다: \emph{William E. Forbath, Law and the Shaping of the
  American Labor Movement} apps. A, C, at 177--92, 199--203 (1991).}
부정적 결과로 이어질 수 있음을 인정한다. 또한 이 제도가 어떤 형태로든
민주적 결손(democratic deficit)을 겪고 있다는 점도 인정한다. 그러나
이들은 그러한 비용(costs)이 종종 과장되거나 잘못 묘사된다고 주장한다.
민주적 절차(democratic process) 역시 결코 완벽하지 않으며, 어차피
사법심사에 대한 민주주의적 반대 역시 다수의 전제(tyranny of the
majority)가 문제될 때에는 의문시될 수 있다. 이들은 논증하기를, 우리는
때때로 발생하는 부정적 결과를 감수할 수 있는데, 왜냐하면 이 제도는
로렌스(Lawrence), 로 대 웨이드(Roe), 브라운 대 교육위원회(Brown) 같은
판결들을 통해,\footnote{\emph{Lawrence v. Texas}, 539 U.S. 558 (2003);
  \emph{Roe v. Wade}, 410 U.S. 113 (1973); \emph{Brown v. Bd. of Educ.},
  347 U.S. 483 (1954).} 편견에 사로잡힌 다수 대중에 맞서 우리 사회가
개인의 권리(individual rights)에 전념하고 있다는 점을 확인시켜 주었기
때문이라는 것이다.

이것이 내가 사법심사(judicial review)에 대해 긍정적으로 말할 수 있는
거의 마지막 말이 될 것이다. (나는 이 제도가 남긴 많은 판례의 가치와,
관련된 절차적 쟁점들이 얼마나 복잡한지를 먼저 인정하고 싶었다.) 이
에세이가 논증하려는 바는, 입법에 대한 사법심사는 자유롭고 민주적인
사회에서 최종적인 결정 방식(mode of final decisionmaking)으로서 적절하지
않다는 것이다.

\textbf{(p.~1349)} 이와 같은 결과를 도출하는 논변들은 과거에도 자주
제기되어 왔으며, 이와 같은 제도적 실천을 논할 때 자연스럽게 떠오를
수밖에 없다. 자유주의 정치이론(liberal political theory)에서
입법우위(legislative supremacy)는 흔히 국민 자치(popular
self-government)와 연관된다.\footnote{이 개념의 고전적 정립은 다음에
  있다: \emph{John Locke, The Second Treatise of Government}, in
  \emph{Two Treatises of Government} 265, 366--67 (Peter Laslett ed.,
  Cambridge Univ. Press 1988) (1690).} 그리고 민주주의적 이상(democratic
ideals)은, 선출되지 않은 판사들이 선출된 입법부의 권한을 승인 여부에
따라 제한할 수 있다고 주장하는 어떤 실천과도 본질적으로 긴장 관계에
놓인다. 알렉산더 비클(Alexander Bickel)은 이 쟁점을 다음과 같은 유명한
표현으로 요약했다: ``반(反)다수결의 어려움''(the counter-majoritarian
difficulty).\footnote{\emph{Alexander M. Bickel, The Least Dangerous
  Branch} 16--17 (2d ed.~1986). → ``사법심사는 우리의 체계 내에서
  다수결에 반하는(counter-majoritarian) 힘이다. \ldots{} 대법원이 입법
  행위를 위헌이라고 선언할 때, 그것은 현재 여기에 존재하는 실제 인민의
  대표들의 의지를 좌절시키는 것이다.''} 비클은 이 어려움을 완화할 수
있는 방법으로, 기존의 입법 절차가 대중 혹은 다수 의지를 완전하게
대표하지 못한다는 점을 보여주는 전략을 제안했다. 그러나 그는 다음과 같이
덧붙였다:

\begin{quote}
현대 정치학(modern political science)이 놀라울 정도로 정교하고 열정적인
노력으로 탐구해낸 제도적 복잡성과 난점들---그리고 그 일부는 단순한
발견의 열의를 넘어서는 풍부함으로 오히려 복잡성을 증폭시키는 경향을
띠기도 하지만---그 어떤 복잡성도 미국 민주주의(American democracy)에서
사법심사(judicial review)가 일탈적 제도(deviant institution)라는 본질적
현실을 바꿀 수는 없다.\footnote{Id. at 17--18.}
\end{quote}

사법심사(judicial review)를 통해 입법이 무효화될 수 없는 국가들에서는,
낙태(abortion), 소수집단 우대정책(affirmative action), 학교 바우처
제도(school vouchers), 동성혼(gay marriage)과 같은 사안에 대해 국민
스스로가 통상의 입법 절차(ordinary legislative procedures)를 통해
최종적으로 결정할 수 있다. 또한, 공공장소에서의 인종적 증오 표현(racial
hatred)의 처벌이나 선거에서 후보자 지출 제한과 같은 문제들에 대해서도
국민이 스스로 결정할 수 있다. 만일 이러한 문제들에 대해 국민 사이에 의견
차이가 있다면, 그들은 대표자(representatives)를 선출하여 의회에서
토론(deliberation)과 표결(voting)을 통해 해당 쟁점을 해결할 수 있다.
예컨대 1960년대 영국에서는 낙태법의 자유화, 동의에 의한 성인 간 동성애
행위의 비범죄화, 사형제도의 폐지를 둘러싼 논의가 의회에서 진행되었고,
이는 모두 이러한 방식으로 처리되었다.\footnote{\emph{Abortion Act,
  1967}, c.~87; \emph{Sexual Offences Act, 1967}, c.~60; \emph{Murder
  (Abolition of Death Penalty) Act, 1965}, c.~71.} 각각의 사안에서,
폭넓은 공적 논의(public deliberation)가 하원(House of Commons) 내에서의
진지한 토론으로 반영되었다. 이러한 토론의 질은 (캐나다, 호주, 뉴질랜드
및 기타 국가에서도 유사한 토론이 이루어졌으며), 입법자들이 이러한
사안들을 책임감 있게 다루지 못한다는 주장을 무의미하게 만들며, 이와 같은
입법 결과들의 \textbf{(p.~1350)} 자유주의적 성격(libera outcomes)은 다수
대중(popular majorities)이 소수자(minorities)의 권리(rights)를 지지하지
않을 것이라는 익숙한 명제에 의문을 제기한다.

반면, 미국에서는 주 및 연방 차원의 입법부(state and federal
legislatures)에서 국민 또는 그 대표자들이 위의 쟁점들에 대해 논의할 수는
있지만, 그 결정이 유지될 것이라는 보장이 없다. 어떤 입법적 결론에
동의하지 않는 사람이 그것을 법정에 제소하기로 결정하면, 최종적으로
승리하는 견해는 판사(judges)들의 견해가 된다. 사법심사의 옹호자인 로널드
드워킨(Ronald Dworkin)의 표현을 빌리자면, ``수 세기 동안 철학자, 정치가,
시민들이 논의해온 난해하고, 논쟁적이며, 심오한 정치적 도덕성(political
morality)의 문제들''에 대해, 국민과 그 대표자들은 단지 ``대법관 다수의
판단(deliverances of a majority of the justices)을 받아들여야만'' 하며,
그들의 통찰력이 이와 같은 중대한 문제들에 대해 특별히 탁월하다고 할 수는
없다.\footnote{\emph{Ronald Dworkin, Freedom's Law: The Moral Reading of
  the American Constitution} 74 (1996).}

최근 수년간 미국 내에서는 사법심사를 비판하는 책들이 다수
출간되었다.\footnote{예: \emph{Larry D. Kramer, The People Themselves:
  Popular Constitutionalism and Judicial Review} (2004); \emph{Mark
  Tushnet, Taking the Constitution Away from the Courts} (1999).}
오랫동안 이 제도에 대한 지지는 자유주의자(liberals)들로부터 나왔고,
반대는 자유주의 법원이 옹호한 권리(rights)에 반대하는
보수주의자(conservatives)들로부터 나왔다. 그러나 최근에는
렌퀴스트(Rehnquist) 대법원이 자유주의 입법정책의 주요 성과들 중 일부를
무효화함에 따라, 사법심사에 대한 자유주의 진영의 반대(liberal
opposition)도 증가해 왔다.\footnote{예: \emph{United States v.
  Morrison}, 529 U.S. 598 (2000) (여성폭력방지법의 일부 조항을 위헌으로
  판시); \emph{United States v. Lopez}, 514 U.S. 549 (1995) (의회가 학교
  주변 총기 소지를 금지할 권한이 없다고 판시). 또한 참조: \emph{Mark
  Tushnet, Alarmism Versus Moderation in Responding to the Rehnquist
  Court}, 78 \emph{Ind. L.J.} 47 (2003).} 물론 이 제도를 열정적으로
옹호하는 방어(defenses)도 있었다.\footnote{예: \emph{Dworkin}, supra
  note 10; \emph{Christopher L. Eisgruber, Constitutional Self
  Government} (2001); \emph{Lawrence G. Sager, Justice in Plainclothes:
  A Theory of American Constitutional Practice} (2004).} Marbury v.
Madison 판결 200주년은 이 제도의 기원과 원초적 정당성(original
legitimacy)을 둘러싼 논의들을 촉발시켰으며, Brown v. Board of Education
판결 50주년은 20세기 중반, 인종분리(segregation)와 인종차별적 법률에
맞선 투쟁을 주도했던 미국 법원(courts)의 역할을 되새기게 했다.

이처럼 논쟁의 전선은 형성되었고, 양 진영의 전략은 익숙하며, 입장 역시
널리 이해되고 있다. 그렇다면 필자가 이 글에서 개입(intervention)을
시도하는 \textbf{(p.~1351)} 이유는 무엇인가? 나 자신도 이 문제에 관해
이미 상당히 많이 글을 써 왔다.\footnote{예: \emph{Jeremy Waldron, Law
  and Disagreement} 10--17, 211--312 (1999); \emph{Jeremy Waldron,
  Deliberation, Disagreement, and Voting}, in \emph{Deliberative
  Democracy and Human Rights} 210 (Harold Hongju Koh \& Ronald C. Slye
  eds., 1999) {[}이하 ``Waldron, Deliberation, Disagreement, and
  Voting''{]}; \emph{Jeremy Waldron, Judicial Power and Popular
  Sovereignty}, in \emph{Marbury Versus Madison: Documents and
  Commentary} 181 (Mark A. Graber \& Michael Perhac eds., 2002) {[}이하
  ``Waldron, Judicial Power and Popular Sovereignty''{]}; \emph{Jeremy
  Waldron, A Right-Based Critique of Constitutional Rights}, 13
  \emph{Oxford J. Legal Stud.} 18 (1993) {[}이하 ``Waldron, A
  Right-Based Critique''{]}.} 사법심사를 비판하는 또 하나의 글을 굳이
쓰는 이유는 무엇인가?

내가 하고자 하는 바는, 사법심사에 반대하는 핵심 논거(a core argument
against judicial review)를 식별하는 것이다. 이 논거는 사법심사의 역사적
형태(historical manifestations)나 그 구체적 효과(particular
effects)---즉, 그것이 낳은 좋은 판결과 나쁜 판결, 그로 인해 발생한
고통과 긍정 등---에 종속되지 않는 것이다. 나는 판사들이 권한을 어떻게
행사하는지, 그리고 그들이 법률을 심사할 때 어떤 태도(신중함 또는
적극성)로 접근하는지와는 무관하게, 사법심사 자체에 대한 반대 논거에
집중하고자 한다. Mark Tushnet과 Larry Kramer의 최근 저작들은 이 제도에
대한 이론적 비판을 사법심사의 역사적 기원과 미국 헌법이 보다
비사법화된(less judicialized) 방식으로 구성될 경우 어떤 모습일지를
설명하는 논의와 뒤섞고 있다.\footnote{\emph{Kramer}, supra note 11;
  \emph{Tushnet}, supra note 11.} 이것은 Tushnet과 Kramer에 대한 비판이
아니다. 그들의 저작은, 이론적 논쟁에 풍부한 역사성과 색채를 부여한다는
점에서 가치가 있다. Frank Michelman이 《The People Themselves》의 책
뒷면 추천문에서 말하듯이, ``Kramer의 역사 서술은 사법심사와 대중
헌법주의(popular constitutionalism)를 둘러싼 논쟁에 생명을
불어넣는다.''\footnote{Frank Michelman, 위 Kramer 저서에 대한 책 커버
  발언.} 그리고 실제로 그렇다. 하지만 나는 그 ``살(flesh)''을 일부
걷어내고, 규범적 논거(normative argument)를 그 ``뼈대(bare
bones)''로까지 끌어내려, 사법심사(judicial review)가 실제로 어떤
전제(premise)에 기초하고 있는지를 바로 들여다보고자 한다.

찰스 블랙(Charles Black)은 한때, 실제 현실에서 사법심사(judicial
review)에 대한 반대는 대체로 ``간헐적인 반응(a sometime thing)''에
그친다고 지적한 바 있다. 사람들은 자신이 소중히 여기는 몇몇
판례---예컨대 \emph{Brown} 또는 \emph{Roe}---에 대해서는 사법심사를
지지하고, 유감스러운 결과가 나왔을 때에만 이를 반대하는 경향이 있다는
것이다.\footnote{\emph{Charles L. Black, Jr., A New Birth of Freedom:
  Human Rights, Named and Unnamed} 109 (1997).} 정치 영역에서는
사법심사에 대한 지지가 특정 판결에 대한 지지와 강하게 얽히는 경우가
많다. 이는 특히 낙태권(abortion rights)을 둘러싼 논쟁에서 두드러지는데,
pro-choice 진영에서는 사법심사에 대한 비판을 고려하는 것조차 두려워하는
경향이 있다. 그 이유는 그러한 비판이 \emph{Roe v. Wade} 판결을 보수적
입법자들의 권리를 침해한 정당하지 못한 간섭(unwarranted intrusion)으로
간주하는 이들에게 위안을 주고 힘을 실어줄 것이라는 우려 때문이다. 나는
사법심사의 개별적 결과로부터 추상화된 방식으로 그 핵심적인 비판
논거(core case against judicial review)를 제시함으로써,
\textbf{(p.~1352)} 이러한 공황상태(panic)의 일부를 극복하는 데 기여할 수
있기를 바란다. 물론 일부 국가에서는 성(sex), 인종(race),
종교(religion)와 관련된 입법의 병리적 현상(legislative
pathologies)으로부터 방어하기 위한 보호 수단(protective measure)으로서
사법심사가 필요할 수도 있다. 그러나 그렇다 하더라도, 그와 같은 옹호가
사법심사 자체의 본질적인 정당성에 해당하는지, 아니면 예외적인 상황에서
일반적으로 설득력 있는 규범적 논거(normative argument)에 반하여
적용되어야 하는 특수한 사정에 불과한 것인지 검토할 필요가 있다.

이론적 비판(theoretical critique)의 ``살(flesh)''을 걷어내고
``뼈대(bones)''만 남기려는 또 하나의 이유는, 사법심사가 미국 이외의 여러
나라에서도 문제로 제기되고 있기 때문이다. 각국은 미국과는 전혀 다른
역사, 사법 문화(judicial culture), 입법 제도에 대한 경험을 가지고 있다.
예컨대 영국에서 판사들이 입법을 심사할 수 있는 권한이 상대적으로
제한되어 있는 문제를 두고 논쟁할 때, 사람들은 1805년
공화주의자(Republicans)들이 연방주의자(Federalists)에게 어떤 주장을
했는지, 혹은 \emph{Brown v. Board of Education}의 유산(legacy)에 대해
별다른 관심을 갖지 않는다. 이처럼 문화적, 역사적, 정치적
선입견(preoccupations)으로 오염되지 않은 보다 일반적인 이해(general
understanding)가 필요하다.\footnote{다시 말하자면, 보다 구체적인 설명을
  무시하려는 것은 아니다. 이 논문이 의도하는 바는 이론적 논거를 명확히
  하여, 이를 영국, 미국, 캐나다, 남아프리카공화국 등에서의 실제 논의와
  나란히 놓고 비교하는 것이다.}

나 자신의 저술은 대체로 다소 추상적이었다. 하지만 나는 사법심사를
법철학(jurisprudence)과 정치철학(political philosophy)의 다른 쟁점들과
얽힌 방식으로 논의해 왔다.\footnote{나는 다음 논문에서 개인의
  \emph{권리(right)}라는 개념이 사법심사를 요구하는지를 물었다:
  \emph{Waldron, A Right-Based Critique}, supra note 14. 또한 다음
  문헌들에서 시민공화주의적 사상과의 관계를 고찰하였다: \emph{Jeremy
  Waldron, Judicial Review and Republican Government}, in \emph{That
  Eminent Tribunal: Judicial Supremacy and the Constitution} 159
  (Christopher Wolfe ed., 2004); 벤담주의적 민주주의 개념과 루소주의적
  개념의 차이와 관련해서는 \emph{Jeremy Waldron, Rights and Majorities:
  Rousseau Revisited}, in \emph{Nomos XXXII: Majorities and Minorities}
  44 (John W. Chapman \& Alan Wertheimer eds., 1990) {[}이하 ``Waldron,
  Rights and Majorities''{]}; 대륙 유럽의 인민주권 이론과의 관계는
  \emph{Waldron, Judicial Power and Popular Sovereignty}, supra note 14.
  나는 또한 사법심사 논쟁이 메타윤리학에서의 실재론(realism)과 가치의
  객관성(objectivity)에 관한 논쟁과 어떤 관련을 가지는지 검토하였다:
  \emph{Jeremy Waldron, The Irrelevance of Moral Objectivity}, in
  \emph{Natural Law Theory} 158 (Robert P. George ed., 1992) {[}이하
  ``Waldron, The Irrelevance of Moral Objectivity''{]}; \emph{Jeremy
  Waldron, Moral Truth and Judicial Review}, 43 \emph{Am. J. Juris.} 75
  (1998) {[}이하 ``Waldron, Moral Truth and Judicial Review''{]}. 또한
  여러 형태의 사법심사 옹호론에 대해 비판하였다. 예:
  선결약정(precommitment)에 대한 논의는 \emph{Jeremy Waldron,
  Precommitment and Disagreement}, in \emph{Constitutionalism:
  Philosophical Foundations} 271 (Larry Alexander ed., 1998) {[}이하
  ``Waldron, Precommitment and Disagreement''{]}; Ronald Dworkin이
  \emph{Freedom's Law}에서 제시한 민주주의와의
  \emph{양립가능성(compatibility)} 논증은 \emph{Jeremy Waldron, Judicial
  Review and the Conditions of Democracy}, 6 \emph{J. Pol. Phil.} 335
  (1998)에서 검토하였다.} 그럼에도 불구하고, 나는 아직까지 사법심사에
대한 근본적 반론이 무엇인지 \textbf{(p.~1353)} 명확하고 간결하게
제시하지 못했다고 느끼며, \emph{Law and Disagreement}나 그 외의
저작들에서 내가 제시한 주장에 비판을 제기한 이들에 대해 만족스러운
답변을 제공하지 못했다고 생각한다.

이 에세이에서 나는 사법심사가 두 가지 차원에서 비판에 취약하다는 점을
논증할 것이다. 첫째, 종종 주장되는 바와는 달리, 사법심사는 시민들이
권리에 관해 의견이 다를 때 핵심 쟁점(real issues at stake)에 명확히
집중할 수 있는 수단을 제공하지 않는다. 오히려 그것은 판례(precedent),
법률 문언(texts), 해석(interpretation) 등 부차적인 쟁점(side-issues)들로
시민들의 관심을 분산시킨다. 둘째, 사법심사는 민주주의적 가치(democratic
values)의 관점에서 정치적으로 정당하지 않다(politically illegitimate).
소수의, 선출되지 않았고 책임성(accountability)도 없는 판사들 사이의
다수결(majority voting)에 특권을 부여함으로써, 일반 시민들의
참정권(enfranchisement)을 박탈하고, 대의(representation)와 정치적
평등(political equality)이라는 소중한 민주주의 원칙들을 권리 문제에 대한
최종 판단 과정에서 배제한다.

에세이의 전개는 다음과 같다. 제 I부에서는 내가 비판하고자 하는 대상인
'강한 형태의 사법심사(strong judicial review of legislation)'를
정의하고, 내가 비판의 대상으로 삼지 않는 다른 실천들과 구분할 것이다. 제
II부에서는 내 주장의 전제가 되는 몇 가지 가정을 제시할 것이다. 이
에세이의 주장은 무조건적인 것이 아니라, 현대 자유민주주의(modern liberal
democracies)의 제도적·정치적 특징에 의존한다. 제 III부에서는 전체 주장의
일반적 성격을 간단히 검토할 것이다. 이 주장은 결과(outcome)와
절차(process) 두 측면의 이유 모두를 포함하며, 각각 제 IV부와 제 V부에서
논의된다. 제 VI부에서는 대의제 기관(representative institutions)의
결정을 허용하지 않으려는 가장 흔한 논거---즉 그러한 시스템이 필연적으로
다수의 전제(tyranny of the majority)로 이어진다는 주장---의 오류를
드러낼 것이다. 마지막으로 제 VII부에서는 핵심 논거(core argument)가
의존하는 가정들에서 벗어나야 할 이유가 존재하는 비핵심 사례들(non-core
cases)에 대해 간략히 논할 것이다.

\subsection{I. 사법심사의 정의 (DEFINITION OF JUDICIAL
REVIEW)}\label{i.-uxc0acuxbc95uxc2ecuxc0acuxc758-uxc815uxc758-definition-of-judicial-review}

우선, 내가 '사법심사(judicial review)'라는 개념으로 무엇을 의미하는지를
간략히 설명하고자 한다. 이 에세이는 입법에 대한 사법심사(judicial review
of legislation)에 관한 것이며, 행정부의 행위나 행정적
의사결정(administrative decisionmaking)에 대한 사법심사를 다루는 것이
아니다.\footnote{유럽인권재판소가 수행하는 다수의 사법심사는 행정부의
  행위에 대한 것이다. 그 중 일부는 입법행위에 대한 사법심사이고, 일부는
  사법행위 자체에 대한 사법심사이다. 다음을 참조하라: \emph{Seth F.
  Kreimer, Exploring the Dark Matter of Judicial Review: A
  Constitutional Census of the ???os}, 5 \emph{Wm. \& Mary Bill Rts. J.}
  427, 458--59 (1997). 이 글은 미국 연방대법원의 헌법 판례 다수가
  입법부가 아니라 하급 행정부 관료의 행위를 다루고 있다고 주장한다.}
내가 다루고자 하는 질문은 \textbf{(p.~1354)} 한 정치 공동체(polity)의
선출된 입법부(elected legislature)가 제정한 \textbf{일차적 입법(primary
legislation)}을 대상으로 한다. 이러한 논점 중 일부는 행정부(executive
action)에 대해서도 동일하게 적용된다고 생각할 수도 있다. 왜냐하면 행정부
또한 일정한 선출 정당성(elective credentials)을 지니고 있어서 판사들의
결정에 반대할 수 있기 때문이다. 그러나 거의 모든 곳에서 행정부의 선출
정당성은 법치(rule of law) 원칙에 종속된 것으로 받아들여지고 있으며, 그
결과로 공직자들은 법률에 의거한 정당한 권한(legal authorization)에 따라
행동해야 한다는 법원의 명령을 받을 수 있다.\footnote{Seana Shiffrin,
  Richard Pildes, Frank Michelman 등은 다음과 같이 제안하였다: 본 논문이
  입법에 대한 사법심사에 반대하는 논거가, 오래 전에 제정된 법률이나
  해석을 요하는 법률에 따라 행해지는 행정부의 행위에 대한 사법심사에도
  확대 적용될 수 있는지를 검토하라는 제안. 이에 대해서는 보다 많은
  논의가 필요하다. 이 방향으로 논의를 전개하는 것이 이 논거에 대한
  \emph{귀류법(reductio ad absurdum)}이 될 수도 있고, 혹은 유용한
  응용으로 평가될 수도 있다.} 입법부에 대해서도 유사한 주장, 즉
``사법심사란 곧 입법부를 법치에 종속시키는 것이다''라는 입장이 제기되어
왔다. 그러나 입법부의 경우에는 이 주장이 논쟁의 여지가 없지 않다. 사실
바로 이 점이 우리가 여기서 다루고자 하는 핵심 쟁점이다.

입법에 대한 사법심사(judicial review of legislation)는 전 세계적으로
다양한 방식으로 이루어지고 있으며, 이를 일반적인 범주 아래에 분류할 수
있다. 이러한 실천들은 몇 가지 차원에서 구분될 수 있는데, 가장 중요한
구분은 내가 \textbf{강한 사법심사(strong judicial review)}와
\textbf{약한 사법심사(weak judicial review)}라고 부를 두 가지 유형 간의
차이이다. 이 에세이에서 내가 비판하고자 하는 대상은 바로 \textbf{강한
사법심사(strong judicial review)}이다.\footnote{강한 사법심사(strong
  judicial review)와 약한 사법심사(weak judicial review)의 구분은,
  사법우위(judicial supremacy)라는 개념과는 별개이다. 사법우위란 다음 세
  가지 요소로 구성된다: (1) 법원이 정치 체제 전반의 중요한 쟁점을
  결정하며, (2) 그 결정이 정치 체제 내의 다른 모든 행위자에게 절대적으로
  구속력을 갖고, (3) 법원이 해당 쟁점에 관해 다른 정부 부서의 입장을
  전혀 존중하지 않으며(과거 자의 판례조차 제한된 \emph{선례구속의
  원칙(stare decisis)} 아래에서만 존중한다는 점에서), 독자적으로
  판단하는 체제. 참조: \emph{Barry Friedman, The History of the
  Countermajoritarian Difficulty, Part One: The Road to Judicial
  Supremacy}, 73 \emph{N.Y.U. L. Rev.} 333, 352 \& n.63 (1998);
  \emph{Jeremy Waldron, Judicial Power and Popular Sovereignty}, supra
  note 14, at 191--98.}

강한 사법심사의 체계에서 법원(courts)은 특정 사건에서 해당 법률이 명백히
적용되는 경우에도 그 법률의 적용을 거부하거나, 법률이 명시적으로
예상하지 않는 방식으로 그 적용 효과를 수정하여 개인의 권리(individual
rights)에 부합하도록 만들 수 있는 권한을 가진다. 더 나아가, 이 체계의
법원은 특정 법률이나 조항을 법적으로 적용 불가하다고 선언할 수 있으며,
그 결과 판례구속력(stare decisis)과 쟁점 차단(issue preclusion)의 원칙에
따라 해당 법률은 사실상 사문화(dead letter)된다. 심지어 더 강한 형태의
사법심사라면, 법원이 해당 입법 조항을 법전(statute-book)에서 완전히
제거할 수 있는 권한까지 갖는다. 일부 유럽 국가의 법원들은 이와 같은
\textbf{(p.~1355)} 권한을 실제로 갖고 있다.\footnote{Mauro Cappelletti
  \& John Clarke Adams, Comment, \emph{Judicial Review of Legislation:
  European Antecedents and Adaptations}, 79 \emph{Harv. L. Rev.} 1207,
  1222--23 (1966). 법률이 위헌으로 선언되었을 경우, 그 법률이 제정
  당시부터 무효였다고 간주되는지 여부에 관해 추가적인 논의가 존재한다.}
미국 법원은 그러한 형식적 권한을 가지지는 않지만,\footnote{이 문제는
  명확하지 않다. 위헌으로 선언된 법률이 법전에서 삭제되지 않는다는
  주장에 대한 근거로 다음 사례를 고려할 수 있다: \emph{Dickerson v.
  United States}, 530 U.S. 428 (2000). 이 사건에서 연방대법원은 과반수
  의견으로, 미란다 경고 없이 이루어진 자백도 자발적이면 증거로
  인정된다고 한 연방법(18 U.S.C. § 3501)을 위헌이라고 판시하였다. 이에
  대한 Scalia 대법관의 반대의견 마지막 부분은 위헌으로 판단된 입법이
  여전히 법적 참조 대상으로 남아 있을 수 있음을 시사한다. 그는 다음과
  같이 말했다: ``나는 오늘의 판결에 반대하며, § 3501 조항이 폐지되기
  전까지는 피고인의 자백이 자발적이었다는 지속 가능한 판단이 내려진 모든
  사건에서 이 조항을 계속 적용할 것이다.'' (\emph{Id.} at 464). 반대되는
  인상을 주는 사례로는 \emph{McCorvey v. Hill}, 385 F.3d 846, 849 (5th
  Cir. 2004)가 있다. 이 사건에서 제5순회항소법원은 \emph{Roe v.
  Wade}에서 문제가 된 텍사스주의 낙태법은 암묵적으로 폐지된 것으로
  간주되어야 한다고 판시하였다. 그러나 이 판례를 면밀히 검토하면, 암묵적
  폐지는 \emph{Roe} 판결 자체가 아니라 그 이후 제정된 텍사스 낙태
  규제법들에 의해 이루어진 것으로 보인다. (이 참고는 Carol Sanger에게
  감사를 표한다.)} 실질적인 효과는 그와 크게 다르지 않다.\footnote{Richard
  H. Fallon, Jr., Commentary, \emph{As-Applied and Facial Challenges and
  Third-Party Standing}, 113 \emph{Harv. L. Rev.} 1321, 1339--40 (2000).}

반면 \textbf{약한 사법심사(weak judicial review)} 체계에서는, 법원이
입법이 개인의 권리와 합치하는지 여부를 검토할 수는 있지만, 단지 권리가
침해된다는 이유만으로 그 법률의 적용을 거부하거나 그 적용을 완화할 수는
없다.\footnote{Stephen Gardbaum, \emph{The New Commonwealth Model of
  Constitutionalism}, 49 \emph{Am. J. Comp. L.} 707 (2001).} 그럼에도
불구하고, 이러한 검토는 일정한 효과를 가진다. 예컨대 영국에서는 법원이
어떤 조항이 유럽인권협약(European Convention on Human Rights)에서 보장된
권리(Convention right)와 양립하지 않는다고 판단하는 경우 ``양립불가능성
선언(declaration of incompatibility)''을 내릴 수 있다. 이러한 권리는
인권법(Human Rights Act)을 통해 영국 국내법에 통합되어 있다. 이 법은
그러한 선언이 ``해당 조항의 유효성(validity), 계속적인 효력(continuing
operation) 또는 집행(enforcement)에 영향을 미치지 않으며, 해당 소송
당사자들을 법적으로 구속하지 않는다''고 명시하고 있다.\footnote{\emph{Human
  Rights Act}, 1998, c.~42, § 4(2), (6).} 그럼에도 불구하고, 이 선언은
실제적인 효과를 낳는다. 예컨대 장관(minister)은 이러한 선언을 근거로
신속입법절차(fast-track legislative procedure)를 개시할 수 있으며,
이러한 절차는 바로 그 사법심사 절차 없이는 행사할 수 없는
권한이다.\footnote{\emph{Id.} § 10.}

\textbf{(p.~1356)} 더 약한 형태의 사법심사조차 인정하지 않는 경우도
있다. 예컨대 영국의 경우와 유사하게, 뉴질랜드(New Zealand)의 법원도
인권을 침해하는 법률을 적용하지 않는 결정을 내릴 수는 없다. (여기서의
인권은 1990년 「권리장전법(Bill of Rights Act)」에 명시된 권리를
의미한다.\footnote{\emph{New Zealand Bill of Rights Act 1990}, 1990
  S.N.Z. No.~109, § 4. ``본 법안의 발효 전후를 불문하고 모든
  입법은\ldots{} 본 권리장전(Bill of Rights)의 어떠한 조항과
  불일치한다는 이유만으로, 암묵적으로 폐지 또는 무효화되었거나, 어떤
  방식으로든 무효 또는 효력을 상실한 것으로 판단되어서는 아니 된다.''})
그러나 법원은 그 침해를 회피할 수 있는 해석을 찾아내기 위해 상당히
노력할 수 있다.\footnote{\emph{Id.} § 6. ``어떠한 입법이라도 본
  권리장전에 명시된 권리 및 자유와 양립가능한 의미로 해석될 수 있는
  경우, 그 의미가 다른 모든 의미에 우선한다.''} 뉴질랜드 법원은 경우에
따라 자발적으로 양립불가능성 선언을 내릴 준비가 되어 있다고 밝힌 적도
있지만, 그러한 선언은 입법 과정에 아무런 법적 효과를 미치지
않는다.\footnote{\emph{Moonen v. Film \& Literature Bd. of Review},
  {[}2000{]} 2 N.Z.L.R. 9, 22--23 (C.A.).}

중간적 형태의 사례들도 존재한다. 예컨대 캐나다에서는 입법에 대한
사법심사 제도가 존재하며, 미국 법원과 마찬가지로, 캐나다 법원은 어떤
연방 또는 주 법률이 「권리와 자유에 관한 헌장(Charter of Rights and
Freedoms)」의 규정을 위반한다고 판단할 경우 그 법률의 적용을 거부할 수
있다. 그러나 캐나다의 입법은 형식상 그러한 검토로부터 벗어날 수 있도록
구성될 수 있다. 즉, 캐나다 입법기관은 헌장의 권리 조항들에도 불구하고
법률을 제정할 수 있는 \textbf{``불구조항(notwithstanding clause)''}을
발동할 수 있다.\footnote{\emph{Canadian Charter of Rights and Freedoms},
  \emph{Part I of the Constitution Act, 1982}, being \emph{Schedule B to
  the Canada Act 1982}, ch.~11, § 33(1)--(2) (U.K.). (1) 의회나 주
  의회는 해당 법률 또는 조항이 본 권리장전(Charter)의 제2조 또는
  제7조\textasciitilde 제15조의 조항에도 불구하고 효력을 갖는다고
  명시적으로 선언할 수 있다. (2) 이 조항에 따라 선언이 이루어진 법률
  또는 조항은 해당 선언이 없다면 무효일 조항에도 불구하고 효력을 가진다.}
그러나 실제로 이 불구조항은 거의 사용되지 않는다.\footnote{이 조항이
  실제로 적용된 사례 대부분은 퀘벡 정치 맥락에서 이루어졌다. 참조: Tsvi
  Kahana, \emph{The Notwithstanding Mechanism and Public Discussion:
  Lessons from the Ignored Practice of Section 33 of the Charter}, 44
  \emph{J. Inst. Pub. Admin. Can.} 255 (2001).} 따라서 이후 논의에서는
캐나다의 제도를 \textbf{(p.~1357)} \textbf{강한 사법심사(strong judicial
review)}의 한 형태로 간주하되, 형식상 존재하는 이 예외조항이 내 논거에
미치는 영향은 미미하다고 보겠다.\footnote{Jeffrey Goldsworthy는
  ``예외조항(notwithstanding provision)''이 강한 사법심사(strong
  judicial review)의 민주주의적 정당성에 대해 우려하는 이들에게 충분한
  답이 될 수 있다고 주장하였다. Jeffrey Goldsworthy, \emph{Judicial
  Review, Legislative Override, and Democracy}, 38 \emph{Wake Forest L.
  Rev.} 451, 454--59 (2003). 그는 이 조항이 자주 사용되지 않더라도
  문제되지 않는다고 말한다: ``분명 그것은 유권자의 민주적 특권이며,
  Waldron은 이를 존중해야 할 것이다. 유권자가 헌법적 권리의 외양적
  객관성에 기만당하거나, 사법부의 신비성 또는 판사의 법적 전문성과
  공정성에 대한 순진한 믿음에 현혹될 수 있다는 식으로 반박해서는 안
  된다. 그런 반박은 헌법에 명시된 권리를 지지하는 이들과 똑같이,
  유권자의 계몽된 자기통치 능력에 대한 불신을 드러내는 것이기
  때문이다.'' (\emph{Id.} at 456--57). 그러나 나는 진짜 문제는 이 33
  조항이 입법부로 하여금 권리에 관한 입장을 왜곡하도록 강요한다는 점에
  있다고 본다. 헌장(Charter)을 무시하고 법을 제정한다는 것은, 해당
  권리가 헌장이 부여한 중요성을 갖지 않는다고 말하는 것이나 다름없다.
  그러나 사법부와 입법부 간의 전형적 대립은 헌장 권리를 중요하게 여기는
  판사들과 그렇지 않은 입법자들 간의 충돌이 아니다. 그것은 보통, 모두
  헌장 권리를 중요하게 여기지만 해당 권리를 어떻게 이해할지에 관해
  \emph{의견불일치(disagreement)}하는 입법 다수, 소수와 사법 다수,
  소수들 사이의 문제다. Goldsworthy도 이를 인정한다: ``사법부가\ldots{}
  헌장 조항의 `진정한' 의미와 효과에 대해 입법부와 의견불일치할 것으로
  예상될 때, 입법부는 헌장을 우회하는 것으로 보이지 않고는 자신의 입장을
  관철시킬 수 없다. 그리고 그것은 입법부가 헌법적 권리를 노골적으로
  전복하고 있다는 정치적으로 치명적인 비판에 노출된다.'' (\emph{Id.} at
  467). 그러나 아마도 이 문제를 피할 수 있는 문구는 존재하지 않을
  것이다. 실질적인 정치의 영역에서, 입법부는 항상 사법부가 발표하는
  권리장전 해석의 영향력에 일정 부분 종속된다. (이 점에 대해 John
  Morley에게 감사를 표한다.)}

사법심사(judicial review)의 유형을 구분하는 두 번째 기준은, 특정 사회의
헌법 체계(constitutional system) 내에서 \textbf{개별적 권리(individual
rights)}가 차지하는 위상에 주목한다. 미국에서는 법률(statutes)이
헌법(Constitution)에 명시된 개별적 권리에 부합하는지를 기준으로
심사된다. 권리 지향적 사법심사(rights-oriented judicial review)는
일반적인 헌법심사(constitutional review)의 필수적 요소(part and
parcel)이며, 법원(courts)은 연방주의(federalism)나 권력분립
원칙(separation of powers principles)의 침해를 이유로 법률을 무효화하는
것과 동일한 정신(spirit)으로 권리 침해를 이유로 법률을
무효화한다.\footnote{사법심사를 옹호한 가장 유명한 판례인 \emph{Marbury
  v. Madison}은 개인의 권리와는 무관하였다. 그것은 하원의 평화 판사를
  임명하고 해임하는 의회의 권한에 관한 것이었다.} 이로 인해 미국에서
사법심사를 옹호하는 논변들은 독특한 색채를 띠게 된다. 철학적 옹호는 흔히
\textbf{(p.~1358)} 사법부(judiciary)가 권리에 관한 명제들을 다루는 데에
특히 능숙하다는 점에 초점을 맞추지만, 실제로는 이러한 논변조차도 헌법
규범(constitutional rules)을 유지하기 위한 법원의 구조적 역할(structural
role)을 방어하는 논변에 종속된다. 이 두 가지 방어 논리가 일치하는 경우도
있지만, 서로 분리되기도 한다. 예컨대, \textbf{문언주의(textualism)}는
구조적 사안에 대해 적절한 접근으로 여겨질 수 있지만, 권리 문제를
사유하는 데 있어서는, 설령 그 권리가 권위 있는 문서에 담겨 있다고 해도,
부적절한 기반으로 간주될 수 있다.\footnote{\emph{Dworkin}, supra note 3,
  at 11--18; \emph{Andrei Marmor, Interpretation and Legal Theory}
  156--57 (rev. 2d ed.~2005).} 다른 국가들에서는, 사법심사는 구조적
헌법(constitution)의 일부로 특별히 지정되지 않은 권리장전(bill of
rights)에 근거하여 이루어진다. 영국에서 \textbf{인권법(Human Rights
Act)}에 따라 이루어지는 약한 사법심사(weak judicial review)는 이러한
유형에 속한다. 강한 사법심사(strong judicial review)의 대부분은
헌법심사(constitutional review)와 관련되어 있으므로, 이후 논의는 주로
이들 사례에 초점을 맞출 것이다. 그러나 구조적 제약(structural
constraints)에 중점을 두는 접근이 권리의 문제를 다루는 데 적절하지 않을
수 있다는 점과, 권리 지향적 사법심사(rights-oriented judicial review)에
대한 여러 비판이 다른 형태의 헌법심사에도 제기될 수 있다는 점은 기억할
필요가 있다. 예컨대 최근 몇 년간, 미국 연방대법원(Supreme Court of the
United States)은 연방주의에 대한 자신들의 해석과 충돌한다는 이유로 여러
법률들을 무효화하였다.\footnote{예: supra note 12.} 누구나 동의하듯,
현재의 주(state)와 연방(central government) 간의 관계는, 대부분의 헌법
조문이 비준되었던 18세기 말이나, 연방 구조에 관한 조문이 실질적으로
마지막으로 수정된 19세기 중엽과는 전혀 다른 방식으로 운영되고 있다.
하지만 새로운 국가 구조(state/federal relations)의 기초가 무엇이어야
하는지에 대해서는 의견이 분분하다. 헌법 조문 자체는 이 문제를 해결해주지
않기 때문에, 결국 이 문제는 대법관들 사이의 표결(voting among
Justices)에 의해 결정된다. 일부는 특정한 연방주의 개념을 지지하고 이를
헌법 속에 ``읽어넣으며(read into the Constitution),'' 다른 이들은 또
다른 개념을 지지하며, 다수 의견이 결국 판례를 형성한다. 그러나 이러한
방식이 자유롭고 민주적인 인민 사이의 결사 조건(terms of association)을
결정하는 데 적절한 기반이라고 보기는 어렵다.\footnote{주 및 연방 의회 간
  권한 배분의 구조적 한계를 감독하는 사법심사의 필요성은 종종 권리 기반
  입법 제한을 옹호하는 이들에 의해 (기회주의적으로) 인용된다. 그들은
  말한다: ``입법부는 어차피 연방주의(federalism) 문제로 인해 사법심사의
  대상이다. 그렇다면 그 관행을 권리 기반 사법심사로 확대해도 되지
  않겠는가?'' 본 논문에서 제시된 권리 기반 사법심사의 바람직성에 대한
  분석은 이러한 혼합형 또는 기회주의적 논증에도 적용될 수 있다.}

세 번째 구분 기준은, 미국식의 사후적 사법심사(a posteriori review)와
사전적 사법심사(ex ante review) 사이의 차이에 있다. 미국식 사후적
사법심사는 특정 사건에 대한 법적 절차의 맥락 속에서 이루어지며, 법률이
제정된 이후 상당한 \textbf{(p.~1359)} 시간이 지난 뒤에 이루어지기도
한다. 이에 반해 사전적 사법심사는, 법률안의 제정 말미에서 그 법률안에
대해 추상적 평가(abstract assessment)를 수행하도록 특별히 설계된
헌법재판소(constitutional court)가 수행하는 방식이다.\footnote{1유형
  제도들 중 일부는 제한된 상황에서 사전 자문 의견(ex ante advisory
  opinions)을 제공할 수 있도록 규정하고 있다. 예컨대 매사추세츠 주
  헌법은 다음과 같이 명시한다: ``입법부 각 지부는 물론, 주지사나
  위원회는 중요한 법률 문제 및 중대한 사안에 대해 대법원 판사들의 의견을
  요구할 수 있다.'' \emph{Mass. Const.}, pt.~II, ch.~III, art. II
  (1964년 개정). 이 절차는 본 논문 서두에서 논의된 \emph{Goodridge} 판결
  이후 몇 달간 실제 사용되었다. \emph{Opinions of the Justices to the
  Senate}, 802 N.E.2d 565 (Mass. 2004) 사건에서, 매사추세츠 대법원은
  동성 커플에 대해 시민 결합(civil union)을 허용하고 차별을 금지하는
  입법이 헌법에 명시된 동성결혼 금지에 대한 반대 사유를 피하는 데
  충분하지 않다고 판단하였다.} 사전적 사법심사는 그 성격을 어떻게 이해할
것인지에 대해 논쟁이 있다. 예컨대, 해당 재판소가 다원적 입법
절차(multicameral legislative process)의 마지막 단계에서 전통적인
상원(senate)처럼 기능한다면, 이는 실제로는 사법심사라기보다는 또 다른
입법기관의 역할에 가깝다(비록 선출되지 않은 기관에게 권한을 부여하는
것에 대한 비판 논리는 유사할 수 있더라도).\footnote{Jeremy Waldron,
  \emph{Eisgruber's House of Lords}, 37 \emph{U.S.F. L. Rev.} 89 (2002).}
나는 이 문제에 대해 더 이상 논하지 않겠다. 일부 사법심사
옹호론자들에게는, 특정 사건에 뿌리내린 사법심사의 사후적 성격(a
posteriori character)---즉 구체적 사건에의 연결성(rootedness in
particular cases)\footnote{infra Section IV.A 참조.} ---이 핵심적인
의미를 지니며, 본 에세이에서도 그 측면에 초점을 맞출 것이다.

네 번째 구분은 세 번째 구분과 연계되어 있다. 사법심사는 일반
법원(ordinary courts)에 의해 수행될 수도 있고(앞서 언급한 매사추세츠
사건처럼), 헌법문제만을 전담하는 전문 헌법재판소(specialized
constitutional court)에 의해 수행될 수도 있다. 이 구분은 내가 후속
장에서 전개할 논점과 관련될 수 있다. 즉, 일반 법원 체계 내의 판사들이
권리에 관해 사유할 수 있는 능력은 과장되어 있으며, 판결 실무에 따르는
여러 규범적 부담들로 인해 이들은 도덕적 논증(moral arguments)을
직접적으로 다룰 여유를 갖지 못한다. 전문 헌법재판소는 보다 나은 기능을
수행할 수 있을지도 모르지만, 경험에 따르면 이들 역시 자신들의 독자적
이론(doctrines)과 판례(precedents)의 발전에 몰두하게 되고, 그로 인해
권리 기반의 논증(rights-based reasoning)에 왜곡된 필터를 씌우게 되는
경우가 많다.

\subsection{II. 네 가지 전제 (FOUR
ASSUMPTIONS)}\label{ii.-uxb124-uxac00uxc9c0-uxc804uxc81c-four-assumptions}

나의 주장을 보다 분명히 하고, 사법심사(judicial review)에 대한 반론이
가장 명확하게 적용되는 \textbf{핵심 사례(core case)}와, 특수한 병리적
상황을 다루기 위한 예외적 장치(anomalous provision)로서 사법심사가
적절하다고 여겨질 수 있는 \textbf{비핵심 사례(non-core cases)}를
구분하기 위해, 나는 몇 가지 전제(assumptions)를 설정하고자
한다.\footnote{이 가정들은 \emph{Jeremy Waldron, Some Models of Dialogue
  Between Judges and Legislators}, 23 \emph{Sup. Ct. L. Rev.} 2d 7,
  9--21 (2004)에 제시된 내용에서 변형된 것이다.}

\textbf{(p.~1360)} 이 전제들 중 일부는 독자에 따라 \textbf{순환
논법(question-begging)}처럼 보일 수도 있다. 그러나 여기서 나는 어떤 식의
기만(subterfuge)도 의도하지 않는다. 이러한 전제에서 출발하는 이유는 글의
흐름 속에서 자연스럽게 드러날 것이며, 이 전제들 중 하나 이상이 충족되지
않는 경우를 \textbf{비핵심 사례(non-core cases)}로 이해하는 가능성 또한
충분히 인정하며, 이는 제7부에서 논의될 것이다. 실질적으로, 내가 주장하는
바는 다음과 같다: \textbf{사법심사에 대한 비판은 조건부 주장이다}. 이
조건들 중 하나라도 성립하지 않는다면, 그 주장은 성립하지 않을 수
있다.\footnote{infra note 136 및 본문 참조.} 덧붙이자면, 이 에세이에서
내가 비판하고자 하는 것은 사법심사에 대한 \textbf{결론 중심(bottom-line
mentality)} 사고방식의 일종이다.\footnote{정치철학에서의
  '결과지상주의(bottom-line mentality)'에 대한 일반적인 비판은 다음에서
  확인할 수 있다: \emph{Jeremy Waldron, What Plato Would Allow}, in
  \emph{Nomos XXXVII: Theory and Practice} 138 (Ian Shapiro \& Judith
  Wagner DeCew eds., 1995).} 일부 독자들은 이 전제들을 훑어보며,
그것들이 자신이 이해하는 바의 미국 혹은 영국 사회에 적용되지 않는다는
이유만으로 나의 핵심 논거를 무시할 수도 있다. 그들에게 중요한 것은,
어떻게든 사법심사를 옹호하고 이에 대한 도전을 차단하는 것이다. 방법은
중요하지 않다. 이것은 불행한 접근 방식이다. 보다 나은 방법은, 먼저
사법심사에 대한 핵심 반론의 기반이 무엇인지 이해하고, 그것이 자체 기준에
따라 타당한지 검토한 후, 그러한 논거의 적용이 문제가 되는 사례들을
살펴보는 것이다.

이제 내가 설정할 네 가지 전제를 요약해서 제시하겠다. 다음 조건들을 갖춘
사회를 상정해 보자: (1) \textbf{민주적 제도(democratic institutions)}가
비교적 잘 작동하고 있으며, 여기에는 보통선거(universal adult suffrage)를
통해 구성된 대의 입법기관(representative legislature)이 포함된다; (2)
\textbf{사법 제도(judicial institutions)} 역시 비교적 잘 작동하며,
대표성 없는 방식(nonrepresentative basis)으로 설계되어, 개인 간의 소송을
심리하고 분쟁을 해결하며 \textbf{법치(rule of law)}를 유지하는 기능을
수행한다; (3) 사회 구성원의 다수와 그 공직자들 대부분이, \textbf{개인적
권리(individual rights)} 및 \textbf{소수자 권리(minority rights)}의
관념에 진지하게 전념하고 있다; (4) 권리에 대한 해석---즉, 권리에 대한
전념이 실제로 무엇을 의미하며 어떤 함의를 갖는지를 둘러싼---사회 내의
구성원들 사이에 \textbf{지속적이며 실질적이고 성실한 의견
차이(disagreement in good faith)}가 존재한다.

나는 이 네 가지 전제가 성립할 경우, 해당 사회는 \textbf{권리에 관한 의견
차이}를 입법 제도(legislative institutions)를 통해 해결해야 한다는 점을
논증할 것이다. 이러한 전제들이 성립한다면, 그러한 의견 차이를 최종적으로
사법기관(judicial tribunals)에 위임해야 한다는 주장은 설득력이 약하며
정당성을 갖기 어렵다. 입법부가 권리에 관해 내린 결정을 사법부가 재심하는
것(second-guessing) 역시 필요하지 않다. 그리고 사법부가 이와 같은 사안에
대해 입법부의 결정을 \textbf{무효화하거나 대체하는 방식으로 결정권을
행사하는 것}은 정치적 정당성(political legitimacy)의 핵심 기준을
충족하지 못한다는 점을 논증할 것이다. 다음으로, 이 네 가지 전제를 하나씩
구체적으로 설명하겠다.

\subsubsection{A. 민주적 제도 (Democratic
Institutions)}\label{a.-uxbbfcuxc8fcuxc801-uxc81cuxb3c4-democratic-institutions}

\textbf{(p.~1361)} 나는 우리가 상정하는 사회가
\textbf{민주사회(democratic society)}라고 가정한다. 이 사회는 현대 서구
세계의 대부분과 마찬가지로, 과거에는 군주제(monarchy), 전제정(tyranny),
독재(dictatorship), 또는 식민 지배(colonial domination)의 다양한 형태를
경험했으나, 이제는 그 법률이 제정되고 공적 정책(public policies)이
국민과 그 대표자들에 의해 \textbf{선거를 통한 제도(elective
institutions)}를 매개로 설정되는 상태에 도달해 있다. 이 사회는
\textbf{보통선거(universal adult suffrage)}를 특징으로 하는 폭넓은
민주적 정치 체제(democratic political system)를 가지고 있으며, 공정하고
정기적인 방식으로 선출되는 \textbf{대의제 입법부(representative
legislature)}를 갖추고 있다.\footnote{따라서, 비민주적 사회나 이와 매우
  다른 제도를 가진 사회에 본 논문의 논증을 적용하는 것은 논외이다.} 나는
이 입법부가 규모 있는 \textbf{심의 기구(deliberative body)}이며,
정의(justice)와 사회정책(social policy)이라는 중요한 사안들을 포함한
어려운 쟁점들을 다루는 데 익숙하다고 전제한다. 입법자들(legislators)은
공적 사안에 대해 심의(deliberation)하고 표결(vote)하며, 입법 절차는
정교하고 책임성 있으며,\footnote{Jeremy Waldron, \emph{Legislating with
  Integrity}, 72 \emph{Fordham L. Rev.} 373 (2003).} 다음과 같은 여러
안전장치들을 내포한다: \textbf{양원제(bicameralism)},\footnote{양원제(bicameralism)에
  대한 전제는 문제적으로 보일 수 있다. 세계에는 잘 작동하는
  단원제(unicameral) 입법부도 다수 존재하는데, 특히 스칸디나비아
  국가들---덴마크, 노르웨이, 스웨덴---이 그렇다. 그러나 단원제는 다른
  입법적 병리 현상을 쉽게 악화시킬 수 있다. 뉴질랜드에서 이 문제가
  어떻게 나타나는지를 주장한 논변은 Jeremy Waldron, \emph{Compared to
  What? - Judicial Activism and the New Zealand Parliament}, 2005 N.Z.
  L.J. 441 참조.} 강력한 위원회 심사(robust committee scrutiny), 그리고
다단계의 검토, 토론, 표결 과정. 이러한 과정들은 공식적으로는
공청회(public hearings) 및 협의 절차(consultation procedures)를 통해,
비공식적으로는 사회 전반의 보다 넓은 논의들과 연결되어 있다고 가정한다.
입법부 구성원들은 자신들을 다양한 방식으로
\textbf{대표자(representatives)}라고 인식하는데, 때로는 지역구
유권자(constituents)의 이익과 의견을 자신들의 활동의 핵심으로 삼고,
때로는 사회 전체의 이익과 의견을 \textbf{가상적 대표(virtual
representation)}한다는 관점에서 활동하기도 한다. 나는 또한
정당(political parties)이 존재하며, 입법자들의 정당 소속이 그들의 견해
형성에 있어 지역 유권자들의 이해관계를 넘어서는 보다 넓은 시야를
제공한다고 가정한다.

이상의 전제는 논쟁적이기보다는 일반적으로 \textbf{민주주의적
입법부(democratic legislatures)}가 통상적으로 작동하는 방식을 설명한
것이다. 전반적으로, 나는 이 민주주의 제도들이 비교적 잘 작동하고 있다고
전제한다. 물론 이들이 완전하지는 않을 수 있고, 개선 방법에 대해서는
지속적인 논의가 있을 수 있다. 나는 이러한 논의가 \textbf{책임 있는
심의(responsible deliberation)}와 \textbf{정치적 평등(political
equality)}이라는 민주주의 문화(culture of democracy)의 가치를 바탕으로
진행된다고 본다. 이 중에서 두 번째 가치, \textbf{(p.~1362)} 즉 정치적
평등은 특히 강조할 가치가 있다. 나는 이 사회에서 입법의 제도, 절차,
실천이 지속적으로 정치적 평등의 관점에서 검토되고 있다고 전제한다. 만일
대표성(representation)에 있어 심각하게 정치적 평등의 이상을 훼손하는
불균형이 존재한다면, 그것이 정당한 비판이라는 점을 사회 구성원 전체가
인식하고 있으며, 필요하다면 입법부와 선거 제도 자체가 이를 교정하기 위해
변경되어야 한다는 점도 널리 받아들여지고 있다고 본다. 또한 나는 입법부가
이러한 변경을 자발적으로 혹은 국민투표(referendum)를 통해 조직할 수 있는
능력을 갖추고 있다고 가정한다.\footnote{때때로 선출 제도(elective
  institutions)는 자기 개혁이 불가능하다는 주장이 제기된다. 이는
  입법자들이 현상 유지에 구조적으로 이해관계를 갖기 때문이라는 것이다.
  이는 미국의 병리적인 선거 및 입법 구조에 대해서는 어느 정도 사실일 수
  있다. (그러나 미국 내에서도 이 주장이 가장 타당하게 적용되는 쟁점은,
  법원이 거의 개입하지 않은 사안들이다. 예컨대, 미국의 선거구 재획정
  제도의 심각한 상태를 고려하라.) 그러나 다른 국가들에는 명백히 적용되지
  않는다. 예컨대, 뉴질랜드에서는 1993년에 의회가 선거제도를
  단순다수제(first-past-the-post system)에서 비례대표제(proportional
  representation)로 변경하는 입법을 하였으며, 이는 기존의 현역 의원
  구조를 불안정하게 만들었다. 관련 법률로는 Electoral Act 1993, 1993
  S.N.Z. No.~87; Electoral Referendum Act 1993, 1993 S.N.Z. No.~86 참조.}

내가 이렇게 민주주의 문화와 선거 및 입법 제도의 기능에 대해 장황하게
설명하는 이유는, 그것이 이후 전개될 논증의 핵심이 되기 때문이다. 논증의
초기 구조는 다음과 같은 질문을 던지는 것으로 시작될 것이다: 이 첫 번째
전제가 성립한다고 가정할 때, 왜 선출되지 않은 기구(nonelective
process)를 설정하여 입법부의 활동을 심사하거나 때로는 무효화하는 제도를
만들어야 할 이유가 있는가? 반면, 나는 이 첫 번째 전제를 통해 어떠한
문제를 은폐하려는 것도 아니다. 이 전제를 곧바로 두 번째 전제와 균형 있게
제시하고자 한다. 즉, 우리가 상정하는 사회에는 또한 \textbf{잘 작동하는
법원(courts in good working order)}이 존재하며, 이들은 사법부가 잘
수행할 수 있는 기능들을 비교적 잘 해내고 있다고 전제한다. 우리가
상정하는 이 사회는, 만일 사법심사(judicial review)가 타당하게 정당화될
수 있다면, 이를 작동시킬 수 있는 조건을 갖추고 있는 사회다.

한 가지 주의할 점이 있다. 내가 여기서 제도들이 ``잘 작동하고 있다(good
working order)''고 말할 때, 그것은 해당 입법부가 제정하는 법률들이 그
\textbf{내용(content)} 면에서 대체로 선하거나 정당하다는 것을 의미하지
않는다. 나는 일부 법률은 정당하고 일부는 부정당하며(무엇이 어느 쪽인지는
사람마다 견해가 다를 수 있다), 이것은 사법심사의 대상이 될 수도 있는
법률뿐 아니라 아무도 사법심사의 대상으로 삼지 않는 법률에 대해서도
마찬가지라고 본다. 지금까지 입법 및 선거 제도가 잘 작동하고 있다고 말한
것은 \textbf{결과 가치(outcome values)}가 아니라 \textbf{절차적
가치(process values)}에 대한 것이다. 그러나 제5부(Part V)에서는, 이러한
과정에서 기대할 수 있는 \textbf{추론의 양상(reasoning)}에 대해 더
구체적으로 논의할 것이다.

\subsubsection{B. 사법 제도 (Judicial
Institutions)}\label{b.-uxc0acuxbc95-uxc81cuxb3c4-judicial-institutions}

\textbf{(p.~1363)}우리가 상정하는 사회에는 법원(courts), 즉 소송을
심리하고 분쟁을 해결하며 \textbf{법치(rule of law)}를 유지하기 위해
설립된, 잘 정비되어 있고 정치적으로 독립적인
\textbf{사법부(judiciary)}가 존재한다고 전제한다. 이 사법 제도는 비교적
잘 작동하고 있으며, 행정행위(executive actions)에 대한 사법심사(judicial
review)를 수행할 법적 권한을 이미 부여받고 있으며, 이는
법률적(statutorily) 및 헌법적(constitutionally) 기준에 따라 이루어진다고
가정한다.

이 사법 제도는 앞 절에서 다룬 입법 제도들과 달리, 대부분
\textbf{선출되거나 대표성을 지닌 제도(elective or representative
institutions)}가 아니라고 전제한다. 이는 단지 판사(judges)의 직위가
대부분 선거직이 아니라는 것만을 의미하는 것이 아니다. 나아가 사법부
전체가 선거, 대표성(representation), 선거 책임성(electoral
accountability)의 윤리(ethos)에 관통되어 있지 않다는 점에서도 입법부와
구별된다. 사법심사를 옹호하는 많은 이들은 이것을 큰 장점으로 본다.
왜냐하면 법원이 대중의 압력(popular pressures)에서 자유롭고, 여론의
분노에 영향받지 않으며, 원칙에 기반한 문제(principled issues)에 대해
집중적으로 숙고할 수 있기 때문이다. 그러나 때때로 사법심사에 대한
민주주의적 비판에 반박할 필요가 있을 때, 옹호자들은 오히려 판사가
선출되는 주(state)를 자랑스럽게 언급하기도 한다. 이는 미국의 일부 주에서
실제로 일어나는 일이다. 하지만 판사가 선출된다고 하더라도, 법원의 운영은
일반적으로 입법부처럼 대표성과 선거 책임성의 윤리(an ethos of
representation and electoral accountability)하에서 이루어지지 않는다.

나는 우리가 상정하는 사회에서, 법원이 사법심사(judicial review)의 제도
하에서 부여받게 될 기능들을 수행할 능력이 있다고 전제한다. 즉, 법원은
입법을 심사할 수 있다. 그러나 그들이 \textbf{입법을 심사해야 하는가},
그리고 \textbf{그들의 결정이 대표기관들(representative branches of
government)에 대해 최종적이고 구속력을 가지는가}가 문제다. 나는, 만일
법원에 이러한 기능이 부여된다면, 그것을 \textbf{법원이 본질적으로 자신의
기능을 수행하는 방식}으로 실행할 것이라고 전제한다. 사법 절차(judicial
process)의 성격과, 법원이 제도적으로 적합하거나 그렇지 않은 과제에 관한
방대한 법학 논문(law review literature)이 존재하지만,\footnote{예: Henry
  M. Hart, Jr.~\& Albert M. Sacks, The Legal Process: Basic Problems in
  the Making and Application of Law 640-47 (William N. Eskridge, Jr.~\&
  Philip P. Frickey eds., 1994); Lon L. Fuller, The Forms and Limits of
  Adjudication, 92 Harv. L. Rev.~353 (1978).} 여기서 그 논의로 깊이
들어가고자 하지는 않는다. 앞서 언급했듯이, 나는 다음과 같은 법원의
일반적 특성을 전제로 삼는다: (1) 법원은 \textbf{자기발의(self-motion)}
또는 \textbf{추상적 기준(abstract reference)}에 따라 행동하지 않으며,
특정 원고(litigants)가 제기한 특정한 청구(claim)에 반응한다. (2) 문제는
이항적(binary), \textbf{대립적 절차(adversarial presentation)}의
맥락에서 다루어진다. (3) 법원은 관련 사안에 대해 자신들의 선례(past
decisions)를 참조하고 이를 발전시키는 방식으로 판단한다. 나는 또한,
익숙한 \textbf{법원 위계 구조(hierarchy of courts)}가 존재하며,
\textbf{(p.~1364)} 항소 절차(provisions for appeal)가 마련되어 있고,
상급심에서는 다수의 판사들(예: 5인 또는 9인의 재판부)이 사건을 심리하며,
하급심 법원들은 상급심 판례에 실질적으로 구속된다는 점도 전제한다.

일부 사회에서는 판사들이 별도의 교육을 통해 특별히 양성되며, 다른
사회에서는 유능한 변호사나 법학자들 중에서 판사를 선출한다. 어느
경우이든, 이들은 정치 체제 내에서 높은 지위를 가지며, 특정한 정치적
압력으로부터는 어느 정도 보호받는 위치에 있다고 전제한다. 또한 이들은 그
사회의 다른 고학력 고지위 계층(high-status and well-educated members)과
유사한 사회적 성격을 갖고 있다고 본다. 이러한 전제는 두 가지 이유로
중요하다. 첫째, 그 사회가 \textbf{민주사회(democratic society)}로서의
자부심을 지니고 있으므로, 나는 판사들 또한 이러한 자부심을 공유하고
있으며, 자신들이 입법에 대해 사법심사를 수행할 경우 \textbf{자신들의
정당성(legitimacy)}에 대해 의식적일 것이라고 가정한다. 이는 그들이 해당
권한을 어떻게 행사하는지에 영향을 줄 수 있다.\footnote{Jesse H. Choper,
  Judicial Review and the National Political Process (1980) (이 맥락에서
  대법원의 정당성을 논의함); 또한 Planned Parenthood of Se. Pa. v.
  Casey, 505 U.S. 833, 864-69 (1992) 참조 (동일 논의).} 둘째, 판사들
역시 다른 시민들과 마찬가지로 권리에 전념할 가능성이 충분히 있지만,
\textbf{개인적 권리(individual rights)}와 \textbf{소수자 권리(minority
rights)}의 의미와 함의에 관해서는 서로 견해가 다를 수 있다. 즉, 판사들
역시 내가 네 번째 전제로 제시한 \textbf{권리에 대한
의견불일치(disagreement about rights)}의 조건에 해당하며, 이것이 또한
사법심사 권한을 행사할 때 영향을 미칠 수 있다. 구체적으로 말하자면,
입법자들과 마찬가지로, \textbf{다수의 판사들로 구성된 재판부(multi-judge
tribunal)} 내에서 권리에 대한 의견이 갈릴 경우, 그에 대한 결정
절차(decision-procedure)가 필요하다. 이때 가장 일반적으로 사용되는
방식은 단순한 \textbf{다수결(simple majority voting)}이다. 나는
제5부에서 이러한 방식이 판사들이 사용하기에 적절한 절차인지 여부에 대해
다룰 것이다.

\subsubsection{C. 권리에 대한 전념 (A Commitment to
Rights)}\label{c.-uxad8cuxb9acuxc5d0-uxb300uxd55c-uxc804uxb150-a-commitment-to-rights}

내가 전제하는 바는 우리가 상정하는 사회의 대부분 구성원들은
\textbf{개인적 권리(individual rights)} 및 \textbf{소수자 권리(minority
rights)}의 관념에 대해 강한 전념(commitment)을 가지고 있다는 것이다.
그들은 넓은 의미에서의 공리주의적 개념(utilitarian conception)에 따라
공공선(the general good)을 추구해야 한다고 믿고 있으며, 또한 정치의
일반적 원칙으로서 다수결(majority rule)을 수용하지만, 그럼에도 불구하고
개인은 일정한 이익(interests)을 가지며 특정한 자유(liberties)를
보장받아야 하고, 그것이 단지 다수에게 편리하다는 이유만으로 부정되어서는
안 된다는 점을 받아들인다. 또한 이들은 소수자(minorities)는 단순한 수적
기반이나 정치적 영향력에 의해 자동으로 보장되지 않는 일정 수준의
\textbf{지원(support)}, \textbf{인정(recognition)},
\textbf{보호(insulation)}를 받을 권리가 있다고 믿는다.

\textbf{(p.~1365)} 이 사회에서 유행하는 권리 이론(theory of rights)의
세부 내용은 여기서 따로 다룰 필요가 없다. 나는 이 사회 전체의 권리에
대한 전념이 다음과 같은 인식을 포함하고 있다고 본다: 세계 인권
합의(worldwide consensus on human rights)에 대한 이해와, 권리에 관한
사유의 역사(history of thinking about rights)에 대한 인식.\footnote{이러한
  인식은, 인권 협약이 1945년 이후 세계적으로 부상해 왔다는 점과 그
  역사가 1776년의 미국 독립 선언문 및 1789년의 프랑스 인권 선언과 같은
  문서에서 암시되는 자연권(natural right) 개념으로 거슬러 올라간다는
  막연한 이해 이상의 것이 아닐 수도 있다.} 또한 이 전념은 살아 있는
합의(living consensus)로서, 권리를 옹호하는 이들이 서로 어떤 권리를
가지고 있으며 그 권리가 무엇을 의미하는지를 끊임없이 대화하는 가운데
발전하고 진화하고 있다고 본다. 나는 이 권리에 대한 전념이 단순한
\textbf{수사적 선언(lip service)}에 그치지 않고, 사회 구성원들이 권리를
\textbf{진지하게 받아들이고(take rights seriously)} 있다고 전제한다. 즉,
그들은 권리에 대해 관심을 갖고 있으며, 자신의 견해와 타인의 권리관을
끊임없이 숙고하고 활발히 논의하며, 사회 내의 모든 공적 결정과 논의에서
권리 문제에 민감하게 반응한다.

물론 어느 사회에나 권리에 대해 회의적인 사람(skeptics)은 존재하겠지만,
나는 그와 같은 입장은 \textbf{주류에서 벗어난 예외(outlier)}라고 본다.
어떤 이들은 모든 정치적 도덕성(political morality)을 거부하는 입장에서
권리를 부정하고, 또 어떤 이들은 공리주의(utilitarianism),
사회주의(socialism), 혹은 기타 이념에 따라 정치적 도덕성의 관점에서
권리를 부정한다. (예: 권리는 지나치게 개인주의적이며, 권리의 우선성은
효율성의 합리적 추구를 방해한다는 식의 주장.) 그러나 나는 \textbf{개인적
권리와 소수자 권리에 대한 일반적인 존중(general respect)}이 그 사회의
폭넓은 합의의 진지한 일부이며, 정치적 여론 중 가장 보편적인 입장의
일부이며, 공식적 이데올로기(official ideology)의 일부라고 가정한다.

이 세 번째 전제를 보다 구체화하기 위해, 우리는 그 사회가 익숙한 형태의
공식적인 \textbf{권리장전 또는 권리선언(Bill or Declaration of
Rights)}을 채택할 정도로 권리를 소중히 여긴다고 가정할 수 있다. 이
글에서는 이를 해당 사회의 \textbf{권리장전(Bill of Rights)}이라 부를
것이다. 이는 예컨대 다음과 같은 것에 대응할 수 있다: 미국 헌법과 그 수정
조항들에 포함된 권리 조항, 캐나다의 \textbf{「권리와 자유에 관한
헌장(Charter of Rights and Freedoms)}」, 유럽인권협약(European
Convention on Human Rights) (예: 「인권법(Human Rights Act)」을 통해
영국 법에 통합된 형태), 또는 뉴질랜드의 「권리장전법(Bill of Rights
Act)」. 이 중 마지막 예시를 잘 아는 사람이라면 내가 여기서 그 권리장전이
\textbf{헌법에 포함되었거나(entrenched)} 또는
\textbf{성문헌(constitutional document)}의 일부라고 전제하고 있지 않다는
점을 이해할 것이다. 나는 그 점을 열어 두고 싶다. 내가 이 시점에서
전제하는 것은 단지 해당 사회가 권리에 대한 전념을 구체화하기 위해 어떤
형태로든 \textbf{권리장전(Bill of Rights)}을 제정했다는 사실뿐이다. 이는
해당 사회가 자발적으로 어느 시점에 제정했을 수도 있고, 다른 국가를
모방해서 도입했을 수도 있으며, 또는 국제 인권법(human rights law)에 따른
대외적 의무(external obligations)를 이행한 결과일 수도 있다.

\textbf{(p.~1366)} 이러한 전제들이 독자에게는 다소 의아하게 느껴질 수도
있다. 한편으로는, 내가 권리장전(Bill of Rights)의 존재를 전제로
삼음으로써, 사법심사(judicial review)에 유리하도록 논거의 전제를
설정하는 것처럼 보일 수 있다. 다른 한편으로는, 뭔가 \textbf{교묘한
술수(sneaky)}가 예고되는 것처럼 느껴질 수도 있다. 실제로 일부 독자들은
내가 과거에 \textbf{사법심사란 권리 옹호자와 권리 반대자 사이의 대결이
아니라, 서로 다른 권리관(rights) 사이의 대결로 이해되어야 한다}는 점을
논증한 바 있다는 사실을 알고 있을 수도 있다.\footnote{Waldron, A
  Right-Based Critique, supra note 14, at 28--31, 34--36 참조.} 이러한
두 가지 관찰에 대해 내가 강조하고 싶은 점은 다음과 같다: 문화적 차원과
제도적 차원 모두에서, \textbf{권리에 대한 전념(commitment to
rights)}(설령 그것이 문서화된 전념이라 하더라도)과 그러한 전념이 취할 수
있는 특정한 \textbf{제도적 형태(institutional form)}(예: 입법에 대한
사법심사) 사이에는 명확한 구분이 존재한다는 점이다. 나는 사법심사를
반대하는 사람들이 '권리에 회의적인 자들(rights-skeptics)'로 폄하되는
것을 지겹도록 들어왔다. 이에 대한 최선의 대응은, \textbf{권리에 대한
강력하고 광범위한 전념을 전제로 하여 사법심사에 대한 비판적 논거를
구성하는 것}이다.

이 세 번째 전제는, \textbf{권리에 대한 전념이 희박하거나 불안정한
사회}를 \textbf{비핵심 사례(non-core cases)}로 정의한다. 이러한 방식의
논의가 다소 이상하거나 부당하게 보일 수도 있다. 왜냐하면 사법심사를
옹호하는 이들 중 일부는, 우리가 권리에 대한 전념을 강화하거나, 새로이
민주화된 사회의 시민들에게 권리의 가치를 교육하거나, 다수결 체계에서
보호받지 못할 수 있는 소수자들에게 보장을 제공하기 위해
\textbf{사법심사가 필요하다}는 점을 논증하기 때문이다. 이러한 논변들은
흥미롭지만, 미국, 영국, 캐나다와 같은 국가들에서 사법심사를 정당화하는
핵심 주장에 직접적인 영향을 주지는 않는다. 이들 국가에서는 사법심사가
\textbf{기존의 권리 전념을 제도화하거나 운영하기 위한 적절한 방식}이라고
주장된다. 나는 이러한 주장을 액면 그대로 받아들이며, 그것이 바로 내가 이
세 번째 전제를 설정하는 이유이다.\footnote{여기서의 접근은 John Rawls와
  유사하다. 나는 핵심사례(core case)라는 장치를 이용하여 일종의 공적으로
  수용된 정의 이론을 갖춘 잘 정돈된 사회(well-ordered society)를
  규정하고자 한다. 예: John Rawls, Political Liberalism 35--36 (1993)
  {[}이하 Rawls, Political Liberalism{]}; Rawls는 입법에 대한
  사법심사(judicial review)가 잘 정돈된 사회에도 적절하다고 보는 듯하다.
  Id. at 165--66, 233--40 참조; 또한 John Rawls, A Theory of Justice
  195--99, 228--31 (1971) {[}이하 Rawls, A Theory of Justice{]} 참조.
  내가 목표하는 바는 Rawls의 이러한 판단이 잘못되었음을 보여주는 것이다.}

\subsubsection{D. 권리에 대한 의견불일치 (Disagreement About
Rights)}\label{d.-uxad8cuxb9acuxc5d0-uxb300uxd55c-uxc758uxacacuxbd88uxc77cuxce58-disagreement-about-rights}

내 마지막이자 핵심적인 전제는, 권리에 대한 합의(consensus)가 오늘날의
자유주의 사회들에서 모든 주요 정치적 쟁점들에 나타나는 일반적인
의견불일치(disagreement)의 발생에서 예외적이지 않다는 것이다. 다시 말해,
나는 \textbf{(p.~1367)} 어떤 권리가 존재하며 그것이 어떤 내용을 갖는지를
둘러싼 \textbf{실질적인 의견불일치(substantial dissensus)}가 존재한다고
전제한다. 이러한 의견불일치 중 일부는 철학적 수준에서 명백히 드러난다
(예: 사회경제적 권리(socioeconomic rights)가 권리장전(Bill of Rights)에
포함되어야 하는가), 일부는 추상적인 권리 원칙을 특정 입법 제안과
연관시키려 할 때 드러난다 (예: 종교의 자유로운 실현(free exercise of
religion)이 일반적으로 적용되는 법률로부터의 예외를 요구하는가), 또
일부는 개별적인 곤란한 사건의 맥락에서만 분명해진다 (예: 국가적 위기
시기에도 이견 표현(dissident speech)을 어느 정도까지 허용해야 하는가).

나는 이러한 권리-의견불일치가 대부분 좁은 의미의 법률적 해석 문제(issues
of interpretation)가 아니라고 전제한다. 물론 처음에는 해석의 문제로
제시되겠지만, 실제로는 공동체 전체에 중대한 실천적 함의를 지닌 쟁점들을
내포한다. 나는 다른 글에서 이와 같은 쟁점을 \textbf{``분수령적 권리
쟁점(watershed issues of rights)''}이라고 부른 바 있다.\footnote{Waldron,
  Judicial Power and Popular Sovereignty, supra note 14, at 198 참조.}
이들은 수많은 사람들의 삶에 중대한 영향을 미치는 정치철학(political
philosophy)의 주요 주제들로, 단지 특정 사회에 국한된 특이한 문제들이
아니다. 오히려 이러한 쟁점들은 모든 현대 사회가 직면해야 하는 중대한
선택지를 정의하며, 기존의 도덕적·정치적 논의의 맥락에서 적절히 이해되고,
많은 사회들에서 도덕적·정치적 의견불일치의 핵심 지점들이 된다. 즉시
떠오를 수 있는 예시들로는 낙태(abortion), 소수우대(affirmative action),
정부의 재분배 및 시장 개입의 정당성, 피의자의 권리(rights of criminal
suspects), 종교 관용(religious toleration)의 의미, 소수문화권의
권리(minority cultural rights), 선거운동에서의 표현 및 지출의 규제 등이
있다.

이러한 예시들이 보여주듯, 권리에 대한 의견불일치는 종종 주변적인 적용이
아니라 \textbf{중심적인 적용(central applications)}에 관한 것이다. 나는
이미 권리에 대한 일반적 헌신을 전제했기 때문에, 이 헌신이 각 권리의
중심부(core)를 포함하며, 논란은 오직 그 권리의 외곽 영역에서만
발생한다고 추론하고 싶어질 수 있다. 그러나 이것은 잘못된 생각이다.
권리에 대한 헌신은 완전하고 성실한(wholehearted and sincere) 것이면서도,
여전히 분수령적 사안에서는 논란이 있을 수 있다. 예컨대, 인종차별적
혐오발언(racist hate speech)에 대한 규제가 허용 가능한가에 관해 의견이
다른 두 사람은 모두 표현의 자유(free speech)가 이 사안을 사고하는 데
핵심적 권리임을 인정할 수 있으며, 동시에 그들이 논쟁 중인 사안이 표현의
자유라는 권리에 비추어 주변적인 사안이 아니라 중심적인 사안임을 받아들일
수 있다. 이로부터 우리는 아마도 그들이 서로 다른 \textbf{권리
개념(conception of the right)}을 갖고 있음을 알 수 있을 것이나, 이것이
곧 그들의 권리에 대한 헌신이 불성실하다는 이유가 되지는
않는다.\footnote{권리 개념(the concept of a right)과 그것의 다양한
  관념들(conceptions of it)의 구분에 대해서는 Ronald Dworkin, Taking
  Rights Seriously 134--36 (1977) 참조.}

\textbf{(p.~1368)} 일반적으로 말해, 사람들이 권리에 관해 의견이 다르다고
해서 그 중 한 사람은 권리를 진지하게 여기지 않는다고 결론지어야 하는
것은 아니다. 물론 어떤 입장은 권리를 아예 무시하는 \textbf{불성실한
자(scoundrels)}나, 권리의 중요성을 이해하지 못하는 \textbf{도덕적
문맹자(moral illiterates)}에 의해 부정직하거나 무지하게 주장되고 있을
수도 있다. 그러나 나는 대부분의 경우, 이와 같은 의견불일치는
\textbf{합리적이고 성실한 태도(reasonably and in good faith)}로 제기되고
있다고 전제한다. 이들은 심각한 쟁점들이며, 이에 대해 사회 구성원들 간에
합의(consensus)가 있을 것이라 기대하는 것은 비현실적이다. 다시 말해,
나는 John Rawls가 말한 \textbf{판단의 부담(burdens of judgment)}이라는
개념을 전제하는데, Rawls는 이를 '좋음(the good)'의 문제에 주로
적용했지만, 나는 그것을 '옳음(the right)'의 문제에도 확장 적용하고자
한다.\footnote{Rawls, Political Liberalism, supra note 53, at 54--58
  (``판단의 부담들(burdens of judgment)''에 대한 논의). Rawls는 ``우리의
  가장 중요한 판단 중 다수는, 합리적 능력을 갖춘 양심적인 사람들이
  자유로운 토론을 거친 후에도 동일한 결론에 도달할 것으로 기대하기
  어려운 조건하에서 이루어진다''고 주장한다. Id. at 58. 이 논의를
  권리(right)뿐 아니라 선(the good)에도 적용한 논변은 Waldron, supra
  note 14, at 149--63 참조.} 즉, 복잡하고 긴장된 권리 쟁점들에 관해
사람들의 견해가 언제나 하나의 합의로 수렴할 것이라 기대하는 것은
합리적이지 않다. Rawls가 강조하듯, ``{[}이러한 문제들에 대한{]} 우리의
모든 차이가 무지, 고의적 왜곡, 혹은 권력·지위·경제적 이득에 대한
경쟁에서 비롯된 것이라고 가정하는 것은 비현실적이다.''\footnote{Rawls,
  Political Liberalism, supra note 53, at 58.}

이러한 의견불일치 전제는 도덕적 상대주의(moral relativism)와는 아무
관련이 없다. 우리는 권리와 정의에 관한 의견불일치의 존재를 인식할 수
있으며, 심지어는 실천적 정치 목적상 그것들이 해결 불가능하다는 점을
인정할 수도 있다. 그러나 그렇다고 해서 논쟁이 벌어지는 쟁점에 대해 어떤
\textbf{진리(fact of the matter)}도 존재하지 않는다는 \textbf{메타윤리적
주장(meta-ethical claim)}을 받아들여야 하는 것은 아니다. 권리와
헌정주의(constitutionalism)의 원칙에 관해 어떤 사실이 존재한다는 점은,
그 진리가 합리적으로 부정될 수 있는 방식으로 드러나지 않는 한,
의견불일치의 인정과 양립 가능하다.\footnote{Waldron, The Irrelevance of
  Moral Objectivity, supra note 19, at 182 참조.}

권리장전(Bill of Rights)이 존재한다면, 나는 그것이 의견불일치가 초래한
쟁점에 관련성을 가지긴 하지만, 그 자체로 쟁점을 해결하지는 않는다고
전제한다. 나는 몇 문단 앞에서 예시들을 언급한 바 있다. 미국의 경우,
권리장전 조항들이 각 쟁점을 해결함에 있어 일정한 관련성을 가진다는 점은
부정할 수 없지만, 동시에 그 조항들 자체가 그 쟁점을 합리적인 논쟁의
여지가 없을 정도로 \textbf{결정적으로 해결(determine a
resolution)}한다고 보기도 어렵다. 따라서 나는 이러한
\textbf{의견불일치의 범위(extent of disagreement)}가, 우리가 만들어내는
추상적 표현의 정교함을 압도한다고 본다. 의견불일치가 권리장전의 제정을
방해하지는 \textbf{(p.~1369)} 않는다.\footnote{Thomas Christiano,
  Waldron on Law and Disagreement, 19 Law \& Phil. 513, 537 (2000) 참조.}
그러나 의견불일치는 여전히 해결되지 않은 채 남아 있고, 따라서 특정 권리
침해의 가능성이 제기되었을 때, 해당 사안이 권리장전의 조항과 관련 있다는
점은 누구도 부인하지 않지만, 그 관련성의 구체적 의미가 무엇인지, 즉
그것이 해당 입법 조항을 금지하거나 제한해야 하는지를 두고는 합리적인
사람들 사이에 여전히 논쟁이 존재한다.\footnote{다시 말하지만, 나는
  권리장전(Bill of Rights)의 조항들이 중심사례를 포괄하며, 그 적용의
  주변에서만 의견불일치(disagreement)가 발생한다고 말하는 것이 아니다.
  이 조항들은 대체로 모호하고 추상적이어서, 명백히 논쟁의 여지가 없는
  사례가 있다고 해도, 사람들이 동일한 추상적 표현을 각기 다른 실질적
  접근에 적용할 수 있으며, 그럼에도 불구하고 여전히 모두가 권리를
  진지하게 고려하고 있다고 말할 수 있다.}

이렇다고 해서 권리장전이 해당 쟁점에 어떤 영향을 미치는지를 두고
\textbf{결정적인 논변(conclusive arguments)}이 제시되지 않는다는 말은
아니다. 적어도 그러한 논변을 제시하는 사람들에게는 그것이 그렇게 보일
것이다. 만약 사법심사가 제도화된다면, 변호사들은 권리장전의
문면(text)뿐만 아니라 그 \textbf{중력(gravitational force)}까지 활용하여
이러한 쟁점들에 관해 논증할 것이다. 사실, 변호사들에게는 활짝 열린
놀이터가 될 것이다. 각 쟁점에 대해 양측 모두가, 권리장전의 모호한
헌신(bland commitments)이 관대하게(혹은 제한적으로) 읽히기만 한다면,
자신의 입장을 그 안에서 읽어낼 수 있다고 주장할 것이다. 그러나 어느 쪽도
공개적으로 다음과 같은 점을 인정하지 않으려 할 것이다. 바로 지금 내가
당연하다고 전제하고 있는 것---즉, 권리장전의 모호한 수사적 표현(bland
rhetoric)은, 권리를 진지하게 받아들이는 사람들 사이에서 불가피하게
발생하는 \textbf{현실적이고 합리적인 의견불일치(real and reasonable
disagreements)}를 일단 충분히 무마시켜 줄 수 있는 수준에서 법안을
통과시키기 위해 고안된 것이었다는 점이다. 사법심사는 이러한 의견불일치에
우리가 직접 마주하도록 독려하기보다는, 오히려 그러한 쟁점들을 권리장전의
모호한 문구 해석 문제로 틀 지우게 한다(frame). 이런 방식이 그 쟁점들이
제기하는 도덕적 문제들에 대해 숙고하기에 바람직한 맥락인지 여부는,
제5부에서 우리가 검토하게 될 사안 중 하나이다.

\subsection{III. 논증의 형식 (The Form of the
Argument)}\label{iii.-uxb17cuxc99duxc758-uxd615uxc2dd-the-form-of-the-argument}

이제 우리는 네 가지 전제를 세웠다. 그렇다면 이러한 전제들이 규정하는
상황에서 우리는 무엇을 해야 하는가? 이 공동체의 구성원들은
권리(rights)에 헌신하고 있으나, 동시에 권리에 대해 의견이
다르다(disagree about rights). 대부분의 권리 문제는
\textbf{결정(settlement)}이 필요하다. 여기서 결정이 필요하다는 것은 그
쟁점을 제거한다는 뜻이 아니라\footnote{Jon Stewart 외, America (The
  Book): A Citizen's Guide to Democracy Inaction 90 (2004) 참조 (Roe v.
  Wade에 대한 논의에서 ``헌법상의 권리로서 태아는 인격체가 아니며,
  사생활의 권리가 여성의 낙태 결정을 보호한다는 판결로 인해, 한때
  논쟁적이었던 쟁점에 대한 모든 논의가 종결되었다''고 서술함).},
\textbf{행동이 필요한 상황에서 공동의 행동을 위한 기반(common basis for
action)}을 마련해야 한다는 의미다. 현실에는 \textbf{(p.~1370)} 공동체
전체가 결정할 필요가 없는 모든 종류의 문제들이 있다. 예를 들어,
화체변화(transubstantiation), 『햄릿』의 의미, 순수한 관조적 삶의 가치
등은 그에 대한 합의가 거의 불가능한 분야이기 때문에 사회 전체가 반드시
입장을 정해야 하는 것은 아니다. 이는 다행스러운 일이다. 하지만 권리
문제는 다르다. 우리는 권리와 관련된 쟁점들에 대해 결정이 필요하지만,
안타깝게도 이 영역에서도 합의의 전망은 그리 밝지 않다. 결정이 필요하다고
해서 의견불일치(disagreement)의 사실이 사라지는 것은 아니다. 오히려 이는
그 의견불일치의 한가운데에서 공동의 행동 근거를 형성해야 한다는 것을
뜻한다.

현실 세계에서 이러한 결정의 필요성은 입법 영역에서 가장 뚜렷이 드러난다.
우리는 특정 영역에 대해 입법을 하고, 그 입법은 권리 문제를 야기한다.
이러한 권리 문제들은 입법 자체의 표면에 명확히 드러나지 않을 수도 있다.
예컨대, 입법은 결혼의 형식요건(marriage formalities), 최소
근로시간(minimum working hours), 선거자금 개혁(campaign finance reform),
도시 중심지의 역사적 보존(historic preservation) 등을 다룰 수 있다.
그러나 실제로는, 누군가 이 입법의 적용이 권리 문제를 야기한다고
지적하면서, 그 입법이 법문에 따라 적용되어야 하는가 아닌가라는 문제로서
결정의 필요성이 발생한다.

이에 대해 사법심사(judicial review)의 일종을 정당화하려는 논증이
존재하며, 나는 그것을 일정 부분 존중한다. 그 논지는 다음과 같다.
입법자가 제출된 입법 제안 속에 어떤 권리 문제가 내포되어 있는지 항상
명확히 인식할 수 있는 것은 아니다. 또한 그 입법이 실제 적용될 때 어떤
권리 문제가 발생할 수 있을지도 항상 예견할 수 없다. 그러므로 시민들이
이러한 권리 문제를 인지하고 공론화할 수 있도록 해주는 메커니즘은
유용하다. 그러나 이는 \textbf{약한 사법심사(weak judicial review)}에
대한 논증일 뿐이며, 법원이 식별한 추상적 권리 문제를 그 법원이
적절하다고 판단하는 방식으로 최종적으로 결정하는 \textbf{강한
사법심사(strong judicial review)}를 정당화하는 논증은 아니다. 이러한
논증은 영국에서처럼 법원이 ``해당 사안은 중대한 권리 문제를 내포하고
있다''는 \textbf{권리선언(declaration)}을 내릴 수 있도록 하는 제도에
대한 옹호라고 볼 수 있다.\footnote{위 주석 26--28 관련 본문 참조.} 혹은,
사법심사가 훨씬 약한 체계에서 볼 수 있는 또 다른 제도에 대한 옹호일 수도
있다. 예컨대, 법무장관(attorney general)이 초당적으로 입법 제안을
검토하고, 그 제안이 제기할 수 있는 권리 문제를 공적으로 지적할 책무를
갖는 경우이다.\footnote{New Zealand Bill of Rights Act 1990, 1990 S.N.Z.
  No.~109, § 7 참조 (``법안이 하원(House of Representatives)에 제출되면,
  법무장관(Attorney-General)은 \ldots{} 가능한 한 신속히 해당 법안에
  포함된 조항이 이 권리장전에 포함된 어떠한 권리 및 자유와도 양립하지
  않을 것으로 보일 경우 이를 하원의 주의에 환기시켜야 한다''). 이 권한
  행사에 대한 논쟁적인 사례로는 Grant Huscroft, The Defeat of Health
  Warnings a Victory for Human Rights? The Attorney-General and
  Pre-Legislative Scrutiny for Consistency with the New Zealand Bill of
  Rights, 14 Pub. L. Rev.~109 (2003) 참조.} 이와 같은 제도는 제3
전제에서 제시된 권리 문제에 대한 민감성(alertness to issues of rights)을
제도화한 한 방식이라 할 수 있다.

\textbf{(p.~1371)} 이제 다음과 같은 상황을 가정해 보자. 입법부는 특정
법안이 제기할 수 있는 권리 문제를 충분히 인식하고 있으며, 그 문제에 대해
숙고한 결과, 토론과 투표를 통해 그 문제를 특정 방식으로
\textbf{해결(resolve)}하였다. 즉, 입법부는 제4 전제에서 상정된 권리에
대한 의견불일치들 중 하나 혹은 그 이상에 대해 입장을 정한 것이다.
그렇다면 이제 우리가 직면하는 질문은 다음과 같다. 입법부의 이 결정이
\textbf{최종적(dispositive)}이어야 하는가? 아니면 그것이 사법부에 의해
재검토되고 심지어는 뒤집힐 수 있는 \textbf{이유(reason)}가 존재하는가?

이 질문에 우리는 어떻게 답해야 할까? 많은 사람들이 이렇게 말하는 것을
들은 적 있다: ``입법부의 결정은 유효하되, 그것이 권리를 침해할 경우는
예외다.'' 그러나 이는 명백히 충분치 않다. 왜냐하면 우리는 이미 사회
구성원들이 특정 입법이 권리를 침해하는지 여부에 대해 서로 의견이
다르다는 것을 전제로 하고 있기 때문이다. 우리는 그 의견불일치를 해결할
수 있는 \textbf{결정 절차(decision-procedure)}를 필요로 한다. 이 점은
토마스 홉스(Thomas Hobbes) 이래로 오래된 통찰이다. 즉, 우리가 결정
절차를 마련하는 이유는 그 절차의 작동이 쟁점을 해결하게 하기 위함이지,
쟁점을 다시 불붙게 하기 위함이 아니다.\footnote{Thomas Hobbes, Leviathan
  123 (Richard Tuck ed., 1996) (1651) 참조.} 따라서 우리가 상정한 사회
구성원들이 권리에 대해 의견불일치하더라도, 그들은 이 의견불일치를 해결할
결정 절차가 정당하다는 이론(theory of legitimacy)을 공유할 필요가 있다.
결정 절차를 구성해야 할 \textbf{이유(reasons)}를 사유할 때, 우리는 그
의견불일치의 어느 한 편에 있는 사람뿐 아니라 양편 모두가 받아들일 수
있는 이유를 고려해야 한다.\footnote{이를 달리 표현하면, 규범적 정치
  이론(normative political theory)은 특정 결정의 당위성에 대한 정당화
  기반 이상을 포함해야 한다는 것이다. 즉, 단순히 정의 이론이나 일반적
  선에 관한 이론을 넘어서, 의견불일치(disagreement)에 직면했을 때 정치적
  결정을 내리기 위해 사용되는 결정 절차(decision-procedures)의
  정당성이라는 규범적 쟁점을 다루어야 한다. 그렇지 않은 규범적 정치
  이론은 본질적으로 불완전하다.}

나는 정당한 결정 절차의 필요성을 \textbf{도덕적 의견불일치(moral
disagreement)}에 대한 대응으로 제시하고 있다. 하지만 어떤 철학자들은
정치에서 의견불일치는 만연하기 때문에, 굳이 그것에 너무 신경 쓸 필요가
없다고 말한다. 우리가 어떤 결과에 대한 정당화만큼이나 정당한 절차에
대해서도 의견불일치하고 있다면, 그리고 (내 설명에 따르면) 우리가 어차피
어떤 입장을 정해야 한다면, 왜 내용(substance) 자체에 대한 입장을 그냥
정해버리면 안 되느냐는 것이다.\footnote{Christiano는 이 점을 절차의
  무한회귀(regress of procedures)라는 방식으로 서술한다: ``Waldron이
  옳다면 우리는 모든 단계에서 의견불일치(disagreement)를 예상할 수 있고,
  따라서 각 분쟁이 발생할 때마다 상위 절차에 의존해야 한다면, 절차의
  무한회귀를 멈출 수 없게 될 것이다.'' Christiano, supra note 59, at
  521. 그러나 Christiano는 이것이 악성 회귀(vicious regress)라는 점을
  증명하려고 시도하지 않는다. 이에 대한 논의는 Waldron, supra note 14,
  at 298--301 참조.} 그러나 이에 대한 대답은 분명하다. 우리는 설령 그
영역에서도 의견불일치가 예상되거나 말거나, \textbf{정당성의 문제(issue
of legitimacy)}로 향해야 한다는 \textbf{(p.~1372)} 것이다. 하나는,
우리는 실제로 결정 절차를 \textbf{설계(design)}해야 하며, 그 설계에
관련된 이유들을 숙고해야 한다. 다른 하나는, 공정성(fairness),
참여(participation), 목소리(voice)와 같이 정당성과 관련된 이유들은 바로
의견불일치가 존재하기 때문에 제기되는 것이며, 의견불일치가 없다면
제기되지도 않았을 이유들이다. 우리가 이에 대해서도 의견불일치할 수는
있지만, 그럼에도 그 이유들을 고려하지 않을 수는 없다. 우리가 그 이유들에
대해 의견불일치할 것이라는 사실 자체가, 그것들을 무시하고 선행되거나
본질적인 의견불일치에 대해 단순히 한쪽 입장을 취하는 방식으로 나아가야
한다는 근거가 되지는 않는다.

그 어떤 결정 절차도 완벽하지 않다. 그것이 사법의 통제를 받지 않는 입법
과정이든, 사법심사 과정이든, 때로는 잘못된 결정을 내릴 수 있으며, 이는
권리를 수호하기보다는 \textbf{배신(betray)}할 수도 있다.\footnote{어떤
  사람들은 오류가 입법부 측에서 더 심각할 가능성이 항상 높다고 말한다:
  입법부는 권리를 실제로 침해할 수 있지만, 법원이 할 수 있는 최악은
  개입하지 않아 권리를 보호하지 못하는 것뿐이라는 것이다. 이는 잘못된
  주장이다. 사법심사의 권한을 행사하는 법원이 권리를 보호하려는 입법을
  무효화함으로써 권리를 침해할 수도 있다. 이 점은 제 IV부 말미에서 더
  논의할 것이다.} 이것은 정치의 현실이다. 모든 사람은 자신이 옳다고
여기는 선택과, 자신이 정당하다고 여기는 결정 절차가 산출하는 선택 사이에
\textbf{괴리(dissonance)}가 발생할 수 있다는 점을 인정해야 한다. 리처드
울하임(Richard Wollheim)은 이를 \textbf{``민주주의 이론의 역설(paradox
in the theory of democracy)''}이라 불렀다.\footnote{Richard Wollheim, A
  Paradox in the Theory of Democracy, in Philosophy, Politics and
  Society 71 (Peter Laslett \& W.G. Runciman eds., 2d ser. 1969).} 한
시민이 ``정책 A는 입법되어서는 안 된다''고 말하면서도, 동시에 ``정책 A는
입법되어야 한다''고 주장할 수 있다는 것이다. 전자는 그가 반대한
정책이고, 후자는 다수결로 선택된 정책이기 때문이다. 하지만 이 역설을
민주주의만의 문제로 환원한 것은 울하임의 오류다. 이것은 \textbf{``무엇이
행해져야 하는가''에 대해 의견불일치가 있을 경우, 그에 대한 결정을 어떻게
내려야 하는가''}에 관한 모든 정치 이론이 겪는 일반적인 역설이다.

이제 그러한 주의를 염두에 두고, 권리(rights)에 관한 의견불일치를
해결하기 위한 \textbf{결정 절차(decision-procedure)}를 설계하거나 평가할
때 고려해야 할 이유들은 무엇일까? 두 종류의 이유가 고려될 수 있다. 나는
이를 각각 \textbf{``결과 관련 이유(outcome-related reasons)''}와
\textbf{``절차 관련 이유(process-related reasons)''}라고 부르겠다. 이 두
가지 이유 모두 결정 절차의 정당성 문제와 관련되어 있다.

\textbf{절차 관련 이유(process-related reasons)}란, 특정한 결과가 어떤
것이어야 하는지에 대한 고려와는 무관하게, 어떤 사람이 그 결정을 내릴
권한을 갖거나 그 결정에 참여할 권리를 가진다고 주장하는 이유를 말한다.
개인의 삶에서 우리는 종종 부모가 자녀에 대해 징계를 할지 말지를 결정할
권리가 있다고 말한다. 이 결정은 길을 지나가던 행인이나 버스의 다른
승객이 대신해서 내려서는 안 된다는 것이다. 우리는 \textbf{(p.~1373)}
아이에게 징계가 필요하다고 판단하기 이전에, 이 결정을 부모가 내려야
한다고 말할 수 있다. 실제로 우리는 행인이 부모보다 더 나은 결정을 내릴
가능성이 있더라도 여전히 이 권한은 부모에게 있다고 말할 수 있다. 정치
영역에서는 이러한 절차 관련 이유들이 평등한 참정권(political equality)과
민주적 투표권(right to vote), 즉 자신의 의견이 타인의 반대에도 불구하고
존중받아야 한다는 권리에서 가장 뚜렷하게 나타난다.

\textbf{결과 관련 이유(outcome-related reasons)}는 그와는 반대로, 올바른
결과(good, just, or right decision)가 산출될 수 있도록 결정 절차를
설계해야 한다는 이유들이다. 지금 우리가 다루는 주제는 \textbf{권리에
대한 의견불일치(rights-disagreements)}이다. 권리가 중요하므로, 우리는
권리에 대한 판단을 정확히 내리는 것이 중요하며, 따라서 결과 관련
이유들을 매우 진지하게 받아들여야 한다. 정책(policy)과 관련된 문제에서는
잘못된 판단이 어느 정도 용인될 수 있지만, 원칙(principle)의 문제에서는
잘못된 판단은 곧 권리의 침해(violation of rights)로 이어질 수 있다.
우리가 상정하는 이 사회의 구성원들은 이러한 결과를 피하거나 최소화하는
것이 얼마나 중요한지를 잘 이해하고 있다.

물론, 권리 문제에 관하여 대립하는 입장을 가진 사람들이 공통적으로 수용할
수 있는 결과 관련 이유를 도출하는 것이 쉬운 일은 아니다. 앞서
언급했듯이,\footnote{위 주석 64 관련 본문 참조.} 결정 절차의 설계는
그것이 해결하려는 특정 의견불일치로부터 독립적이어야 한다. 그렇지 않으면
그 절차 자체가 그 의견불일치를 다시 불붙게 만들 뿐이다. 따라서 우리는
특정 논쟁적 결론을 지지하기 위한 결과 관련 이유는 피해야 한다. 예컨대,
``이 결정 절차는 낙태 권리(pro-choice)를 더 잘 보호해줄 가능성이
크다''고 해서 그 절차를 선택하는 방식은, 낙태 반대자(pro-life
advocate)들의 동의를 얻기 어려울 것이다. 중요한 점은, 그 의견불일치
속에서도 정당성 있게 받아들여질 수 있는 절차를 설정해야 한다는 것이다.

그럼에도 보다 온건한 방식으로 결과 관련 이유를 정당화할 수는 있다. 즉,
논쟁적인 권리 집합을 산출할 가능성이 높은 정치 절차를 선택하자고
주장하는 대신, \textbf{권리에 대한 진실(the truth about rights)}을
밝혀줄 가능성이 더 높은 절차를 선택하자고 주장할 수 있다. 에일린
카바나(Aileen Kavanagh)는 다음과 같이 말한다:

\begin{quote}
``우리는 우리가 어떤 권리를 가지고 있고, 그것이 어떻게 해석되어야
하는지에 대한 정확한 이론 없이도 \textbf{수단주의적 주장(instrumentalist
claims)}을 할 수 있다. 많은 수단주의적 논증은 특정한 권리의 내용에 대한
지식이 아니라, \textbf{(p.~1374)} 입법부와 사법부가 결정을 내리는 방식,
이들의 결정에 영향을 미치는 요인, 그리고 개인이 어느 기관에 어떻게
자신의 청구를 제기할 수 있는지를 포함하는 \textbf{일반적인 제도적
고려사항(general institutional considerations)}에 기초한다.''\footnote{Aileen
  Kavanagh, Participation and Judicial Review: A Reply to Jeremy
  Waldron, 22 L. \& Phil. 451, 466 (2003).}
\end{quote}

이러한 종류의 이유들은 진지하게 고려할 만한 가치가 있다. 조셉
라즈(Joseph Raz)는 더 나아가, 이러한 결과 관련 이유들이 유일하게 고려할
가치가 있는 이유들이라고 주장한 바 있다.\footnote{J. Raz, Disagreement
  in Politics, 43 Am. J. Juris. 25, 45-46 (1998); 또한 다음을 보라:
  Rawls, A Theory of Justice, supra note 53, 230쪽 (``어떤 절차를
  판단하는 근본적인 기준은 그 절차가 산출할 것으로 예상되는 결과의
  정의(justice)이다.'').} 이러한 독단적 주장(dogmatism)은 아마도 해당
쟁점들이 갖는 중요성에서 비롯된 것일 것이다. 권리에 대한 판단은 그만큼
중대한 결과를 낳기 때문이다. 그러나 절차 설계에 있어 우리가 고려해야 할
이유들은 결과 관련 이유만이 아니다. 예컨대
\textbf{자기결정(self-determination)}의 원칙을 생각해보자. 권리에 대한
의견불일치를 해결하는 결정은 각 사회의 정치 체계 내에서 해결되는 것이
바람직하며, 외부의 명령(diktat)에 의해 결정되어서는 안 된다는 이유가
존재한다. 일부는 이것이 결정적(reason-conclusive)인 이유는 아니라고 말할
수 있다. 예컨대 어떤 경우에는 \textbf{국제적 권위(international
authority)}가 인권 문제에 있어 국민자결권(national
self-determination)이나 주권(sovereignty)보다 우선해야 한다고 주장할 수
있다.\footnote{예를 들어 다음을 보라: Louis Henkin, That ``S'' Word:
  Sovereignty, and Globalization, and Human Rights, Et Cetera, Fordham
  University School of Law, Robert R. Levine Distinguished Lecture
  Series 발표문 (1999년 2월 23일), 수록: 68 Fordham L. Rev.~1 (1999).}
하지만 극소수를 제외하면, 이러한 절차 관련 이유가 전혀 중요하지 않다고
부정하는 사람은 없다. 라즈는, 권리 판단의 절차 설계에서 결과 관련 이유가
매우 중요하다는 점은 인정하면서도, 그 외의 이유들도 상당히 중요할 수
있다는 점에는 충분한 주의를 기울이지 않았다.

이처럼 다양한 이유들이 존재한다는 사실을 인식하게 되면, 우리는 이들
이유들의 \textbf{규범적 성격(normative character)}을 검토해야 한다. 이는
각 이유가 서로 어떤 관계를 갖는지를 이해하는 데 중요하기 때문이다.
``결과 관련''이라는 용어는 다소 결과주의적(consequentialist)처럼 들릴 수
있지만, 우리가 피하고자 하는 결과가 \textbf{권리
침해(rights-violations)}이기 때문에, 그 회피에는 권리의 맥락에서
나타나는 \textbf{의무적 긴급성(deontological urgency)}이 부여된다. 물론
그것이 직접적 권리 침해를 금지하는 원칙만큼 강력한 것은 아닐 수 있다.
결정 절차를 설계하는 사람은 그러한 침해에 대해 직접적으로 책임지지는
않기 때문이다. 그러나 여전히 이들은 \textbf{권리에 기반한
책임(rights-based responsibility)}을 지니며, \textbf{주의할 의무(duty to
take care)}를 갖는다.\footnote{특정한 권리(right)가 다양한 파장의
  의무(duty)를 유발한다는 사상은 다음에서 논의된다: Jeremy Waldron,
  Rights in Conflict, 99 Ethics 503, 509--12 (1989).}

\textbf{(p.~1375)} 그렇다면 절차 관련 이유(process-related reasons)의
규범적 성격은 어떠한가? 절차 관련 이유는 종종 의무론적
긴급성(deontological urgency)의 문제이기도 하다. 로널드 드워킨(Ronald
Dworkin)은 참여적 이유(participatory reasons)의 성격을 ``정치적 절차의
참여적 결과(participatory consequences of a political process)''라고
표현함으로써 이를 잘못 기술했다고 생각한다.\footnote{Ronald Dworkin,
  Sovereign Virtue: The Theory and Practice of Equality 187 (2000).}
그는 공동체의 정치적 의사결정(political decisionmaking)에 시민 개개인이
참여할 기회를 갖도록 허용하는 것이 하나의 결과---좋은 결과---를
가져온다고 주장한다. 즉, 그것은 그들이 공동체 내에서 동등한 구성원 혹은
지위(equal membership or standing)를 가진 존재임을 확인시켜 준다는
것이다. 이는 다른 사람들이 그들의 의견(opinions)과 선택(choices)을 가치
있는 것으로 간주한다는 점에서 그들을 안심시키는 효과가 있다. 또한,
사람들에게 참여를 허용하는 것은 정치적 결정(political decisions)의
결과와 시민들이 자신을 동일시하도록 돕고, 그러한 결정들을 어떤
의미에서든 자신의 결정으로 받아들이게 하여 정당성(legitimacy)에 긍정적인
파급 효과(knock-on effects)를 준다(여기서 '정당성'은 사회학적 의미에서의
용어이다).\footnote{Dworkin의 입장에 대한 요약적 진술은 Kavanagh, supra
  note 70, 458--59쪽을 참조하여 작성되었다.} 이러한 설명은 모두 의심할
바 없이 중요하다. 그러나 이는 마치 교장이 학생회(school council)를 통해
교육 문제에 대해 학생들에게 발언을 부여함으로써 얻을 수 있는 이점을
설명하는 듯한 인상을 준다. 드워킨의 설명은 참여할 권리(right to
participate), 즉 한 사회의 의사결정에 있어 평등하게 다뤄져야 한다는
당위(imperative), 그리고 ``이 결정은 나에게 영향을 미치고 나는 그에
구속되는데, 도대체 무슨 권리로 나의 발언을 배제하고---나를 시민권 박탈
상태(disenfranchise)에 놓는가?''라고 분노하며 항의할 때 문제되는
원칙의식(sense of principle)을 근본적으로 과소평가한다.

그렇다면 이제 우리는 절차 관련 이유와 결과 관련 이유를 어떻게 비교
평가할 수 있을까? 이는 흔히 있는 딜레마다. 가령 ``가장 빠르면서 가장
저렴한 차''를 찾는 것처럼, 우리는 두 가지 서로 다른 가치를 동시에
극대화하려 한다. 다음과 같은 방식으로 질문을 설정할 수 있다: ``권리에
대한 진실에 도달할 가능성이 가장 크면서, 동시에 그에 영향을 받는
사람들의 목소리가 \textbf{들릴 동등한 청구권(equal claim to be heard)}을
적절히 존중하는 방법은 무엇인가?''\footnote{Frank I. Michelman, Brennan
  and Democracy 59--60 (1999)에서는 이 문제를 다음과 같이 제기한다.}
또는 반대로, ``영향받는 사람들의 목소리가 들릴 동등한 청구권을 가장 잘
존중하면서도, 권리에 대한 진실에 도달할 현실적 가능성이 있는 방법은
무엇인가?'' 나는 이러한 고르디우스의 매듭(Gordian knot)을 잘라낼 수
있다고 생각한다. 내가 제IV부에서 논증할 바는 다음과 같다. \textbf{결과
관련 이유는 최대한 긍정적으로 평가해도 여전히 결정적이지
않다(inconclusive)}는 것이다. 그들은 중요하지만, 흔히 주장되는 것처럼
사법심사를 정당화하는 명확한 근거를 제공하지는 않는다. 반면, 절차 관련
이유는 매우 일방적이다(one-sided). 이들은 사법심사의 정당성을 훼손하며,
입법에 의한 결정은 온전히 살아남게 만든다. 따라서 나는 어떤 방식으로
질문을 정식화하든 입법 과정이 우위에 선다고 본다. 이것이 바로
\textbf{사법심사에 반대하는 핵심 논증(the core of the case against
judicial review)}이 될 것이다.

\subsection{IV. 결과 관련 이유들 (Outcome-Related
Reasons)}\label{iv.-uxacb0uxacfc-uxad00uxb828-uxc774uxc720uxb4e4-outcome-related-reasons}

Raz에 따르면, ``진행의 자연스러운 방식은, 근본적 권리(fundamental
rights)의 집행(enforcement)을, 해당 시간과 장소의 상황 속에서 그것들을
가장 잘 집행하고 부작용을 최소화할 가능성이 가장 높은 정치적 결정
절차(political decision-procedure)에 위탁하는 것이라고 가정하는
것''이다.\footnote{Raz, supra note 71, 45쪽.} 이 지점에서 논의는
법과정학파(legal process school)가 시작한 법원의 제도적
역량(institutional competence)에 대한 보다 광범위한 논쟁과 연속되어야 할
것이다.\footnote{Hart \& Sacks, supra note 49, 640--47쪽을 보라.}
법원(court)은 어떤 쟁점을 다루는 데에는 능숙하지만, 다른 쟁점에는 그렇지
않다. 기술적으로는, 법원에서 청구(claim)를 제기하려면 권리(right)에 대한
주장을 명시해야 하므로, 법원이 통상적으로 결정하는 쟁점을 지칭하기 위해
``권리(rights)''라는 용어를 사용한다. 그러나 Lon Fuller가 지적했듯이,
그렇다고 해서 법원이 여기서 논의되는 비기술적 의미에서의 권리 청구(claim
of right)를 다루기에 적절한 포럼이라는 결론이 도출되는 것은
아니다.\footnote{Fuller, supra note 49, 368--70쪽.} 일부 권리 청구는
법원이 다루기에 적합한 이진적(binary) 성격을 가지지만, 다른 것들은 보다
다면적(multifaceted) 성격을 지니며, 이는 전통적으로 사법 구조에서
결정되기에 부적합한 것으로 여겨져 왔다. 이 문제는 추가적인 고려를 필요로
하며, 여기서 더 언급하기보다는 이 글의 주요 주제인 도덕적 쟁점들에 관한
법원과 입법기관의 역량에 대한 보다 구체적인 주장으로 논의를 전환하고자
한다.

결과 관련 이유(outcome-related reasons)를 사법심사(judicial review)에
대한 옹호와 연결짓고, 절차 관련 이유(process-related reasons)를
사법심사에 대한 반대와 연관시키는 유혹이 존재한다. 이는 오류이다. 중요한
절차 관련 이유들 중 많은 수가 참여(participation)에 기반을 두고 있으며,
따라서 선출되거나 대표적인 기관을 지지하는 것이 사실이다. 그러나
그렇다고 해서 모든 결과 관련 이유들이 반대 방향으로 논증하는 것은
아니다. 결과 관련 이유들은 실제로 양방향으로 작용한다. 입법기관은 때때로
권리가 방어하고자 하는 유형의 압력에 취약할 수 있는 특성을 지니며, 반면
법원은 권리에 대한 의견불일치(rights-disagreements)가 제기하는 도덕적
쟁점들을 직접적으로 다루는 데 어려움을 겪는 측면도 있다.

Raz는 결과 관련 이유가 양쪽 모두에 무게를 싣는다는 점을 인정한다. 그는
익숙한 방식으로 다음과 같이 논증한다:

\begin{quote}
\textbf{(p.~1377)} ``많은 국가들에서는 입법기관 구성원들이 종파적
이익(sectarian interests)에 의해 상당히 영향을 받아, 어떤 이들이 가지는
권리(rights)가 무엇인지조차 확인하려는 시도를 하지 않을 것이라고 의심할
만한 이유가 충분하다. 우리는 특정 요인들이 판단을 흐리게 할 가능성이
있다는 것을 알고 있다. 예를 들어, 그들은 자신들의 이해관계에 편향될 수
있다. 그러므로 우리는 결정권자가 자신의 결정에 의해 영향을 받지 않거나,
직접적으로 영향을 받지 않는 절차를 선호할 수도 있다. 판단에 영향을
미친다고 알려진 다른 요인들도 존재하며, 그 성격과 존재는 해당 권리의
구체적 내용에 대한 지식 없이도 확인할 수 있다.''\footnote{Raz, supra
  note 71, 46쪽.}
\end{quote}

이러한 비판을 고려할 때, 우리는 그것이 우리의 세 번째 전제와
양립가능성(compatibility)을 가지는지 물어야 한다. 이러한 종파적 편향은
모든 사회의 입법기관에서 전형적인 것인가? 아니면, 권리에 대해 전반적으로
무관심한 사회라는 비핵심 사례(non-core case)와 연결되어야 하는가? 이에
대해서는 제7부에서 더 논할 것이다.\footnote{아래 본문 중 주석 137--141에
  해당하는 텍스트를 보라. 이 부분에서는 (핵심 사례가 아닌 경우에)
  소수자의 권리(rights)에 공감하는 판사가 입법자보다 대중적 편견에
  저항할 위치에 더 낫다는 주장에 대해 논의한다.} 그러나 글자 그대로
받아들인다 해도, Raz의 논변은 일방적인 경향을 갖고 있지 않다. 동일한
종파적 압력은 종종 법원이 권리를 경시하는 이유이기도 하다. 우리는 미국의
Korematsu, Schenck, Dred Scott, Prigg 사건 등에서 이를
보아왔다.\footnote{Korematsu v. United States, 323 U.S. 214 (1944)
  (제2차 세계대전 중 일본계 미국인에 대한 강제 수용을 보호하지 않음);
  Schenck v. United States, 249 U.S. 47 (1919) (제1차 세계대전 중 징병
  비판을 ``극장 안에서 '불이야'라고 외치는 것''과 유사하다고 판시함);
  Dred Scott v. Sanford, 60 U.S. (19 How.) 393, 425--27 (1857); Prigg v.
  Pennsylvania, 41 U.S. (16 Pet.) 539, 612 (1842) (아프리카계 미국인을
  도망노예 사냥꾼으로부터 보호하려는 주 법률을 무효화함).} 더 최근에는,
보통 사법심사의 강력한 옹호자로 알려진 Laurence Tribe조차도 9/11 이후의
미국의 공황 상태를 언급하며 ``시민권과 자유에 대해 우려하는 이들이 위기
상황에서 사법부에 과도한 기대를 거는 것은 심각한 실수일 것''이라
경고하였다.\footnote{Laurence Tribe, Trial by Fury: Why Congress Must
  Curb Bush's Military Courts, New Republic, Dec.~10, 2001, 18, 19쪽;
  또한 다음을 보라: Ronald Dworkin, The Threat to Patriotism, N.Y.
  Rev.~Books, Feb.~28, 2002, 44, 46--47쪽 (위기 상황에서 권리 침해에
  대한 과거 법원의 관용을 지적).}

어쨌든 Raz는 결과 관련 이유가 반대 방향으로 논증할 수 있다는 점도
인정한다:

\begin{quote}
때때로 \ldots{} 한 결정에 의해 자신의 이해관계가 영향을 받지 않을
사람들은 그 상황에서 무엇이 정의로운지에 대해 진지하게 파악하려
\textbf{(p.~1378)} 하지 않을 것이라는 이유가 존재한다. 때때로 사람은
자신이 동일한 계층에 속해 있지 않으면 어떤 계층의 고통을 이해할 수
없으며, 그러므로 그 결정은 영향받지 않는 이들이 아닌, 영향을 받는
이들에게 맡겨야 한다.\footnote{Raz, supra note 71, 46쪽.}
\end{quote}

입법기관은 바로 이러한 이해 능력을 증진하기 위해 대표 구조(structures of
representation)를 갖추고 있다.

민주적 참여 구조가 적절한 결과 확보의 독립적 중요성을 전혀 고려하지 않고
단순히 다수에게 권력을 부여한다고 주장되기도 한다. 이는 터무니없는
주장이다. 모든 민주주의 체제는 일정 수준의 성숙한 판단을 보장하기 위해
참정권(franchise)을 다양한 방식으로 제한한다. 예컨대, 어린이는 해당
결정의 영향을 받음에도 불구하고 투표권이 배제된다. 또한 입법기관은
사회의 다양한 집단에 의해 받아들여질 수 있는 선택지들에 관한 정보가
의사결정 절차에 전달되도록 설계된다. 그리고 일반적으로 의사결정은
양원제(bicameral) 구조 내에서 이루어지며, 각각 다소 상이한 선출 주기를
지닌 두 개의 의회에서 다수의 지지를 받아야만 입법안이
통과된다.\footnote{영국과 같은 일부 양원제 체계(bicameral systems)는
  선출되지 않은 상원(upper house)을 가지며, 하원이 갈등 상황에서 결국
  우위를 점하도록 허용하는 규정(Parliament Acts)과 영국 헌법의
  관습들(conventions)을 포함한다.} 게다가, 약한 사법심사 또는 사법심사가
존재하지 않는 체제에서도 입법 과정에서 권리 쟁점들이 강조되도록 별도의
장치를 마련하는 경우가 있다.\footnote{supra note 63 및 이에 해당하는
  본문을 보라.} 대부분의 민주주의 국가에서는 선거 시기를 중심으로 면밀히
조정된 논의 절차뿐만 아니라, 입법부 내의 공식 토론과 입법부 외부에서
이루어지는 비공식 토론 및 정보 축적 간의 다양한 연결 고리들을 갖추고
있다. 이 모든 것이 민주적 절차에 대한 결과 관련 조정들이다. 참여
측면에서 우리가 보게 되는 것은 Rawls가 '순수한 절차적 정의(pure
procedural justice)'라 부른 것이 아니라, 일종의 '불완전한 절차적
정의(imperfect procedural justice)'에 가깝다.\footnote{Rawls, A Theory
  of Justice, supra note 53, 84--85쪽. 순수 절차적 정의(pure procedural
  justice)는 정의로운 절차를 철저히 따랐다는 사실 자체가 결과의 정의를
  보장한다고 보는 개념이다. 반면, 불완전한 절차적 정의(imperfect
  procedural justice)는 산출된 결과의 정당성을 절차뿐 아니라 그 자체의
  내용에 의해서도 평가해야 한다는 점을 강조한다.}

일반적으로, 사법심사를 지지하는 결과 관련 논변을 읽다 보면, 사람들은
결과 관련 논변은 법원을 지지하는 논변일 수밖에 없다고 \textbf{(p.~1379)}
가정하는 경향이 있다. 그 이유는 사법심사를 반대하는 가장 익숙한 논변들이
비결과 관련적(non-outcome-related)이기 때문이다. 이들은 결과 관련
이유들을 사법부와 연관시키려 무리하게 노력하며, 그 과정에서 사법적
의사결정(judicial decisionmaking)에 대해 현실과 동떨어진 서사를
퍼뜨리기도 한다.\footnote{미국 연방대법원(Justices)의 구성원 간 도덕적
  숙고가 신중하게 이루어진다는 주장에 대해 일반적 비판을 가한 논의로는
  Kramer, supra note 11, 240쪽을 보라. Kramer는 대법관의 정치적 의제와
  각 재판부 내 이념적으로 동기화된 로클럭들(phalanxes of ideologically
  motivated clerks)이 진정한 동료 간 숙고(collégial deliberation)를
  방해한다고 지적한다.} 사법심사 반대자들이 입법기관에 대해 순진하게
낙관적인 시각을 가진 것으로 비판받기도 하지만, 때로 우리는 사법부
옹호자들이 법원의 현실을 솔직히 설명하지 않는 데 대한 대응으로
의도적으로 낙관적 대조를 이루기도 한다.\footnote{Jeremy Waldron, The
  Dignity of Legislation 2 (1999).}

이 절의 남은 부분에서는, 법원을 지지하는 근거로 종종 제시되는 세 가지
결과 관련 이점을 보다 자세히 검토하고자 한다. (a) 권리 쟁점이 특정
사건의 맥락에서 법원에 제기된다는 점, (b) 법원이 권리 쟁점을
권리장전(Bill of Rights)의 문언에 근거하여 다룬다는 점, (c) 사법적
심의(judicial deliberation)에서 추론(reasoning)과 이유
제시(reason-giving)가 핵심적인 역할을 한다는 점이 그것이다. 이 세 가지
모두가 사법심사를 정당화하는 요소로 여겨진다. 그러나 나는 이 세 가지
측면 각각에서, 법원이 권리를 다루는 방식에는 중요한 결과 관련 결함이
있으며, 입법기관 쪽에는 중요한 결과 관련 장점이 존재한다는 점을 논증할
것이다.

\subsubsection{A. 개별 사안에 대한 지향 (Orientation to Particular
Cases)}\label{a.-uxac1cuxbcc4-uxc0acuxc548uxc5d0-uxb300uxd55c-uxc9c0uxd5a5-orientation-to-particular-cases}

사람들은 종종 권리에 관한 사법적 사유(judicial reasoning)가 (입법적
사유{[}legislative reasoning{]}와는 달리) 특별히 뛰어난 점은
권리(right)에 관한 쟁점이 살아 있는 구체적 개인의 상황으로 판사 앞에
제시된다는 점이라고 논증한다. 권리는 결국 개별적 권리(individual
rights)이기 때문에, 어떤 법률 조항이 한 개인에게 어떤 영향을 미치는지를
보는 것은 사고를 집중시키는 데 도움이 된다. 마이클 무어(Michael Moore)의
표현을 빌리자면, ``판사들은 도덕적 통찰(moral insight)을 위한 위치에
입법자보다 더 잘 놓여 있다. 왜냐하면 판사들은 도덕적 통찰에 필요한 세부
정보와 구체적인 개인적 관여를 포함한 도덕적 사고실험(thought
experiment)을 매일같이 접하기 때문이다.''\footnote{Michael S. Moore, Law
  as a Functional Kind, in Natural Law Theory, supra note 19, 188,
  230쪽. 이에 대한 반론은 Waldron, Moral Truth and Judicial Review,
  supra note 19, 83--88쪽을 보라.}

그러나 이는 대부분 신화에 가깝다. 사법심사(judicial review)에 대한
논쟁에서 우리가 주로 말하는 사건들이 고등 항소심(high appellate
levels)에 도달할 때쯤이면, 최초의 살아 있는 권리 보유자(flesh-and-blood
right-holders)의 흔적은 거의 사라진다. 그 시점에서의 논증은
\textbf{(p.~1380)} 해당 권리(right)의 추상적인 쟁점만을 중심으로
순환된다. 원고(plaintiff)나 청원인(petitioner)은 대체로 옹호
단체(advocacy groups)에 의해 선택되며, 이는 그들이 일반적인 공공
정책(public policy) 논증의 일부로 강조하고자 하는 추상적 특성들을
구현하기 위해서이다. 개별 소송 당사자의 특이한 정황이나 개성은 대부분
미국 연방대법원(U.S. Supreme Court)에 사건이 도달할 때쯤이면 사라지며,
대법원은 거의 항상 해당 쟁점을 일반적인 개념으로 논의한다.\footnote{Sarah
  Weddington, Roe v. Wade: Past and Future, Suffolk University Law
  School, The Donahue Lecture Series 발표문 (1989년 12월 7일), 수록: 24
  Suffolk U. L. Rev.~601, 602--03 (1990).}

반면, 입법 과정은 개별 사안들을 고려하는 데 개방적이다. 이는 로비
활동(lobbying), 공청회(hearings), 그리고 입법 토론(debate)을 통해
가능하다. 실제로 오늘날 입법은 악명 높은 특정 개인 사건에 기초해
시작되는 경향을 보이기도 한다. 예컨대 '메건법(Megan's Law)'이 그
예이다.\footnote{Megan의 법(Megan's Law)은 성범죄자 등록제를 도입한
  법으로, Megan Nicole Kanka가 성범죄 전과자에 의해 강간 및 살해된 사건
  이후 1994년 뉴저지에서 제정되었다. 1994 N.J. Laws 1152 (현재 N.J.
  Stat. Ann. §§ 2C:7-1 to 7-11 (2005)로 성문화됨). 연방 수준에서도
  Megan의 법이 제정되었으며, 42 U.S.C. § 14071 (2000)이 그에 해당한다.
  해당 법의 제정 경위에 대한 설명은 다음을 참조하라: Daniel M. Filler,
  Making the Case for Megan's Law: A Study in Legislative Rhetoric, 76
  Ind. L.J. 315 (2001).} 물론, ``어려운 사건(hard cases)은 나쁜 법(bad
law)을 만든다''고 흔히 말하기도 한다. 그러나 이 말이 일정 부분
타당하다면, 특정 개별 사건이 수백만 명에게 다방면으로 영향을 미치는
일반적인 권리 문제에 비추어 어떤 의미를 가지는지 평가함에 있어, 입법부는
훨씬 유리한 위치에 있다고 보는 것이 타당하다고 생각된다.\footnote{Eisgruber,
  supra note 13, 173쪽: ``판사들은 구체적 당사자들 간의 논쟁을 해결하는
  과정에서 헌법적 쟁점을 다루게 된다. 그 결과, 이 쟁점들은 불완전한
  방식으로 법원에 제기된다. 모든 이해관계자들이 법정에 출석할 수 있는
  것은 아니다. 판사들은 제한된 수의 당사자로부터 증거를 받고 주장만을
  듣는다. 따라서 전체 사회·정치·경제 체계의 공정성에 대한 포괄적 관점을
  형성하는 데 필요한 정보를 갖지 못할 수 있다.'' Eisgruber는 이러한
  점에서, 판사들이 포괄적 도덕 원칙에 의존하여 문제를 해결하려는 시도는
  현명하지 못하다고 결론짓는다. 위 저작의 165, 171, 173쪽 참조.}

\subsubsection{B. 권리장전 문언에 대한 지향 (Orientation to the Text of
a Bill of
Rights)}\label{b.-uxad8cuxb9acuxc7a5uxc804-uxbb38uxc5b8uxc5d0-uxb300uxd55c-uxc9c0uxd5a5-orientation-to-the-text-of-a-bill-of-rights}

우리는 지금 권리장전(Bill of Rights)이 존재하는 사회를 상정하고 있으며,
입법에 대한 사법심사(judicial review)가 존재한다면, 그것은 아마도
권리장전을 중심으로 이루어질 것이다. 우리가 이미 가정한 바와 같이,
권리장전은 구성원들이 개인의 권리(individual rights)와 소수자의
권리(minority rights)에 대한 공통된 헌신(commitment)을 바탕으로
채택되었으며, 이들은 해당 권리들이 무엇이고 그것이 어떤 함의를
가지는가에 대해 의견불일치(disagreement)를 갖고 있음에도 불구하고 이를
수용하였다. 이제 입법과 관련하여 권리에 대한 의견불일치가 발생할
\textbf{(p.~1381)} 경우, 해당 쟁점이 제기되는 결정
절차(decision-process)에서 권리장전이 어떤 역할을 해야 하는지가
문제된다. 결과 관련 관점(outcome-related point of view)에서 보았을 때,
권리장전의 문언에 기반하여 권리에 대한 의견불일치를 다투는 것이 바람직한
일일까, 아니면 그렇지 않은 일일까?

그것이 바람직하다고 생각할 수 있는 한 가지 이유는, 권리장전의 서면
표현들이 쟁점이 되는 추상적인 권리 문제들에 대해 당사자들이 보다 명확히
집중할 수 있도록 도와줄 수 있기 때문이다. 그러나 반대편에도 강력한
이유들이 존재한다. 권리장전에 사용된 표현 방식은 권리에 대한
의견불일치를 고려하여 선택된 것이 아닐 수 있다. 혹은, 그러한
의견불일치가 존재하였을 때, 그것을 회피(finesse)하기 위해 그러한 표현이
선택되었을 수도 있다. 그러한 상투적 표현들은 권리에 대한 의견불일치를
냉철하고 책임 있게, 그리고 성실하게 탐구하려는 시도를 오히려 방해하는
잘못된 출발점일 수 있다.

권리장전의 서면 표현들은 또한 일정한 경직된 문언 형식주의(textual
formalism)를 조장하는 경향이 있다.\footnote{이 주장은 Waldron, A
  Right-Based Critique, supra note 14에서 내가 전개한 것이다.}
권리장전의 보호를 받는 법적 권리(legal right)는 그 조항들이 명시된
정형화된 문구(canonical form of words) 아래에서 그 보호를 받는다. 미국
헌법 해석의 경험에서 얻을 수 있는 하나의 교훈은, 각 조항의 문구가
독자적인 생명력을 가지게 되어, 해당 권리에 대해 말하고자 하는 모든
내용을 반복적으로 표현하는 강박적 표현(catchphrase)으로 작용하게 된다는
것이다. 입법우위(legislative supremacy) 체제에서는 이러한 위험이 비교적
적은데, 이는 입법자들이 권리장전의 표현을 참조하지 않고서도 스스로
문제를 제기할 수 있기 때문이다. 그러나 법원은 그 작동 방식상 사유의
정당화를 위한 문언적 근거(textual haven)를 추구하기 때문에, 권리장전의
문언에 집착하는 방식으로 사고를 조직하게 되는 경향이 있다.

최소한, 법원은 권리에 대해 논증할 때, 권리장전과 같은 문헌을 어떻게
접근해야 하는지에 대한 부차적인 논증에 휘말리게 될 것이다. 미국의 경험은
이를 뒷받침한다. 대부분의 사법 의견에서 해석 이론(theories of
interpretation)에 대한 논의의 비중이 도덕적 쟁점에 대한 직접적인 논의의
비중보다 훨씬 높으며, 실제 쟁점들이 중요하다고 생각하는 사람이라면
누구도 이를 만족스럽다고 여길 수 없다. 이는 부분적으로, 사법심사 자체의
정당성(legitimacy)이 문제적이라는 점 때문이다. 판사들 또한 우리와
마찬가지로 자신들이 이러한 쟁점을 결정하는 절차가 정당한가에 대해
우려하며, 도덕적 이유에 대한 직접적인 논의로 나아가기보다는 그들을
정당화해주는 문언(authorizing texts)에 집착하고, 그 해석을 둘러싼 논쟁에
매달리게 된다.\footnote{Tushnet, supra note 11, 60쪽: ``법원은
  본질적으로 헌법적 가치보다는 자신의 제한된 능력을 반영하는 방식으로
  일부 법리를 설계할 수 있다.''}

\textbf{(p.~1382)} 마지막으로 하나 더 지적할 점이 있다. 권리장전의
문언은 포함하고 있는 내용뿐 아니라 생략하고 있는 내용을 통해서도 사법적
사유(judicial reasoning)를 왜곡할 수 있다. 특정 사회의 구성원들이
권리장전이 부정적 권리(negative rights, 자유권)뿐 아니라 긍정적
권리(positive rights, 사회경제적 권리)도 포함했어야 했는지에 대해
의견불일치를 갖고 있다고 해 보자.\footnote{Jackson v. City of Joliet,
  715 F.2d 1200, 1203 (7th Cir. 1983) (Posner, J.): 미국 헌법 체계는
  ``긍정적 자유(positive liberties)``가 아니라 ``부정적 자유(negative
  liberties)``의 헌장이라는 관찰을 담고 있음. 이에 대한 반례로는 다음을
  보라: Mark Tushnet, An Essay on Rights, 62 Tex. L. Rev.~1363, 1393--94
  (1984): ``우리는 물론 다른 헌법을 가질 수도 있었을 것이다\ldots{}
  인간성의 당파(the party of humanity)는 권리의 언어를 재구성하기 위해
  노력해야 하며, 그 결과 Posner 판사의 묘사가 자연스럽지 않게 들리고,
  심지어 부적절하게까지 보일 수 있도록 해야 한다는 주장이 가능하다.''}
긍정적 권리가 포함되었어야 한다고 생각하는 이들은, 현재의 권리장전이
그것을 생략함으로써 도덕적 사유를 왜곡하고 있다고 생각할 수 있다. 이에
대한 반론으로, 최악의 경우 그 생략은 적절하게 사법심사가 이루어져야 할
입법에 대해 그것이 이루어지지 않는 결과를 초래할 뿐이며, 그것 자체로
사법심사 일반에 대한 비판은 아니라고 주장될 수 있다. 그러나 그것은
지나치게 단순한 주장이다. 긍정적 권리의 부재는 포함된 권리들에 대한
판사의 이해를 바꾸거나 왜곡할 수 있다. 예컨대, 만약 재산권(property
rights)이나 계약의 자유(freedom of contract)가 명시적인 복지권(welfare
rights)과 병렬적으로 제시되었다면, 판사들은 그것들에 지금보다 덜 중대한
무게를 부여했을 수도 있다. 그러나 그러한 병렬 구성이 존재하지 않을 경우,
판사들은 재산권이나 계약의 자유에 과도한 무게를 부여할 수 있으며, 그
결과 권리장전의 문언에 담기지 않은 권리들을 입법을 통해 실현하려는
입법의도를 가진 법률들을 위헌으로 무효화하게 되는 일이 발생할 수 있다.

\subsubsection{C. 이유 진술하기 (Stating
Reasons)}\label{c.-uxc774uxc720-uxc9c4uxc220uxd558uxae30-stating-reasons}

개별적 권리(individual rights)의 문제에 관한 사법적 결정(judicial
decisionmaking)의 가장 큰 장점은 그에 수반되는 명시적 사유
제시(reasoning and reason-giving)라고 종종 여겨진다. 법원(courts)은
자신들의 판결에 대해 이유를 제시하며, 이는 쟁점의 중요성을 진지하게
받아들이고 있다는 표시이고, 반면에 입법부(legislatures)는 그렇지 않다는
것이다. 하지만 이는 잘못된 대비이다. 입법자들도 판사들과 마찬가지로
자신들의 투표에 대해 이유를 제시한다. 우리는 그것을 '토론(debate)'이라고
부르며, 그 내용은 『한자드(Hansard)』나 『미국 의사록(Congressional
Record)』에 공개된다. 차이점은 법률가들이 판사의 사유(reason)에 대해
정밀한 연구를 하도록 훈련받는 반면, 입법부의 논증에는 그렇게 훈련되어
있지 않다는 데 있다. (물론 해석적 목적을 위해 입법부 논의 내용을
활용하는 경우는 있다.)

어쩌면 이 논점은 사유 제시의 유무가 아니라 그 질(quality)에 관한 것일
수도 있다. 그러나 내 생각에는, 입법에 대한 사법심사(judicial review)를
수행할 때 법원이 제시하는 이유들은, 이상적 정치적 숙의(deliberation)에서
다루어질 법한 완전한 논의에서 검토될 이유들과는 상당히 다르며, 그 이유를
탐색하고 인용하고 평가하고 비교하는 \textbf{(p.~1383)} 방식도 이상적
정치적 숙의자(ideal political deliberator)의 그것과는 매우 다르다. 이는
앞서 언급된 지점과도 관련이 있다. 즉, 그 이유들은 권리장전(Bill of
Rights)의 용어에 맞춰져 있다는 점이다. 만약 우리가 훌륭하고 최신의
권리장전을 가지고 있다면, 사법적 사유(reason-giving)와
합리적·도덕적·정치적 숙의에서 우리가 기대하는 사유 제시 간에 일정한
일치가 있을 수도 있다. 하지만 수백 년 된 구시대 헌법을 가지고 있다면,
이른바 사유 제시는 인위적이고 왜곡된 것이 되기 쉽다. 미국에서는 '사유
제시(reason-giving)'라고 불리는 것이 대체로 법원이 직면한 판결을 18세기
혹은 19세기의 사고가 미치지 못한 낡은 문장들과 연결시키려는 시도일
뿐이다. (예컨대, ``실질적 적법절차(substantive due process)''라는 개념이
모순어법(oxymoron)인지 아닌지를 따지는 논의가 노동법이나 낙태권(abortion
rights)을 논하는 데 적절한 틀인가?)

법원의 사유 제시는 또한 현재의 결정을 과거의 유사한 결정들과 연결하거나
구별하려는 필사적인 유비 혹은 비유비의 구성을 포함한다. 비록 최고 수준의
법관들도 판례(precedent)가 문제를 확정 짓지는 못한다고 인정하면서도,
여전히 판례에 대한 힘겨운 논의가 지속된다.\footnote{예: Henry Paul
  Monaghan, Stare Decisis and Constitutional Adjudication, 88 Colum. L.
  Rev.~723 (1988).} (그와 동시에, 어떤 상황에서 판례를 유지하거나
무시해야 하는가에 대한 논의도 끊임없이 이어진다.\footnote{예: Planned
  Parenthood of Se. Pa. v. Casey, 505 U.S. 833, 854--69 (1992) (헌법적
  선례를 뒤집을 수 있는 조건에 대해 논의함).}) 그러는 사이, 권리에 관한
선의의 의견불일치(good faith disagreement)에서 실제로 쟁점이 되는
핵심적인 문제들은 주변부로 밀려난다. 보통 그러한 핵심 쟁점들은 수십 쪽의
판결문 중 한두 단락 정도만을 차지하며, 그나마도 직접적으로 다루어지는
일은 드물다. 예컨대, \emph{로 대 웨이드(Roe v. Wade)} 판결문은 50쪽이
넘는데, 그 가운데 재생산의 자유(reproductive rights)와 사생활(private
life)의 관계에 관한 도덕적 중요성을 논하는 부분은 몇 단락에 불과하며, 또
하나의 도덕적 쟁점인 태아(fetus)의 권리 지위(rights-status)에 대한 몇
단락은 단지 그 쟁점에 대한 다양한 입장들이 존재함을 보여주는 데 할애되고
있다.\footnote{해당 판결에는 방대한 법률적·사회적 역사 서술이 담겨
  있으나, 실제로 쟁점이 되는 도덕적 문제에 대한 논의는 몇 페이지만을
  차지한다. Roe v. Wade, 410 U.S. 113, 153--55 (1973) (사생활과 재생산
  권리의 중요성에 대해 논의); id. at 159--62 (태아의 권리나 인격성에
  대한 주장에 대해 논의).} 그 단락들을 읽어보라. 결론이 매력적으로 보일
수는 있지만, 그 '논증(reasoning)'은 매우 빈약하다.

나는 이것이 나름 정당한 이유를 지니고 있다고 본다. 법원은 자신들의
결정의 정당성(legitimacy)에 대해 우려하고 있으며, 그래서 자신들의 '사유
제시(reason-giving)'를 \textbf{(p.~1384)} 헌법(constitution),
법률(statute), 혹은 판례(precedent)에 의해 자신들이 법적으로 권한을 갖고
있음을 보여주는 사실들에 집중시키는 것이다. 이는 이해할 수 있는
행위이다. 하지만 입법보다 사법심사가 바람직하다는 결과 중심
논증(outcome-related argument)에서는 이러한 점이 법원에 불리하게
작용한다.\footnote{Eisgruber는 이를 사실상 인정한다: ``{[}판사들은{]}
  논쟁적인 판결을 정당화하기 위해 모호한 선례들을 인용하고, 그들의
  진정한 이유를 이전 판례로부터 인용된 불분명한 문구들로 감추려는 시도를
  자주 한다.'' Eisgruber, supra note 13, 70쪽; 또한 135쪽 참조:
  ``{[}판사들은{]} 종종\ldots{} 자신들이 정치적 판단을 내리고 있지
  않다고 가장하며, 그들의 판결이 단지 문언적 세부사항이나 역사적 사실에
  의해 강요되었다고 주장한다.''} 정당성 문제에 주의를 빼앗긴 법원은 다른
법원이 무엇을 했는가, 권리장전의 문구가 무엇인가에 집중하는 반면,
입법자들은---비록 많은 결함이 있을지라도---문제의 핵심으로 직접
들어가려는 경향을 보인다.\footnote{Mark Tushnet이 강조했듯이, 우리는
  입법자가 판사처럼 추론하지 않는다고 비판해서는 안 된다. 그것이 문제를
  해결하는 현명한 방식이 아닐 수도 있기 때문이다. Tushnet, supra note
  11, at 63--65.}

이 점에서, 사법심사가 없는 국가에서 권리에 관한 중대한 쟁점을 다룰 때
입법 토론에서 보여지는 풍부한 논증의 깊이는 주목할 만하다.\footnote{Waldron,
  supra note 46에서 일부 내용을 수정하여 인용함.} 나는 최근 1966년 영국
하원의 \emph{임신의 의료적 종료 법안(Medical Termination of Pregnancy
Bill)}에 대한 토론을 읽은 적이 있다.\footnote{영국 의회에서는
  제2독회(second reading) 토론이 법안의 주요 원칙들에 대한 심의가
  이루어지는 단계이다.} 이는 낙태법을 완화하려는 법안이었으며, 그
제2독회 토론은 정치 제도가 도덕적 문제와 씨름하는 가장 훌륭한 사례 중
하나였다. 『한자드(Hansard)』에 100쪽에 달하는 토론이 수록되어
있으며,\footnote{732 Parl. Deb., H.C. (5th ser.) (1966) 1067,
  1067--1166.} 친생명(pro-life) 노동당 의원들과 선택 지지(pro-choice)
노동당 의원들, 친생명 보수당 의원들과 선택 지지 보수당 의원들이 낙태를
둘러싼 모든 쟁점들을 진지하게 토론했다. 그들은 권리, 원칙, 실용적 쟁점
모두에 주의를 기울이며 열정적이면서도 철저하고 명예롭게 논쟁하였다. 결국
법안 지지파가 승리했고, 선택 지지 입장이 우위를 점하였다.\footnote{물론
  제2독회 토론이 끝은 아니었다. 장시간의 위원회 단계와 제3독회 토론,
  그리고 상원(House of Lords)에서의 유사한 토론이 이어졌다. 그러나
  궁극적으로는 자유화 입법이 통과되었다.} 인상 깊은 점은, 토론에 참여한
모든 이들---심지어 표결 결과를 예상한 친생명 의원들조차---자신들의
입장이 존중받으며 \textbf{(p.~1385)} 경청되었다는 점에 경의를 표했다는
사실이다.\footnote{예: 732 Parl. Deb., H.C. (5th ser.) (1966) 1152.
  가톨릭 신자이자 법안에 반대 투표를 한 Norman St.~John-Stevas 의원은
  ``이것이 의회의 가장 고귀한 전통에 걸맞는 수준에서 진행된 매우 중요한
  토론이었다는 데에는 우리 모두가 동의할 것이다''고 말하며 자신의 주장을
  시작하였다. 이어 그는 법안을 발의한 의원에게 ``법안을 소개한 방식이
  매우 절제되고 능숙했음을'' 치하하였다. Id.} 생각해 보라. 미국에서
낙태에 반대하는 입장을 가진 이가 \emph{로 대 웨이드} 판결에서 자신의
입장이 정중하게 경청되었고 성실히 다루어졌다고 공공연히 말한 적이
있었는가?\footnote{이 예시를 언급하면 미국 친구들은 영국 의회는
  미국에서는 불가능한 형태의 토론이 가능하도록 조직되어 있다고 말하곤
  한다. 그러나 미국이 병리적 사례인지 여부를 떠나, 이는 사실이 아니다.
  위 토론은 하원(House of Commons)이 권리 문제에 한해서 정당 기강(party
  discipline)을 일시적으로 중단했기 때문에 가능했던 것이다. 의원들은
  실제로 자신들의 의견을 명확하고 강력하게 표현하는 방식으로 토론을
  벌였으며, 이는 오히려 미국식 의회 운영에 가까운 것이었다.}

미국에서는 낙태권과 같은 개별적 권리 문제를 법원에 위임하여 헌법적
판단을 받게 하는 것을 자축하며, 법원을 '원칙의 포럼(forum of
principle)'으로 간주하는 경향이 있다. 이는 로널드 드워킨(Ronald
Dworkin)의 유명한 표현이다.\footnote{Dworkin, supra note 3, at 33,
  69--71.} 우리는 때때로 영국인들이 이러한 방식으로 하지 않는 것을 두고
시대에 뒤떨어졌다고 말하기도 한다.\footnote{Editorial, Half-Measures on
  British Freedoms, N.Y. Times, Nov.~17, 1997, at A22 (Human Rights
  Act가 영국을 강력한 사법심사 체제로 완전히 전환하지 못한 점을 비판함).}
그러나 영국의 입법 토론과 미국의 사법적 사유 사이의 핵심 차이는, 후자가
주로 해석과 교리에 집중하는 반면, 전자는 낙태 그 자체와 그에 수반되는
문제들---태아의 윤리적 지위, 임신한 여성의 상황, 그들의
선택·자유·사생활의 중요성, 그 모든 것에 따르는 도덕적 갈등과 실천적
문제들, 그리고 사적 도덕 문제에 대해 법이 어떤 역할을 해야 하는가---를
중심으로 집중하여 논의한다는 점이다. 낙태권에 관하여 사회가 결정을
내려야 할 때 반드시 논의되어야 하는 쟁점들은 바로 이러한 것이며, 입법
토론에서는 이 문제들이 가장 많이 다루어지고, 사법적 심의에서는 가장 적게
다루어진다.\footnote{Elena Kagan 등은, 법원이 권리에 대해 논의하는
  방식에 대한 본 비판은 우리가 보호하려는 것이 도덕적 권리(moral
  rights)라는 전제에 기반한 것이라고 지적하였다. 반면, 우리가 보호하고자
  하는 것이 헌법적 법적 권리(legal constitutional rights)라면, 이와 같은
  논의 방식은 그다지 부적절하지 않다는 것이다. 그러나 나는 이에 동의하지
  않는다. 우리가 보호하고자 하는 것은 권리이며, 현대 국가 내에서 이를
  가장 잘 보호하고 지적으로 논의할 수 있는 메커니즘이 무엇인지가 문제다.
  내가 주장하는 바는 도덕철학 세미나의 담론 방식을 이상화하자는 것이
  아니다. 다만 우리가 겪는 의견불일치에서 실제 인간의 이익과 자유라는
  문제에 어떻게 접근할 것인지가 중요하다는 점을 강조하고자 한다.
  법률주의적인 방식이 그것에 적합할 수도 있고 아닐 수도 있지만, 법률주의
  자체를 목적으로 삼아야 한다고 말하는 것은 잘못된 태도일 것이다.}

\textbf{(p.~1386)} 결과 중심 논증(outcome-related argument)에 관하여 더
논의할 수 있는 여지는 분명 있다. 분명 법원이 권리 쟁점을 다루는 방식이
실제 쟁점의 본질을 왜곡하는 것과 마찬가지로, 입법의 논증도 당황하거나,
경솔하거나, 단지 인기 있는 혹은 분파적 구호를 반복하는 방식으로 진행되어
수치스러운 경우도 있다. 문제는 다음과 같다. 이러한 심의의 결함들 가운데
어떤 것들이 해당 제도(법원이나 입법부)의 정상적 행위로 간주되어야 하며,
어떤 것이 일탈(aberration)로 간주되어야 하는가? 드워킨의 ``원칙의
포럼''이라는 수사에도 불구하고, 나는 법원이 전례(precedent), 문헌(text),
교리(doctrine), 기타 법형식주의적 요소들에 집중하는 방식으로 행동할
것이라 기대되는 기관이라고 본다. 우리가 법원에 대해 갖고 있는 두 번째
가정(assumption two)은 법원이 바로 그런 방식으로, 즉 법형식주의적 기준에
따라 잘 행동하는 기관이라는 것이다. 반면 입법부의 경우, 경솔하거나
분파적인 입법 행위는 입법부가 수행하도록 설계된 정상적 이론의 일부가
아니다. 그것은 권리를 존중하는 시민이 다수인 사회에서 입법 결정의 핵심
사례(core case)에 대해 우리가 가져야 할 기본적인 전제가 아니다. 물론
어떤 국가들---예컨대 미국---에서는 입법의 병리적 현상이 발생했을 수
있다. 만약 그것이 사실이라면, 미국인들은 사법심사에 대한 비핵심 사례
논증(non-core argument)을 자신들의 예외적 상황에 한정해야 할 것이다.

\subsection{V. 절차 관련 이유 (Process-Related
Reasons)}\label{v.-uxc808uxcc28-uxad00uxb828-uxc774uxc720-process-related-reasons}

우리가 어떤 방식으로 의사결정 절차(decision-procedures)를 구성할지를
결정하는 데 있어 고려하는 이유들 가운데 일부는, 특정 결과나 일반적인
결과와는 거의 무관하다. 이러한 이유들은 오히려 \emph{목소리(voice)},
\emph{공정성(fairness)} 또는 그 밖의 절차 그 자체의 측면들과 관련된다.
앞서 언급했듯, 절차 관련 논거(process-related arguments)는
사법심사(judicial review)에 반대하는 방향으로 명확히 작용한다고 흔히
간주된다. 이는 완전히 사실은 아니다. 이 관행을 옹호하는 사람들 중 일부는
미약한 절차 관련 논거들을 꾸며내기도 했으며, 나는 본 절의 마지막에서
그것들을 검토할 것이다. 그러나 대부분의 경우 그 반대는 사실이다. 절차
관련 이유들의 압도적 다수는 입법부(legislatures)를 지지하는 방향으로
작용한다.

결과에 대한 의견불일치(disagreement)가 존재하는 상황에서 의사결정 절차의
정치적 정당성(political legitimacy)을 묻는 문제는 다음과 같이 제기될 수
있다. (이 논의는 다소 추상적일 수밖에 없음을 양해 바란다.)

\textbf{(p.~1387)} 우리는 특정한 절차에 따라 결정이 이루어지는 상황을
상상하고, 그 결정에 의해 구속되거나 부담을 지게 될 시민 Cn이 해당 결정에
동의하지 않으며, 왜 자신이 그것을 받아들여야 하는지, 왜 준수하거나
감내해야 하는지를 묻는 장면을 상상한다. 그 결정을 지지하는 사람들 중
일부는, 그 결정이 실질적 측면에서 정당하다고 설득하려 할 수 있다. 그러나
그들은 설득에 실패할 수도 있는데, 이는 Cn이 고집스럽기 때문이 아니라,
단지 이 중대하고 난해한 문제에 대해 Cn이 여전히 (그리고 합리적으로) 다른
견해를 유지하고 있기 때문이다. 그렇다면 이제 Cn에게 무엇을 말할 수
있을까? 그녀에게 제시될 수 있는 그럴듯한 응답은 그 결정이 도달된 절차에
관한 것일 수 있다. 비록 결과에는 동의하지 않더라도, Cn은 그것이 공정하게
도달된 것이라고 받아들일 수 있을지도 모른다. 이러한 절차 중심의 응답이
기반하는 이론이 바로 정치적 정당성 이론(theory of political
legitimacy)이다.

정치적 의사결정 절차들은 대체로 다음과 같은 형식을 취한다. 특정 결정에
대해 의견불일치가 있을 경우, 해당 결정은 지정된 일련의 개인들 \{C1, C2,
\ldots, Cm\}에 의해, 지정된 절차에 따라 이루어지게 된다. 정당성 이론이
감당해야 할 과제는, 왜 바로 이들 개인들, 그리고 다른 누군가가 아니라,
이들이 그러한 의사결정에 참여할 특권을 가져야 하는지를 설명하는 것이다.
Cn의 표현을 빌리자면, ``왜 그들인가? 왜 내가 아닌가?''라는 물음에 대해,
정당성 이론은 그에 대한 설명의 기초를 제공해야 한다. 이 문제는 일반적인
것이며, Cn 개인의 특이한 완고함 때문이 아니므로, 동일한 논변이 포함되지
않은 다른 시민들---Co, Cp, 기타 모든 C들---의 유사한 질문에도 적용되어야
한다. 그러나 설령 이 응답이 받아들여진다고 해도 문제는 끝나지 않는다.
정당성 이론은 Cn이 제기할 수 있는 또 다른 질문에도 답해야 한다. 즉, ``그
결정 절차에서, 왜 나와 같은 견해를 가진 결정자들의 의견에 더 큰 비중이
주어지지 않았는가?''라는 질문이다. \{C1, C2, \ldots, Cm\}의 구성에 대한
정당화뿐 아니라, 그들이 사용한 결정 절차 자체에 대한 방어가 필요하다.

이제 이와 같은 추상적인 대수(algebra)를 구체화해 보자. 어떤 시민이
권리에 관한 입법적 결정에 동의하지 않으며, 위에서 상정한 두 가지 질문을
제기한다고 하자. 그녀는 다음과 같이 묻는다. (1) ``나와 2억 5천만 명의
사람들에게 영향을 미치는 권리 문제에 대해, 왜 이 대략 500명 정도 되는
남녀(입법부 구성원들)가 결정을 내릴 특권을 가져야 하는가?'' 그리고 (2)
``비록 이 500명의 특권을 인정한다고 하더라도, 왜 나와 같은 견해를 지닌
입법자들의 의견에 더 많은 비중이 주어지지 않았는가?''

민주주의 체제에서 입법부는 위 두 가지 질문에 대해 비교적 설득력 있는
응답을 제공하는 방식으로 구성된다. 첫 번째 질문에 대한 응답은 입법부
구성에 이르는 공정한 선거(fair elections)에 관한 이론에 의해 제공된다.
그 선거에서는 Cn과 같은 사람들이 다른 모든 시민들과 동등하게
대우받았으며, 이와 같은 유형의 결정에 참여할 소수자들을 누구로 할 것인지
결정하는 데 기여하였다. 두 번째 질문에 대한 응답은 다수결 결정
원칙(majority decision, MD)의 정당성을 뒷받침하는 잘 알려진 공정성
논변들(fairness arguments)에 \textbf{(p.~1388)} 의해 제공된다. 나는
여기서 이를 옹호하는 임무를 수행할 필요는 없다. 다수결 규칙에 대한
공정성/평등성 방어는 익히 알려져 있다.\footnote{사회선택이론(social
  choice theory)에서 단순다수결(MD, majority decision)이 공정성, 평등성,
  합리성의 기본 요건을 충족하는 유일한 방식이라는 정리에 관해서는 다음
  문헌들을 보라: Amartya K. Sen, Collective Choice and Social Welfare
  71--74 (1970); Kenneth O. May, A Set of Independent Necessary and
  Sufficient Conditions for Simple Majority Decision, 20 Econometrica
  680 (1952). 또한 다음 논의들도 유익하다: Charles R. Beitz, Political
  Equality 58--67 (1989); Roberta Dahl, Democracy and Its Critics
  139--41 (1989).} 다수결 결정 원칙(MD)은 다른 어떤 규칙보다도 논쟁 중인
결과들 사이에서 중립적이며, 참가자들을 평등하게 대우하고, 각 의견에 대해
가능한 최대의 비중을 부여하되 모든 의견에 동일한 비중을 주는 데 필요한
양립가능성(compatibility)을 유지한다. 우리가 바람직한 결과에 대해
의견불일치를 겪고 있으며, 사전에 어느 쪽으로도 편향되지 않기를 원하고,
관련된 모든 참가자가 그 절차에서 동등하게 대우받아야 할 도덕적
청구권(moral claim)을 지닌다고 할 때, 다수결 결정 원칙(MD)---혹은 그와
유사한 어떤 것---이 우리가 사용해야 할 원칙이다.\footnote{Ronald
  Dworkin은 개인적 대화에서, 정의(justice)의 제1차적 사안에 대해
  단순다수결(MD)을 적용하는 것은 적절하지 않다고 나를 설득하였다. 만약
  우리가 과밀한 구명보트에 있고 누군가가 보트에서 나가야 하는
  상황이라면, 그 사람을 누구로 할지를 다수결로 결정하는 것은 적절하지
  않다. 그러나 일반적인 규칙들 사이의 선택에서는 MD가 적절한 원칙이 될
  수 있다. 예를 들어 누군가는 제비뽑기를 제안하고 다른 사람은 가장 나이
  많은 사람이 떠나야 한다고 제안할 때, 이러한 규칙들 중에서 선택할 때는
  MD가 공정한 방식일 수 있다.}

그러면 누군가 다음과 같이 반응한다면 어떠한가: 나는 Cn과 같은 개별
시민들이 자신들에게 영향을 미치는 문제에 관한 의사결정 과정에서 동등하게
대우받을 \emph{권리(right)}를 갖는다는 점은 이해할 수 있다. 그러나 왜
입법부의 500명 대표들은 이 과정에서 동등하게 대우받을
\emph{권리(right)}를 갖는가? 그들이 MD를 사용할 정당성은 어디에서
비롯되는가?

그에 대한 대답은 입법부의 경우 제1질문과 제2질문에 대한 응답 간의
연속성에 있다. 입법부에서는 대표자를 선출하기 위해 다수결의 일종을
사용하고, 대표자들 사이의 의사결정에도 다수결의 일종을 사용한다.
이론적으로, 이 두 가지는 함께 시민 전체 사이의 의사결정 절차로서 다수결
원칙을 사용하는 것의 합리적인 근사값을 제공하며, 따라서 다수결 원칙의
기초에 놓인 가치들을 전체 시민에게 적용하는 것의 합리적인 근사값을
제공한다.

일반적으로, 우리가 Cn에게 전하려는 말은 대략 다음과 같다: 이와 같은
의사결정 절차에 대해 도전하는 것은 Cn 혼자가 아니다. 실제로는 수백만
명의 개인들이 같은 질문을 던진다. 그리고 우리는 각 개인에게 그 지점을
인정하고, 그 의사결정에서 일정한 \emph{발언(a say)}을 제공함으로써
응답한다. 사실 우리는 가능한 한 많은 \emph{발언(a say)}을 제공하려고
한다. 물론, 이는 다른 모든 시민들의 \emph{목소리가 들릴 동등한
청구권(the equal claim to be heard of the voices)}을 고려하는 동일한
논변 구조에 따라 공정하게 응답하려는 우리의 시도에 의해 제한된다. 우리는
각 개인에게, \textbf{(p.~1389)} 다른 모든 사람들과
\emph{양립가능성(compatibility)}을 유지하면서도 가장 많은 \emph{발언(a
say)}을 제공하고자 한다. 이것이 우리의 원칙이다. 그리고 우리는 우리의
복잡한 선거 및 대표 제도가 정치적 평등---즉, 동등한 목소리와 동등한 결정
권한(decisional authority)---에 대한 요구를 대략적으로 만족시킨다고
믿는다.

물론 현실 세계에서, 선거, 대표, 그리고 입법 절차를 통한 정치적 평등의
실현은 완전하지 않다. 선거 제도는 종종 결함을 갖고 있다(예: 선거구 경계
설정 방식의 불만족스러움, 선거구 간의 비례성 부족 등). 입법 절차 역시
결함이 있을 수 있다(예: 입법부 내의 연공서열 시스템이 공정성을 해치는
경우). 이러한 모든 점들은 인정 가능하다. 하지만 우리의 제1전제를
기억하라: 평등성과 공정성이라는 민주적 가치와 관련하여 대체로 양호한
상태에 있는 일련의 입법 제도---입법부 선거 제도 및 내부 의사결정 제도를
포함하는---가 존재한다고 가정하는 것이다. 또한, 입법자들과 그들의
유권자들이 이러한 제도들이 해당 원칙에 부합하는지를 검토하며 지속적으로
감시하고 있다고 가정한다. 예컨대, 많은 민주국가에서는 비례대표제, 선거구
구획, 입법 절차 등에 대해 다양한 대안들이 논의된다. Cn은 이 제도들이
완전하지 않으며, 마땅히 이루어졌어야 할 만큼 개혁되지 않았다고 불평할 수
있다. 그러나 (현실 정치체를 위한) 정당성 이론은 불가피한 결함을 수용할
수 있도록 일정한 유연성을 지녀야 한다. 그것은 완전한 공정성이 아니라
합리적인 공정성(reasonable fairness)에 대해 논의해야 한다. 물론 일부
선거 및 입법 제도는 이러한 관대한 기준조차 충족하지 못할 수 있다. 그러나
우리가 논의하는 핵심 사례(core case)는 입법 및 선거 제도가 병적으로 혹은
교정 불가능할 정도로 기능 장애를 일으키는 상황을 다루기 위한 것이
아니다.

이제 다시 우리의 핵심 사례와, 완강한 시민 Cn과의 가상적 대면 상황으로
돌아가 보자. 위에서 서술한 바와 같은 내용이, 비교적 잘 조직된 입법부의
결정에 대해 Cn이 제기한 불만에 대한 응답으로 제시될 수 있다는 점은
정당성 측면에서 중요하다. 하지만 이것이 결정적인 것은 아니다. 왜냐하면
Cn은 입법 절차보다 더 정당한 다른 절차를 상정할 수도 있기 때문이다.
\emph{정당성(legitimacy)}은 부분적으로 \emph{비교적(comparative)}인
개념이다.\footnote{Waldron, supra note 47 참조.} 상이한 제도와 절차가
서로 다른 결과를 낳을 수 있기 때문에, 특정 제도나 절차의 정당성을
옹호하려면 그것이 사용 가능했던 다른 제도나 절차보다 더 공정했거나 더
공정할 것임을 보여주어야 한다.\footnote{Michelman, supra note 76, at
  57--59 참조.}

이제 우리는 다음과 같은 장면을 상상할 수 있다---혹은 미국과 같은
체제에서는 실제로 목격하게 된다: 결정이 입법부가 아니라 법원(이를테면
미국 \textbf{(p.~1390)} 대법원)에서 이루어지며, 그 결정은 권리에 대한
시민들 사이의 \emph{의견불일치(disagreement)}가 존재하는 중대한 사안에
관한 것이다. 그리고 한 시민---다시 Cn이라 부르자---이 그 판결의 실질적
내용에 동의하지 않으며 불만을 제기한다. 그녀는 다음과 같이 묻는다: (1)
왜 이 아홉 명의 남녀(대법관들)가 이 사안을 결정해야 하는가? 그리고 (2)
그들이 결정한다고 하더라도, 왜 그들이 사용하는 현재의 절차를 통해 결정을
내려야 하는가? 왜 내가 지지하는 견해를 가진 대법관들에게 더 많은 비중이
주어지는 다른 절차가 아닌가?

이 질문들은 입법자들에게 제기된 질문보다 훨씬 더 대답하기 어렵다. 우리는
이런 도전들이 종종 법원 외부에서 시끄럽게 제기된다는 점, 그리고
대법관들이 이로 인해 때때로 심리적 불편을 겪는다는 점을 신뢰할 만한
권위자로부터 잘 알고 있다. 그러나 그들 중 일부는 그러한 불편함을
성찰하기도 한다. (이제 눈을 굴리고 잠시 주의를 내려놓아도 좋다. 왜냐하면
내가 지금부터 안토닌 스캘리아(Justice Antonin Scalia)의 인용을 길게
소개할 것이기 때문이다.)

\begin{quote}
실제로, 나는 법원이 느끼는 바로 그 ``정치적 압력''에 대해 법원만큼이나
심각한 불편함을 느낀다. 즉, 우리의 견해를 바꾸도록 유도하기 위해 시위,
항의서한, 그리고 항의 시위가 법원을 향해 쏟아지는 현실 말이다. 이 얼마나
당혹스러운 일인가. 너무도 많은 우리 시민들---선량하고 합법적인 사람들로,
이 낙태 문제를 포함한 여러 사안들에서 다양한 입장을 가진 이들---이
우리가 어떤 객관적 법률을 파악하는 것이 아니라 일종의 사회적 합의를
결정하고 있다고 생각하며, 따라서 대법관들이 자신들의 견해를 고려해야
마땅하다고 생각한다는 점 말이다. 나는 법원이 이 불편한 현상의 ``사실''
자체보다, 그것의 ``원인''에 더 많은 관심을 기울였으면 한다. 그 원인은
오늘의 판결문 전체에 스며들어 있다. 즉, 헌법적 판단의 새로운
방식---문언(text)이나 전통적 관행이 아닌, 법원이 말하는 ``이성적
판단(reasoned judgment)''에 의존하는 방식인데, 결국 이는 철학적 성향과
도덕적 직관에 불과하다.\footnote{Planned Parenthood of Se. Pa. v. Casey,
  505 U.S. 833, 999--1000 (1992) (Scalia, J., 반대의견) (인용 생략).}
\end{quote}

스캘리아 대법관은 이어 다음과 같이 말했다:

\begin{quote}
이 모든 것이 법원에 대한 ``정치적 압력''이 불편하게 적용되는 문제와
관련되는 까닭은 두 가지 사실 때문이다. 첫째, 미국 국민은 민주주의를
사랑하며, 둘째, 미국 국민은 어리석지 않다. 이 법원이 우리 대법관들이
여기서 주로 변호사로서의 업무---문언을 해석하고 그에 대한 우리 사회의
전통적 이해를 밝히는 작업---를 수행한다고 믿었고, 국민들 또한 그렇게
믿었을 때, 대중은 우리를 거의 내버려두었다. 문언과 전통은 연구해야 할
사실이지, 시위해야 할 신념이 아니기 때문이다. 그러나 헌법적 판단의 실제
과정이 주로 가치판단(value judgments)을 내리는 일이라면\ldots{} 그렇다면
\textbf{(p.~1391)} 자유롭고 지성적인 국민의 태도는 달라질 수밖에 없으며,
그렇게 되어야만 한다. 국민은 그들의 가치판단이 법대에서 배운 그 어떤
것과도 동등하거나 어쩌면 더 낫다고 알고 있다. 만약 헌법이 보호하는
``자유(liberties)''가 법원이 말하듯 정의되지 않았고 경계도 없다면,
국민은 우리가 우리의 가치가 아니라 그들의 가치를 구현하지 않는다는
이유로 시위를 하고 항의해야 한다.\footnote{Ibid., 1000--01.}
\end{quote}

따라서 스캘리아(Scalia)가 말했듯이, 정당성(legitimacy)에 관한 질문은
논의의 중심에 위치하며, 사법심사(judicial review)의 옹호자들은 이에 대한
응답을 도출해내야만 한다.

첫 번째로, 왜 이 대법관들(Justices), 그리고 이 대법관들만이 해당 사안을
결정해야 하는가? 하나의 응답은 다음과 같을 수 있다. 대법관들은 일정한
선출 정당성(elective credentials)을 지닌 결정자들과 결정기관들(대통령과
상원)에 의해 지명되고 승인된다. 대통령은 선출된 인물이며, 그가 대법원에
어떤 유형의 인물을 임명할 가능성이 있는지는 유권자들이 대개 인지하고
있다. 또한 임명 승인을 담당하는 미합중국 상원의원들 역시 선출된
인물들이며, 그들 또한 이 문제에 대한 입장이 잘 알려져 있을 수 있다. 물론
대법관들은 입법자들처럼 정기적으로 책임을 지는 구조에 놓여 있지는 않다.
그러나 이미 언급했듯이, 우리는 완전함(perfection)을 요구하는 것이
아니다.

따라서, 시민의 도전에 대해 사법심사의 옹호자가 전혀 무력한 것은 아니다.
할 말은 있다. 그러나 만약 정당성이 비교적인(comparative) 개념이라면,
위의 응답은 실로 놀라울 정도로 불충분하다. 입법 선거 제도 역시 완전하지
않지만, 민주주의와 민주적 가치(democratic values)의 관점에서 보았을 때,
사법부가 지닌 간접적이고 제한된 민주적 정당성에 비해 명백히 우월하다.
입법자들은 자신들의 유권자들에 대해 정기적으로 책임을 지며, 정치적
의사결정에 참여하는 데 있어 자신들의 선출 정당성이 중요하다는 점을
실질적으로 고려하며 행동한다. 반면 이러한 점들은 대법관들에 대해서는
전혀 해당되지 않는다.

두 번째로, 설사 이러한 중대한 \emph{권리(right)} 문제들이 이 아홉 명의
남성과 여성에 의해 결정되어야 한다는 점을 인정한다 하더라도, 왜 그
결정이 단순 다수결(simple majority voting)을 통해 이루어져야 하는가? 이
지점에서 사법심사의 옹호자들에게 상황은 더욱 불리해진다. 나는 항상
법원이 그들의 소수 인원에 대해 다수결 원칙(MD: majority decision)을
적용해 표결을 통해 결정을 내리는 사실에 흥미를 느껴왔다. 물론 그들은
판결문(reasoning)을 작성하고, 위에서 논의한 바와 같은 방식으로 논증을
제시한다. 그러나 결국 결론은 숫자 싸움이다: 미국 대법원에서는 다섯 표가
네 표를 이긴다. 대법관들이 어떤 논증을 구성했는지와 무관하게 말이다. 이
문맥에서 다수결 원칙(MD)에 대한 도전이 제기된다면, 우리는 입법부를
변호할 때 상정했던 방식과 유사하게 이를 방어할 수 있을까? 사실상 그렇지
않다. 우리는 그렇게 할 수 없다. 다수결 원칙(MD)은 \textbf{(p.~1392)}
일정한 의사결정 절차에서 동등하게 간주되어야 할 \emph{도덕적
청구권(moral claim)}을 지닌 개인들에게 적절한 것이다. 그러나 대법관들에
대해서는 그러한 청구권의 도덕적 근거를 발견할 수 없다. 그들은 누구도
대의하지 않는다. 그들의 참여에 대한 \emph{청구권(claim)}은
기능적(functional)일 뿐, 자격(entitlement)에 근거하지 않는다.

나는 여기서, 법원이 다수결 원칙(MD)을 사용하는 것에 대해 이론적으로
진지하게 논의한 학문적 연구가 거의 전무하다는 점 때문에 제약을
받는다.\footnote{이 주제에 대해서는 Waldron, Deliberation, Disagreement,
  and Voting, supra note 14, at 215--24에서 간략히 다룬 바 있다.}
학자들은 법원 내 표결과 표결 전략에 대한 경험적 논의는 어느 정도
해왔으며, 일부는 개별 사안의 전체 결론이 아닌 특정 쟁점별로 판사들의
표를 조합하는 새로운 방식을 제안하기도 했다.\footnote{관련 논의는 다음
  문헌을 보라: Lewis A. Kornhauser \& Lawrence G. Sager, The One and the
  Many: Adjudication in Collegial Courts, 81 Cal. L. Rev.~1 (1993);
  Lewis A. Kornhauser \& Lawrence G. Sager, Unpacking the Court, 96 Yale
  L.J. 82 (1986); David Post \& Steven C. Salop, Rowing Against the
  Tidewater: A Theory of Voting by Multijudge Panels, 80 Geo. L.J. 743
  (1992).} 그러나 나는 사법적 다수주의(judicial majoritarianism)를 위한
기본적인 정당화 논의를 알지 못한다.\footnote{사법심사를 옹호하는 이들은
  권리에 관한 사안에서 재판관들 사이의 단순다수결 투표 문제를 꺼내는
  것을 꺼린다. 그들은 다수결이 입법부의 특징이라고 비판하길 원하기
  때문이다. 물론 그들 역시 대법원에서 판결이 5대4, 6대3 등으로
  결정된다는 사실을 인정하지만, 이를 스스로 논의에 포함시키거나 왜
  그것이 정당한 방식인지 설명하려는 경우는 매우 드물다.} 일반적으로
사용되는 공정성과 평등에 기반한 방어 논거는 여기서는 적용될 수 없다. 내
생각에는, 법원이 다수결 원칙(MD)을 사용하는 것을 정당화하려 한다면,
그것은 이 절차가 단지 이론적 함의를 지니지 않는 단순한 기술적 결정
수단이라는 점을 강조하거나,\footnote{Hannah Arendt, On Revolution 163
  (photo, reprint 1982) (1963): ``다수결의 원리는 결정 절차 그 자체에
  내재하며, 거의 자동적으로 모든 유형의 심의기관과 의회에서 채택된다''고
  진술함.} 아니면 콩도르세의 배심 정리(Condorcet's jury theorem)에
근거해 옹호할 것이다. 이 정리는 다수결이 개별 구성원의 평균 능력보다
집단의 판단 역량을 수학적으로 향상시킨다고 주장한다.\footnote{Marquis de
  Condorcet, Essay on the Application of Mathematics to the Theory of
  Decision Making (1785), 재수록: Condorcet: Selected Writings 33 (Keith
  Michael Baker ed.~\& trans., 1976).} 그러나 이 경우, 다수결에 대한
옹호는 사법적 능력(judicial competence)에 대한 \emph{결과 기반
논거(outcome-related case)}의 일부가 되며, 이는 입법부라는 훨씬 더 큰
표결 집단에 대해 제기할 수 있는 유사한 논거와 경쟁해야 할
것이다.\footnote{콩도르세 정리(Condorcet theorem)는 투표 집단의 규모가
  클수록, 개별 유권자의 평균 역량이 50\%를 초과한다는 조건 하에서,
  다수결이 개별 유권자 평균을 넘는 집단 역량을 보장할 가능성이 커진다는
  것이다. 이 조건의 적용에 대한 콩도르세 자신의 회의는 Jeremy Waldron,
  Democratic Theory and the Public Interest: Condorcet and Rousseau
  Revisited, 83 Am. Pol. Sci. Rev.~1317, 1322 (1989)에서 논의된다.} 이
논변이 어떻게 전개되든 간에, 내가 \textbf{(p.~1393)} 말하려는 요지는
다음과 같다: 입법부에 대해 제시될 수 있는 것과 같은 공정성 논거는 법원에
대한 다수결 원칙 사용에는 적용되지 않는다.

이 마지막 논점들은 우리가 QC(비판적 시민)가 입법 절차와 사법 절차 모두에
대해 제기한 도전에 대해 상정해온 응답들이 독립적으로 존재하는 것이
아님을 상기시켜준다. 우리는 또한 결과 기반 논거(outcome-related case)를
통해 그녀의 도전에 대응할 수 있다. 그러나 나는 이번 장과 그 이전 장에서,
결과 기반 논거는 결정적인 것이 아니며(혹은 입법부에 유리한 옹호 논증을
제기하며), 절차 기반 논거는 거의 전적으로 입법부 측에 유리하다는 점을
보여줄 수 있었다고 생각한다. 제3장에서 말했던 바도 기억하라. 양측의
논거는 모두 \emph{권리(right)}에 관한 것이다. 만약 한 기관이 다른
기관보다 사람들이 실제로 어떤 \emph{권리(right)}를 가지고 있는지를 더
정확히 판단할 수 있다면, 이는 그 기관에게 결정적으로 유리한 근거가 될
것이다. 그러나 현실은 그렇지 않다. 절차 측면에서 볼 때, 이러한 문제들에
대한 최종 권한을 판사들에게 부여하는 제도는 평범한 시민이 제기하는
공정성에 대한 불만에 대해, 정치적 평등(political equality)의
\emph{원칙(principle)}---단순한 \emph{가치(value)}가 아니라---에 기초한
어떠한 적절한 응답도 제공하지 못한다. 만약 사법적 결정에 대해 결과
기반의 설득력 있는 논거가 존재한다면, 이러한 실패는 감수할 수 있을지도
모른다. 사법심사의 옹호자들은 그러한 논거가 존재한다고 주장한다. 그러나
앞서 살펴본 바와 같이, 그것은 단지 근거 없는 주장일 뿐이다.

아마 이러한 점들을 인식하고 있기 때문에, 사법심사의 옹호자들은 자신들이
선호하는 제도를 민주적 가치(democratic values)와 조화시키기 위해 마지막
시도의 일환으로 몇 가지 논변들을 시도해왔다. 나는 그것들을 간략히 살펴볼
것이다. 왜냐하면 그리 중요하거나 강력한 논거들이 아니기 때문이다.

첫째, 사법심사(judicial review)의 옹호자들은 판사들이
\emph{권리(right)}에 대한 자신들의 결정을 내리는 것이 아니라, 단지
국민의 결정이 구현된 \emph{권리장전(Bill of Rights)}을 집행할 뿐이라고
주장한다. 이 권리장전은 입법이나 헌법의 일부로서 일정한 민주적
신용(democratic credentials)을 갖춘 것이라는 주장이다. 그러나 이와 같은
주장은 사법심사에 반대하는 핵심 논거(core case)를 무너뜨리지 않는다.
우리는 \emph{권리장전(Bill of Rights)}이 개인과 소수자
\emph{권리(rights)}에 대해 사회 내부에 존재하는
\emph{의견불일치(disagreements)}를 해결해 주지 않는다는 전제를 두고
있다. 권리장전은 이러한 문제들에 영향을 미치기는 하지만, 그것들을
해결하지는 않는다. 기껏해야 권리장전의 추상적인 조항들은 이러한 문제들에
대한 쟁점이 전개되는 장소로 대중적으로 선택된 것일 뿐이다. 우리가
검토해온 쟁점은, 이러한 쟁점들이 다투어지는 그 장소에서 누가 그 문제들을
최종적으로 해결할 것인가에 관한 것이다.

둘째, 이와 유사한 취지로, 사법심사의 옹호자들은 판사들이 단지 사회가
특정한 \emph{권리(right)} 원칙들에 대해 미리 맺어둔
자기구속(precommitment)을 집행하는 것이라고 주장한다. 사회는 특정한
\emph{권리(right)}의 원칙에 대해 자신의 몸을 돛대에 묶어둔 것이며,
판사들은 마치 율리시스의 선원들처럼 밧줄이 풀리지 않도록 보장하는 역할을
한다는 것이다. 이와 같은 익숙한 비유는 학계에서 이미 철저히 비판된 바
있다.\footnote{Jon Elster, Ulysses Unbound 88--96 (2000) 참조. 이 책은
  Elster의 이전 저작인 Ulysses and the Sirens: Studies in Rationality
  and Irrationality 93 (1984)에서 주장된 일부 논거에 대해 의문을 제기함.
  또한 Waldron, supra note 14, at 255--81 참조.} 간단히 말해,
\textbf{(p.~1394)} 사회가 어떤 \emph{권리(right)}가 무엇을 수반하는지에
대해 특정한 견해에 스스로를 구속시켰다고 말할 수는 없다. 그러므로 시민들
사이에서 이에 대한 \emph{의견불일치(disagreement)}가 발생할 경우,
판사들에게 그 문제를 결정할 권한을 부여하는 것이 어째서 기존의
자기구속을 존중하는 행위로 간주되어야 하는지 불분명하다. 누군가가 사회가
특정한 \emph{권리(right)}에 관해 특정한 견해에 자신을 구속했다고
주장한다 하더라도, 그리고 판사들이 표결을 통해 그 자기구속을 somehow
밝혀낸다고 하더라도, 그 \emph{권리(right)}에 대한 대안적 이해가
현실적으로 제시되는 순간, 왜 그 기존의 자기구속(precommitment)이 그대로
유지되어야 하는지는 여전히 설명되지 않는다. 율리시스 모델은 이러한
자기구속이 일탈(aberration)에 대한 방어일 경우에만 작동한다. 그러나
그것이 어떤 결과가 정당한지에 대한 진지한
\emph{의견불일치(disagreement)} 속에서 마음을 바꾸는 것을 방어하는
수단일 경우에는 적절하지 않다.\footnote{Waldron, supra note 14, at
  266--70 참조.}

셋째, 사법심사의 옹호자들은 만약 입법자들이 어떤 판결에서
\emph{권리(right)}에 대한 판단에 동의하지 않는다면, 권리장전을 개정하여
그 판결을 명시적으로 무력화하는 캠페인을 벌일 수 있다고 주장한다. 그들이
그렇게 하지 않았다는 사실은 묵시적인 민주적 승인(tacit democratic
endorsement)을 의미한다는 것이다. 그러나 이 주장은 사법적 결정이
형성하는 기준선(baseline)을 정당화하지 못한다는 점에서 결함이 있다.
\emph{권리장전(Bill of Rights)}의 개정은 일반적으로
초다수제(supermajority)를 필요로 한다. 또는 그것이 영국이나 뉴질랜드식의
법률 형식일 경우, 해당 법률이 정치문화적 정당성을 갖추고 있어 개정
시도에 상당한 부담을 부과하게 된다. 만약 불만을 가진 시민 Cn이 왜 그
기준선이 이처럼 불균형하게 설정되어야 하느냐고 묻는다면, 우리가 줄 수
있는 유일한 응답은 사법적 결정에 근거한 것이며, 이는 이미 충분치
못하다는 평가를 받았다.

넷째, 사법심사의 옹호자들은 판사들 또한 민주적 신용(democratic
credentials)을 지닌다고 주장한다. 그들은 선출된 공직자들에 의해 지명되고
승인되며, 오늘날에는 정치 후보자가 어떤 종류의 판사 지명을 할 것인지를
유권자들이 중요하게 고려하는 경우가 많아, 판사 지명은 선거 캠페인에서도
주요한 역할을 한다는 것이다.\footnote{Eisgruber, supra note 13, at 4
  (``비록 대법관들이 직접 선거로 선출되지는 않지만, 그들은 정치적·민주적
  절차를 통해 임명된다. 대통령이 지명하고 상원이 승인한다는 점에서,
  이들은 정치적 견해와 정치적 연줄에 기반해 임명되며, 이는 법적 전문성과
  마찬가지로 중요하다.'')} 이는 사실이다. 그러나 (이미 언급했듯이)
문제는 비교적(comparative) 성격을 띠고 있으며, 이러한 정당성은 선출된
입법자들이 지닌 민주적 신용과는 도저히 경쟁이 되지 않는다. 게다가,
판사들을 이와 같은 민주적 신용 때문에 수용한다면, 사법심사가 독자적으로
가치 있는 정치적 결정 형태라는 긍정적 논거 자체를 약화시키는 결과가
된다.

다섯째이자 마지막으로, 사법심사의 옹호자들은 이 제도가 시민들이 정치
체제에 접근할 수 있는 추가적인 수단으로 정당화될 수 있다고
\textbf{(p.~1395)} 주장한다. 시민들은 때로 유권자로, 때로 로비스트로,
때로 소송 당사자(litigant)로 정치 시스템에 접근한다는 것이다. 이들은
특정한 요소의 민주적 신용만을 따지기보다는, 시민 참여의 다양한 방식들이
구성하는 전체적 체계의 정당성을 평가해야 한다고 말한다. 이 논점은 일정
부분 타당하다. 그러나 사법심사를 시민 참여의 더 넓은 체계 안에
포함시킨다고 해서, 이 방식이 정치적 평등(political equality)의 원칙에
의해 규율되지 않는 형태의 참여라는 사실이 달라지지는 않는다. 사람들은
자신들의 의견이 선거 정치에서 부여받는 무게보다 더 큰 무게를 갖길 원할
때 사법심사를 찾는 경향이 있다. 물론 다른 정치적 변화 경로가 차단되어
있을 경우, 이 방식의 접근이 일정한 정당성을 얻을 수도 있다.\footnote{John
  Hart Ely, Democracy and Distrust: A Theory of Judicial Review (1980).}
이에 대해서는 제7장에서 논의할 예정이다. 그러나 우리가 논의하고 있는
핵심 사례(core case), 즉 입법부와 선거제도가 민주적 가치들에 비추어 보아
비교적 잘 작동하고 있는 상황에서는, 사법심사가 동료 시민들에 대해
전달하는 태도는 결코 정당화되기 어렵다.

\subsection{VI. 다수의 폭정(The Tyranny of the
Majority)}\label{vi.-uxb2e4uxc218uxc758-uxd3eduxc815the-tyranny-of-the-majority}

나는 사법심사(judicial review)를 옹호하는 이들에게---우리가 설정한 핵심
사례(core case)를 전제로---마지막 기회를 주고자 한다. 민주적
입법기관(democratic legislature)의 작동에 대해 가장 흔히 제기되는
우려는, 다수결주의(majoritarianism)를 바탕으로 조직된 입법 절차가
``다수의 폭정(tyranny of the majority)''을 구현하게 될 수 있다는 점이다.
이 공포는 매우 널리 퍼져 있으며, 우리 정치문화의 익숙한 요소가 되었고,
``다수의 폭정''이라는 표현은 우리의 입에서 너무도 자연스럽게
흘러나온다.\footnote{John Stuart Mill은 Toqueville의 Democracy in
  America에 대한 유일한 비판으로 ``다수의 폭정(tyranny of the
  majority)''이라는 표현이 보수 세력에게 진보적 입법을 반대하는 수사적
  도구를 제공할 것이라고 지적하였다. John Stuart Mill, M. de Tocqueville
  on Democracy in America, Edinburgh Rev., Oct.~1840, 재수록: 2
  Dissertations and Discussions: Political, Philosophical, and
  Historical 1, 79--81 (photo, reprint 1973) (1859).} 그 결과, 입법
결정을 사법적으로 감시하는 제약을 설정해야 한다는 주장은 거의 자명한
것으로 여겨지게 되었다. 다수의 폭정으로부터 소수자(minorities)를 지켜줄
수 있는 다른 안전장치는 무엇이 있는가?

나는 이 흔한 주장이 심각하게 혼동된 것이라 본다. 일단 다음과 같이 전제해
보자. 즉, \emph{권리(right)}가 부정될 때 누군가에게 '폭정(tyranny)'이
발생한다고 가정하자. 이 정의에 따르면, 폭정은 거의 모든
\emph{권리(right)}에 관한 \emph{의견불일치(disagreement)}에서 언제나
쟁점이 될 수밖에 없다. 어떤 \emph{권리(right)}에 대해 보다 넓은 이해를
지지하는 쪽(혹은 상대가 인정하지 않는 \textbf{(p.~1396)}
\emph{권리(right)}를 주장하는 쪽)은 상대방의 입장이 잠재적으로
폭압적이라고 생각할 것이다. 예를 들어, 페요테를 종교의식에 사용하는
이들은, 그들의 성례가 일반적인 마약 규제법에 의해 제한되는 것을 폭정이라
여길 것이다. 선거 자금법(campaign finance laws)에 반대하는 이들은 그
법을 폭압적이라 생각할 것이다. 하지만 그들이 옳은지는 여전히 개방된
문제다. 이 중 일부는 실제로 옳을 수도 있다. 그러나 단지 그렇게
주장된다는 사실만으로 그것이 옳아지는 것은 아니다. 실제로 어떤
\emph{권리(right)} 문제에서는 양측이 모두 상대방의 입장을 폭정이라
주장하기도 한다. 낙태 \emph{권리(rights)}를 옹호하는 이들은, 생명 보호
입장이 여성에게 폭정이라고 생각하지만, 반대로 생명 보호론자들은 선택의
자유를 옹호하는 입장이 태아(fetuses)를 향한 폭정이라 주장한다(그들에
따르면 태아는 인격자(persons)이다). 어떤 이들은 적극적
우대조치(affirmative action)가 폭압적이라고 생각하며, 또 어떤 이들은 그
조치를 시행하지 않는 것이 오히려 폭정이라 생각한다. 이러한 사례는 무수히
많다.

우리는 3장에서 월하임(Wollheim)의 역설을 논할 때도 인정했듯이, 민주적
제도들은 때때로 잘못된 \emph{권리(right)} 판단을 내리고 이를 강제할 수
있다. 이는 곧 민주적 제도도 때때로 폭압적으로 작동할 수 있음을 뜻한다.
하지만 이 점은 어떤 결정 절차에도 동일하게 적용된다. 법원 역시 때로는
폭정적으로 판단할 수 있다.\footnote{여기서 말하는 것은 입법부의 권리
  침해를 막지 못한 '부작위의 죄(sins of omission)'가 아니라, 사법심사
  권한이 오히려 입법부가 실제로 인정하려 했던 권리(rights)를 막아선
  '작위의 죄(sins of commission)'이다. 예시는 supra note 4 및 해당 본문
  참조.} 우리가 지금 전제하고 있는 정의에 따르면, 폭정은 거의 피할 수
없는 것이다. 문제는 단지 얼마나 많은 폭정이 발생할 가능성이 있는가일
뿐이며, 이것은 우리가 4장에서 다룬 문제이다.

그렇다면 어떤 정치적 결정이 다수에 의해 부과될 때, 그 폭정성이
심화되는가? 법원이 역시 다수결로 판결을 내릴 수 있다는 (형식적인) 지적은
일단 제쳐두자. 그러나 민중 다수(예: 선출된 대표자들의 다수, 그리고 이
대표자들이 각각 자신의 유권자 다수의 지지를 받고 있는 경우)에 의한
폭정은 특별히 더 심각한 폭정인가? 나는 그렇게 볼 수 없다고 생각한다.
우리는 폭정이란 누가 어떤 방식으로 결정을 내리든 그 자체로 폭정이라고 할
수도 있고, 혹은---내 입장은 이쪽이다---다수의 측면이 오히려 폭정의
정도를 약화시킨다고 말할 수도 있다. 왜냐하면 적어도 그 결정은 어떤
사람들을 정치적 평등의 입장에서 참여할 자격 없이 배제하는 방식으로
만들어진 것이 아니기 때문이다.

이 주장이 지나치게 가볍다고 느껴질 수 있으니, 좀 더 신중하게 설명해
보겠다. 입법 권한이 무제한일 경우 가장 자주 제기되는 우려는, 소수자나
개인들이 다수에 의해 억압당할 수 있다는 것이다. 그들은 억압받거나
차별당하거나, 다수에 비해 \emph{권리(rights)}가 부정되거나 침해되거나,
그들의 이익이 부당하게 다수자의 이익에 종속될 수 있다(예:
\textbf{(p.~1397)} 정의에 반할 정도로 해를 입거나 방치됨). 이러한 폭정,
억압, 불의의 양상들을 설명할 때 우리는 ``다수(majority)''와
``소수(minority)''라는 용어를 사용한다. 그러나 이 맥락에서 이 용어들은
반드시 정치적 결정 절차와 관련된 용어는 아니다. 다음과 같이 구분해보자.

불의(injustice)는 소수자의 \emph{권리(right)} 또는 이익이 다수자의 것에
부당하게 종속될 때 발생한다. 우리는 다수결 정치적 결정이 이러한 결과를
초래할 수 있다는 점을 인정했다. 그러나 그런 경우, 우리는 최소한 ``결정의
다수(decisional majority)''와 ``결정의 소수(decisional minority)''를,
그리고 ``주제적 다수(topical majority)''와 ``주제적 소수(topical
minority)''를 구분해야 한다.\footnote{여기서 ``주제적 소수(topical
  minority)''라는 표현을 사용하는 것은, 해당 결정의 주제가 되는 권리와
  이익을 지닌 집단을 의미하기 때문이다. 이 용어는 다소 느슨하게
  정의되며, 그 구성원을 둘러싼 논쟁도 있을 수 있다(``주제적 다수(topical
  majority)''도 마찬가지다). 그러나 그 느슨함은 문제가 되지 않는다. 이
  구분은, 권리에 관한 사안에서 패배한 표를 던졌다고 해서 모두가 해당
  권리에 부정적 영향을 받은 소수에 속한다고는 할 수 없다는 점을 명확히
  해준다. Waldron, supra note 14, at 13--14; Waldron, Precommitment and
  Disagreement, supra note 19; Waldron, Rights and Majorities, supra
  note 19, at 64--66 참조.} 여기서 ``주제적 다수/소수''란, 그 결정에서
\emph{권리(rights)}가 직접적으로 문제된 집단을 말한다. 어떤 경우에는
결정의 다수와 주제적 다수가 일치하고, 결정의 소수와 주제적 소수가
일치하기도 한다. 예를 들어, 인종적 불의의 경우 그러한 경향이 강하다.
백인 입법자들(결정의 다수)이 백인 특권(주제적 다수)을 위한 표결을 하고,
흑인 입법자들(결정의 소수)은 흑인의 평등 \emph{권리(rights)} 확보에
실패한다. 나는 이러한 경우야말로 ``다수의 폭정''이라는 표현이 진정으로
문제 삼아야 할 사례라고 본다.

이 구분을 바탕으로 다시 \emph{권리(rights)}에 관한
\emph{의견불일치(disagreement)} 사례로 돌아가 보자. 어떤 사회 내부에
주제적 소수(topical minority)의 \emph{권리(rights)}에 관해 의견불일치가
있다고 하자. 이 의견불일치를 해결해야 한다면, 사회는 이에 대해 숙고하고
기존의 결정 절차를 적용해야 한다. 이 사회가 다수결(MD)을 사용한다고
하자. 내가 이 결정 과정에 참여했고, 내가 지지한 입장이 패배했다면 나는
그 사안에서 결정의 소수(decisional minority)가 된다. 그러나 이것만으로
나에게 무언가 폭정적인 일이 일어났다고 말할 수는 없다. 그것이 폭정이라는
것을 보이기 위해서는 추가로 두 가지를 더 입증해야 한다: (1) 그 결정이
실제로 잘못되었고, 그 영향이 관련된 사람들의 \emph{권리(rights)}에 대해
폭정적인 것이었다는 점, (2) 내가 그러한 잘못된 결정에 의해 영향을 받은
주제적 소수(topical minority)에 속했다는 점.

\textbf{(p.~1398)} 여기서 기억해야 할 요점은 다음과 같다. 내가 속한
공동체가 내 의견을 채택하지 않았다는 사실만으로 나에게 폭정적인 일이
발생했다고 할 수는 없다. 공동체가 채택한 의견이 나의 이익을 타인의
이익과 함께 적절히 고려한 것이라면, 내 의견이 채택되지 않았다는 사실은
나의 \emph{권리(right)}, 자유(freedom), 복지(well-being)에 대한 위협이
되지 않는다. 내가 주제적 소수에 속한 경우에도 이 점은 반드시 달라지지
않는다. 사람들은---주제적 소수에 속한 사람들을 포함하여---자신이 갖고
있다고 생각하는 \emph{권리(rights)}를 실제로는 갖고 있지 않을 수도 있다.
그들이 착오를 가졌을 수도 있고, 오히려 다수의 입장이 옳을 수도 있다.
``다수의 폭정(tyranny of the majority)''에 대해 책임 있는 논의를 하고자
한다면, 이러한 분석적 구분들을 반드시 염두에 두어야 한다.

요약하자면, 다수의 폭정은 발생할 수 있다. 그러나 이 용어는 단순히
다수결의 결과에 대한 화자의 \emph{의견불일치(disagreement)}를 나타내는
방식으로 사용되어서는 안 된다. 다수의 폭정을 가장 생산적으로 개념화하는
방식은, 그것이 결정의 소수(decisional minority)와 주제적 소수(topical
minority)가 일치할 때 발생한다고 이해하는 것이다. 7장에서는 이 개념이
``명확하고 고립된 소수(discrete and insular minorities)''라 불리는
집단에 어떻게 적용될 수 있는지를 논할 것이다.\footnote{본문 중 note
  137--141을 동반하는 본문 참조.} 그러나 지금 단계에서 우리가 주목할
점은, 우리가 전제로 삼고 있는 핵심 가정(core assumptions) 하에서는 이와
같은 정렬(alignment)이 발생하지 않을 것이라는 사실이다. 세 번째 가정은,
대부분의 사람들이---즉 어떤 결정의 다수에 속한 사람들
대부분이---\emph{권리(rights)}에 대해 결정의 소수에 속한 사람들과
마찬가지로 관심을 갖고 있다는 것이었다. 네 번째 가정은,
\emph{의견불일치(disagreement)}가 대개 이기심에 의해 유발되지 않는다는
것이었다. \emph{의견불일치(disagreement)}는 쟁점의 복잡성과
난해함만으로도 충분히 설명될 수 있다. 롤즈(John Rawls)가 ``판단의
부담(the burdens of judgment)''\footnote{Rawls, Political Liberalism,
  supra note 53, at 54--58.}이라 부른 개념은, 우리가 본 것처럼, 다수의
폭정에 대한 책임 있는 논의가 전제로 삼아야 할 그러한
정렬(alignment)---즉 의견과 이익이 일치한다는 전제---를 오히려 거부하는
논증을 제공한다.

결론은, 그렇다고 해서 다수의 폭정(tyranny of the majority)이 우리가
걱정할 필요 없는 문제라는 것이 아니다. 오히려 결론은, 만약 이 용어가
책임 있게 사용된다면, \emph{다수의 폭정}은 핵심 사례(core case)가 아닌
경우들의 특징이라는 것이다. 즉, 사람들이 자신이나 자기 집단을 제외한
소수자나 개인의 \emph{권리(rights)}에 거의 관심을 기울이지 않는 상황을
의미한다. 나는 그러한 상황이 실제로 발생한다는 점을 부정하려는 것이
아니다. 다만, 내가 강조하고 싶은 것은, 그것이 이 논문에서 전제한 세 번째
가정(third assumption)과 양립하지 않는다는 점이다. 따라서 우리는
\emph{권리(rights)}에 헌신된 사회에 대해 말하면서 동시에 다수의 폭정이
상존하는 가능성이라고 주장하려 해서는 안 된다.

이 장에서 설정한 구분들은 사법심사(judicial review)에 대한 두 가지 다른
주장에 대해서도 대응하는 데 도움을 준다. 첫째, 로널드 드워킨(Ronald
Dworkin)은 『자유의 법(Freedom's Law)』에서 \textbf{(p.~1399)} 다음과
같이 논증한다. 즉, \emph{권리(rights)}가 존중되지 않는 한, 민주적
결정(democratic decisionmaking)은 본질적으로 폭정적이라는 것이다. 그는
이것이 단지 폭정적인 결과를 낳을 수 있기 때문만이 아니라,
\emph{권리(rights)}에 대한 존중이 어떤 정치적 결정 체계의 정당성의 배경
조건(background condition)이라는 점에서 그러하다고 논증한다. 드워킨은
단지 민주주의가 투표권(right to vote), 혹은 간접적으로는 언론의
자유(freedom of speech)나 결사의 자유(freedom of association) 같은
일정한 \emph{권리(rights)}에 본질적으로 의존한다는 익숙한 주장을 하고
있는 것이 아니다. 그의 주장은 그보다 더 정교하다. 그는 다수결(MD)과 같은
절차는, 참여하는 각각의 투표자들이 서로를 동등한 배려와 존중(equal
concern and respect)으로 대한다고 확신할 수 있는 경우가 아니면, 민주적
맥락에서도, 그 밖의 어떤 맥락에서도 정당성을 갖지 못한다고 주장한다.
테러리스트 집단이 다수결로 나의 운명을 결정한다면---심지어 그 다수결
과정에 나에게도 투표권이 주어진다 해도---그 결정은 전혀 정당하지 않다는
것이다. 왜냐하면 그 배경 조건이 충족되지 않기 때문이다. 드워킨의 논증에
따르면, 일반적으로 어떤 개인이 공동체의 다른 구성원들이 자신이나 자신과
같은 유형의 사람에 대해 존중하거나 신중하게 고려하지 않는다는 것을 알게
된다면, 그가 그러한 공동체 내의 다수결 결정을 정당하다고 수용하기를
기대하는 것은 거의 불가능하다는 것이다.\footnote{Dworkin, supra note 10,
  at 25.}

드워킨은 이러한 주장을 통해 사법심사에 대한 민주주의적 반박(democratic
objection)을 반박하고자 한다.\footnote{Dworkin은 이 주장이
  사법심사(judicial review)를 옹호하는 논변은 아님을 명확히 밝히고 있다.
  Id. at 7 (``민주주의는 판사들이 최종 결정을 내려야 한다고 주장하지도
  않지만, 그들이 최종 결정을 내려서는 안 된다고 주장하지도 않는다.'').}
가령 어떤 법률이 선출된 의회에 의해 제정되었고, 한 시민이 그 법률이
\emph{민주적 정당성(democratic legitimacy)}의 조건이 되는 어떤
\emph{권리(right) R}을 침해한다고 주장하며 이의를 제기한다고 하자. 다른
사람들은 이와 \emph{의견불일치(disagreement)}를 가질 것이다. 어떤 사람은
R이 민주주의의 조건이 아니라고 생각하고, 어떤 사람은 R을 전혀 다른
방식으로 이해할 수 있다. 이 문제가 법원에 최종 판단을 위한 사안으로
배정되었고, 법원이 해당 법률을 위헌으로 판단하며 시민의 이의를
수용했다고 하자. 이것이 민주주의에 손해를 주는 일일까? 드워킨의 대답은
명확하다. 이는 전적으로 법원이 올바른 결정을 내렸는지 여부에 달려 있다.
즉, 그 법률이 정말로 다수결(MD)의 정당한 적용을 위한
\emph{권리(rights)}와 양립할 수 없는 것이라면, 법원이 그 법률을 무효화한
것은 오히려 민주주의를 향상시킨 것이 된다. 그 결과 공동체는 해당 법률이
존속되었을 경우보다 더 민주적으로 정당한 상태가 된 것이다.\footnote{Ibid.
  at 32--33 (''{[}법원의 결정이 잘못되었다고 가정할 경우{]}, 위와 같은
  논지는 전혀 성립하지 않는다. 권위 있는 법원이 민주주의의 요건에 대해
  잘못된 결정을 내렸다면, 이는 민주주의를 훼손한다. 그러나 그것은 다수결
  입법기관이 잘못된 헌법적 결정을 내리고 그것이 그대로 유지되는 경우와
  전혀 다르지 않다. 오류의 가능성은 대칭적이다.'').}

\textbf{(p.~1400)} 그러나 이 주장에는 많은 문제가 있으며, 그 중 일부는
내가 다른 곳에서 지적한 바 있다.\footnote{이에 대한 전체적인 반론은
  Waldron, supra note 14, at 282--312 참조.} 우선, 드워킨은 정치적
결정이 민주주의와 관련된 경우에는 그 결정이 어떤 제도적 절차를 통해
이루어졌는지에 대한 흥미로운 논의는 성립하지 않는다고 보는 듯하다. 나는
이것이 잘못되었다고 생각한다. 예컨대 다수결 절차 자체(또는 그 정당성의
조건)에 관한 결정이 여성의 참여를 배제하는 방식으로 내려졌다면, 그
절차에 대한 평등 기반의 이의 제기는 단지 다수결의 정당성이 쟁점이라는
이유만으로 배제되어서는 안 된다. 우리는 절차 자체가 쟁점일 때에도 절차적
가치(process values)를 중요하게 여긴다.

그러나 이 주장에 대한 가장 결정적인 반박은 다음과 같다. 드워킨의 전제를
수용하자. 즉, \emph{민주적 절차(democratic procedures)}는 오직
구성원들이 서로의 \emph{권리(rights)}를 존중하는 경우에만 정당하다는
것이다. 이 전제는 두 가지 방식으로 해석될 수 있다: (1) \emph{민주적
절차}는 오직 서로의 \emph{권리(rights)}에 대한 올바른 견해를 가지고 이를
실천하는 사람들 사이에서만 정당하다는 해석, 혹은 (2) \emph{민주적
절차}는 서로의 \emph{권리(rights)}를 진지하게 받아들이고 그것이 무엇인지
최선을 다해 알아내려는 진지한 시도를 하는 사람들 사이에서 정당하다는
해석이다. 첫 번째 해석은 지나치게 강하다. 상상 가능한 어떤 정치체제도
이를 만족시키지 못한다. 그러나 두 번째 해석에는 아무런 이의가 없다.
그리고 우리가 이 두 번째 해석을 수용한다면, 이 논문에서 가정하는 사회는
드워킨의 전제를 만족시키는 것이다. 사람들이 \emph{권리(rights)}에 관해
\emph{의견불일치(disagreement)}를 가질 수는 있다. 하지만 그들은 서로의
\emph{권리(rights)}를 진지하게 받아들일 수 있다. 결정의 다수(decisional
majority)는 때로 옳고 때로는 틀릴 수 있다. 하지만 그것은 모든 결정
체계가 공유하는 문제이며, 주제적 소수(topical minority)가 자신의
\emph{권리(rights)}가 공동체 구성원 대다수에게 진지하게 고려되고 있다는
확신을 가질 수 있다면, 그 체계의 정당성을 무너뜨리는 근거가 될 수 없다.

둘째, 이 장에서 논의한 구분은 또 다른 주장을 평가하는 데에도 유용하다.
즉, 사법심사가 존재하지 않는 상태에서 \emph{권리(rights)}에 대해
입법기관이 최종 결정을 내리는 것은, 다수가 자기 사건의 재판관이 되는
것과 같다는 주장이다. 이때 인용되는 격언이 ``nemo iudex in sua
causa''(자기 사건의 재판관이 되어서는 안 된다)이다. 이 격언을 적용하는
사람들은, \emph{권리(rights)}에 관한 최종 판단은 국민의 손에 맡겨서는 안
되며, 오히려 법원과 같은 독립적이고 공정한 제도에 위임되어야 한다고
주장한다.

그러나 이 주장이 실제로 갖는 힘은 그다지 크지 않다. 어떤 결정 규칙이든
결국 누군가는 자신의 사건에 대해 결정을 내리는 위치에 서게 될 것이다.
우리가 무한한 항소 절차를 상상하지 않는 한, 항상 어떤 사람이나 기관이
최종 결정을 내리게 된다. 그리고 그 사람이나 기관에 대해서는,
\textbf{(p.~1401)} 그들이 최종 결정을 내리는 순간 자동적으로 자신들의
견해의 정당성을 스스로 판단하게 되는 것이라 말할 수밖에 없다. 이러한
방식으로 ``nemo iudex in sua causa''를 기계적으로 인용하는 것은 정당성의
기본 논리를 망각하는 변명이 될 수 없다. 사람들은
\emph{의견불일치(disagreement)}를 가지고 있고, 그렇기에 우리는 최종적인
결정과 그 결정을 내리는 절차(final decision procedure)를 필요로 한다.

이 두 번째 주장---즉 사법심사의 필요성에 대한 주장---이 의미하는 바는,
주제적 다수(topical majority), 즉 자신의 \emph{권리(rights)}나 이익이
쟁점이 된 집단이 그러한 \emph{권리(rights)}나 이익을 유지할지 여부를
결정하는 데 결정적 표를 행사해서는 안 된다는 뜻일 수 있다. 그리고 주제적
다수와 결정의 다수(decisional majority)가 일치하는 경우에는 우려할 만한
타당한 이유가 존재한다. (이러한 일치가 만연하다면, 나는 그것이 핵심
사례가 아닌 상황이라고 판단하며, 그 이유는 7장에서 설명할 것이다.)
하지만 이런 경우는 실제로 드물며, 특히 미국에서 사법심사가 다루는
일반적인 사건들에서는 더욱 그러하다. 앞서 언급한 두 가지 예를 떠올려
보자: 낙태와 적극적 우대조치(affirmative action). 어느 경우에도 우려할
만한 정렬(alignment)은 보이지 않는다. 많은 여성이 낙태
\emph{권리(rights)}를 지지하지만, 많은 남성도 이를 지지하며, 또 많은
여성은 이에 반대한다. 많은 아프리카계 미국인이 적극적 우대조치를
지지하지만, 많은 백인도 지지하며, 또 많은 아프리카계 미국인은 반대한다.
이는 우리가 2장에서 설정한 세 번째, 네 번째 가정이 충족된 사회에서
기대할 수 있는 모습이다. \emph{권리(rights)}를 진지하게 여기는 사람들은
이에 관해 \emph{의견불일치(disagreement)}를 가질 수 있다. 그러나 바로 그
진지함이 이 \emph{의견불일치(disagreement)}들이 개인의 이해관계와
독립적으로 나타난다는 것을 보여주는 지표가 된다.

\subsection{VII. 핵심 사례가 아닌 경우 (Non-Core
Cases)}\label{vii.-uxd575uxc2ec-uxc0acuxb840uxac00-uxc544uxb2cc-uxacbduxc6b0-non-core-cases}

지금까지 내가 제시한 논증은 네 가지 매우 엄격한 전제에 기반하고 있다.
그렇다면 이러한 전제가 실패할 경우, 혹은 해당 전제가 성립하지 않는
사회에서는 이 논증은 어떤 의미를 가지는가? 특히 내가 염두에 두고 있는
것은 첫 번째 전제와 세 번째 전제이다. 즉, 첫 번째 전제는 정치적
평등(political equality)의 관점에서 보았을 때 민주적이고 입법적 제도들이
제대로 작동하고 있다는 것이며, 세 번째 전제는 우리가 논의하는 사회의
구성원들이 대체로 개인과 소수자의 \emph{권리(rights)}에 헌신하고 있다는
것이다. 많은 이들에게 있어 사법심사(judicial review)에 대한 주장은
이러한 전제를 받아들이지 않는 데서 출발한다. 사법심사는 민주적 제도의
실패에 대한 반응이거나, 많은 사람들이 \emph{권리(rights)}를 충분히
진지하게 받아들이지 않기 때문에 (그들을 대신하여) 법원이 이를 대신해야
한다는 생각의 산물이다. 요컨대, 사법심사 옹호자들은 현실 세계(real
world)에서 입법에 대한 사법심사가 필요하다고 주장한다. 이는 내가 설정한
전제가 규정한 이상적 세계(ideal world)와는 다른 것이다.

\textbf{(p.~1402)} 이에 대해 몇 가지 응답을 제시할 수 있으며, 이후에
핵심 사례가 아닌 경우(non-core cases)에 대한 구체적 쟁점들을 살펴볼
것이다. 첫째, 내가 사용한 전제들이 비현실적인 것은 아니다. 예를 들어 세
번째 전제---사회 내의 전반적 \emph{권리(rights)}에 대한 헌신---는 비교적
쉽게 만족된다. 왜냐하면 사법심사에 대한 논의는 거의 항상, 사법심사가
도입될 사회가 시민들의 견해와 일정한 관련성을 지닌 권리장전(Bill of
Rights)을 갖고 있다는 점을 전제로 하기 때문이다. 첫 번째 전제는 선거 및
입법 절차가 상당히 잘 작동하고 있다는 것이며, 정치적 평등(political
equality)의 이름으로 완전무결함을 요구할 권리가 우리에게 있는 것도
아니다. 또한, 제5장에서 입법부와 법원의 정당성에 대해 논의할 때 나는
다시금 강조하였다. 나의 논증은 시민 개개인이 \emph{목소리가 들릴 동등한
청구권(the equal claim to be heard of the voices)}을 요구하는 데 완벽히
응답해야 한다는 전제에 의존하지 않는다. 입법적 결정의 정당성을 옹호하는
주장은, 입법 제도가 유토피아적으로 완벽하다는 전제에 의존하지 않으며,
선거 절차와 제도적 구조가 정치적 평등의 원칙을 완벽하게 구현하고 있다는
전제에도 의존하지 않는다. 대신 그것은, 해당 제도가 정치적 평등이라는
원칙에 명시적으로 지향되고 있으며, 그 원칙을 충족시키도록 조직되어 있고,
이를 실현하기 위한 합리적 노력을 기울이고 있다는 점에 기반한다.
마지막으로, 나는 철학적 선험적 구성물이 아닌 실제 입법 절차의 사례---즉
1966년 영국 하원의 '임신중절 의료법안(Medical Termination of Pregnancy
Bill)'에 관한 토론---를 인용하여, 입법부가 어떻게 작동할 수 있는지를
보여주었다.

이 모든 말을 인정하더라도, 우리는 여전히 다음과 같은 질문을 던져야 한다:
이러한 전제가 성립하지 않을 경우, 사법심사에 반대하는 논증은 어떻게
되는가?

만약 전제가 성립하지 않는 경우라면, 이 논문에서 제시한 사법심사 반대
논증은 성립하지 않는다. 제2장에서 강조한 바와 같이, 나의 논증은
조건부(conditonal) 논증이다.\footnote{note 43과 관련된 본문을 참조하라.
  이 논변이 실패하는 사례로는 Waldron, supra note 47 참조.} 그러나
그렇다고 해서 전제가 무너질 경우 항상 사법심사가 정당화된다는 결론이
나오는 것도 아니다. 나의 전제에 의존하지 않는 다른 유효한 사법심사 반대
논증이 존재할 수 있다. 또는 사법심사가 특정 상황을 개선하는 데 아무런
희망도 제공하지 못할 수도 있다. 만약 법원 또한 입법부 못지않게
부패하거나 편견에 물들어 있다면, 그러한 사회에서 입법에 대한 사법심사를
설정하는 것은 적절하지 않을 수 있다. 핵심 사례에 대해 우리가 논의했던
많은 논증이 비교적(comparative) 논증이었던 만큼, 이 논리는 핵심 사례가
아닌 경우에도 동일하게 적용된다.

예를 들어, 나의 첫 번째 전제가 실패한 경우를 가정해보자. 즉, 입법부가
대표성과 숙의성을 결여하고 있으며, 선거 제도가 훼손되었고,
\textbf{(p.~1403)} 입법부에서 사용되는 절차들이 더 이상 정치적 정당성과
신뢰할 만한 관계를 맺고 있지 않다. 그렇다면 두 가지 질문이 제기된다. (1)
입법부와 관련된 상황을 개선할 수 있는가? (2) 중요한
\emph{권리(rights)}에 관한 쟁점에 대해 최종 결정 권한을 법원에 부여해야
하는가, 만약 법원이 그 문제를 더 잘 다룰 수 있다면? 이 두 질문은
독립적인 질문이다. 왜냐하면 우리는 어떤 \emph{권리(rights)}의 문제들은
보다 책임 있는 입법부가 나타날 때까지 기다릴 수 없을 정도로 시급하다고
생각할 수도 있기 때문이다. 그러나 이 두 질문은 완전히 분리되어 있지는
않다. 법원에 최종 결정 권한을 부여하는 것은 입법부의 개혁을 더 어렵게
만들 수 있으며, 첫 번째 전제와 아마도 세 번째 전제가 전제하는 입법적
규범(legislative ethos)의 형성을 방해할 수 있다. 미국의 상황에 대해 이와
같은 추측을 들어본 적이 있다. 예를 들어, 미국---특히 주의회---의
입법부가 무책임하게 작동하고 있으며, \emph{권리(rights)}를 진지하게
다루지 않는 방식으로 운영되는 이유는, 법원이 뒤에 버티고 있다는 인식이
책임 있는 입법 문화의 발전을 어렵게 만들기 때문이라는 것이다. 이 말이
실제로 어느 정도 사실인지는 모르겠다. 그러나 충분히 고려해볼 만한
주장이다.

마지막으로, 나의 첫 번째 전제가 실패할 수 있는 잘 알려진 사례 하나를
살펴보며 이 장을 마무리하고자 한다. 내가 염두에 두고 있는 것은 스톤
대법관(Justice Stone)이 유명한 Carolene Products 판결의 각주 4(footnote
four)에서 제안한 내용이다. ``{[}D{]}iscrete and insular minorities에
대한 편견은, 소수자를 보호하는 데 일반적으로 의존되는 정치적 절차의
작동을 심각하게 저해하는 특별한 조건이 될 수 있다\ldots{}''\footnote{United
  States v. Carolene Products Co., 304 U.S. 144, 153 n.4 (1938); 또한
  Keith E. Whittington, An ``Indispensable Feature''? Constitutionalism
  and Judicial Review, 6 N.Y.U. J. Legis. \& Pub. Pol'y 21, 31 (2002)
  참조 (Whittington은 내가 『Law and Disagreement』에서 이 아이디어를
  다루지 않은 것이 ``미국 헌법 이론의 관점에서 보았을 때 주목할
  만하다''고 말한다).} 나의 판단으로는, 이는 입법 결정에 대한 사법심사가
일정한 타당성을 갖는 핵심 사례가 아닌 경우를 묘사하는 훌륭한 방식이다.
이러한 상황에 놓인 소수자들은 비선출적(non-elective) 제도만이 제공할 수
있는 특별한 보호---그들의 \emph{권리(rights)}를 보호하고, 존 하트
일라이(John Hart Ely)가 지적한 바와 같이, 정치체제를 복원하며 대표성을
가능케 하는 방식으로---가 필요할 수 있다.\footnote{Ely, supra note 126,
  at 135--79.}

그러나 우리는 이 논증을 다룰 때 신중해야 한다. 제6장에서 논의한 바에
따르면, 모든 소수자가 이러한 특별한 대우를 받아야 하는 것은 아니다.
결정의 소수(decisional minority) 전체도, 주제적 소수(topical minority)
전체도 예외가 아니다.\footnote{Tushnet, supra note 11, at 159 (``모든
  법은 패배한 소수의 견해를 무시한다. 여기서 우리는 단지 패배한
  사람들과, 정치 과정에서 자신을 보호할 수 없기 때문에 패배한 소수
  집단을 구별해야 한다.'').} 심지어 만성적인 소수(chronic minority)라
하더라도, 그들이 항상 이러한 특별 대우를 받아야 할 \textbf{(p.~1404)}
이유는 없다. 특히 만성적인 결정의 소수---예컨대 미국 내
볼셰비키(Bolsheviks)---는 그렇지 않다.

우리는 너무 자주 ``discrete and insular''라는 표현을 무비판적으로
사용한다. 모든 뚜렷하고 식별 가능한 소수자가 'discrete and insular'한
것은 아니다. 스톤 대법관의 표현에 어떤 마법적 힘이 있는 것도 아니다.
그러나 이 표현이 진지하게 받아들여진다면, ``discrete''와 ``insular''는
유용한 형용사들이다. 이들은 단지 정치적 결정 과정에서 분리되어 존재하는
소수자---즉 주제적 소수(topical minority)---를 뜻하는 것이 아니다. 그
구성원들이 공동체의 나머지 구성원들과 여러 공통의 이익을 공유하지
못하여, 자신들의 이익을 위해 연합(coalition)을 구축하는 것이 불가능한
유형의 고립된(minimally connected) 소수자를 의미한다. 우리가 제6장에서
경고했던 결정의 소수와 주제적 소수가 일치하는 상황은 이러한 의미에서
'고립(insularity)'의 좋은 예이며, 그것은 분명한 우려의 대상이다.

그렇다면 스톤 대법관(Justice Stone)이 언급한 다른 기준---즉, 해당 소수자
집단이 \emph{편견(prejudice)}의 피해자라는 기준---은 어떠한가? 만연한
편견은 분명히 나의 세 번째와 네 번째 전제와 양립할 수 없는 것이다.
그것은 해당 집단 구성원의 \emph{권리(rights)}에 대한 무관심 혹은
적대감을 내포하며, 대다수 집단의 구성원들이 소수 집단 구성원들의 자기
인식과 비합리적인 방식으로 다르게 판단하도록 이끌 수 있다. 그러나
``편견(prejudice)''이라는 용어는 지나치게 협소하며, 하나의 집단과 다른
집단 사이에 고착화되고 무의식적인 반감을 형성하는 깊이를 포착하지 못할
수 있다.\footnote{Charles R. Lawrence III, The Id, the Ego, and Equal
  Protection: Reckoning with Unconscious Racism, 39 Stan. L. Rev.~317
  (1987) 참조. 이 점을 강조해준 Ian Haney-Lopez에게 감사한다.} 여기서
핵심은 반감의 특정 양태를 고집하는 것이 아니라, 그것이 다양한 양태를
지니고 있으며, 정당한 \emph{권리(rights)}에 대한
\emph{의견불일치(disagreement)}와 구별되어야 한다는 점이다.\footnote{편견(prejudice)과
  종교적 또는 윤리적 근거에 기반한 강한 신념은 구별되어야 한다. 우리는
  생명권(pro-life)을 지지하는 이들의 견해가 단지 우리가 그들의 종교적
  신념을 공유하지 않는다는 이유로 편견이라고 간주해서는 안 된다. 권리에
  관한 거의 모든 견해―선택권(pro-choice)을 지지하는 견해를 포함하여―는
  강한 감정적 기반을 가지며, 궁극적으로는 확고하고 깊이 뿌리내린 가치
  신념에 기초한다.}

이러한 경우에는, 내가 앞서 전개한 사법심사에 반대하는 핵심 논증(core
argument)은 유지될 수 없다. 그러나 이 역시 곧바로 사법심사를 찬성하는
논거가 성립된다는 의미는 아니다. 모든 것은 사법부의 다수(decisional
majorities)가 입법부의 다수와 동일한 \emph{편견(prejudice)}에 물들어
있는지 여부에 달려 있다. 만약 동일하다면, 이는 단지 핵심 사례(core
case)가 아닌 것이 아니라, 아예 희망이 없는(non-viable) 경우이다. 해당
사회 전체에서 소수자의 \emph{권리(rights)}에 대한 어떠한 지지도 존재하지
않는다면, 사법심사의 실천은 소수자의 \emph{권리(rights)}에 아무런 도움도
줄 수 없다. 이러한 상황에서 종종 제기되는 사법심사 옹호 논증은, 해당
소수자의 \emph{권리(rights)}에 대한 일정한 존중이 소수자 집단 내부
구성원을 넘어 사회 내의 정치 엘리트(political elites) 사이에
제한적으로나마 존재한다는 점을 전제로 한다. \textbf{(p.~1405)}
가정하건대, 일반 대중 대다수는 그러한 동정(sympathy)을 공유하지 않는다.
이러한 동정을 공유하는 엘리트 집단 구성원들(이를 나는 \emph{엘리트
동조자(elite sympathizers)}라 부르겠다)은 입법부에 있을 수도 있고,
사법부에 있을 수도 있다. 사법부에 최종 결정 권한을 부여하자는 논증은
다음과 같다: 사법부 내 \emph{엘리트 동조자(elite sympathizers)}들이,
선출된 입법부 내의 \emph{엘리트 동조자}들보다, 비인기 소수 집단의
구성원들에게 \emph{권리(rights)}를 부여할 때 스스로를 더 잘 보호할 수
있다는 것이다. 그들은 대중의 분노에 덜 노출되며, 보복에 대해서도 걱정할
필요가 없다. 그러므로 이들은 소수자를 더 잘 보호할 가능성이 높다.

이러한 사법심사 옹호 논증이 소수자의 \emph{권리(rights)}에 대한 지지
분포(distribution of support)에 관한 특정한 가정에 의존하고 있다는 점에
주목하라. 이 논증은 해당 동정(sympathy)이 정치 엘리트 내에서 가장 강하게
존재한다고 가정한다. 만약 이 가정이 거짓이라면---즉, 동정이 일반 대중
사이에서 더 강하게 존재한다면---앞 단락의 논증을 받아들일 이유는 없다.
반대로, 선출 기관(elective institutions)이 소수자의
\emph{권리(rights)}를 보호하는 데 더 효과적일 수 있으며, 선거 제도는
소수자의 \emph{권리(rights)}에 대한 대중적 지지를 입법부로 연결할 수
있는 통로를 제공하는 반면, 사법부에는 그러한 통로가 존재하지 않는다.
물론 소수자 \emph{권리(rights)}에 대한 지지의 분포는 사안마다 다를 수
있다. 그러나 내가 흥미롭게 여기는 점은 다음과 같다. 사법심사를 옹호하는
다수의 이론가들이, 만약 해당 사회 내에 소수자의 \emph{권리(rights)}에
대한 일정한 지지가 존재한다고 가정할 경우, 그 지지가 어디에 있든 간에
반드시 엘리트층에 존재할 것이라고 확신한다는 점이다. 이들은 이것을
실증적 주장(empirical claim)으로 방어하려 하겠지만, 나는 그것이 민주적
결정(democratic decisionmaking)에 대한 오래된 편견과도 정확히 조응한다고
말하고 싶다.

또 하나 고려해야 할 요소는, 사법심사의 확립된 관행이 장기적으로 선출 및
입법 기능의 실패를 바로잡는 데 더 용이하게 작용할지, 혹은 오히려 더
어렵게 만들지를 따져보는 일이다. 특정한 상황에서는, \emph{discrete and
insular minorities}(분리되고 고립된 소수자들)이 그들의
\emph{권리(rights)}를 보호받기 위해 사법적 개입(judicial
intervention)으로부터 이익을 얻을 수 있다. 그러나 제도적 차원에서는,
사법적 배려(judicial solicitude)가 상황을 더 악화시키거나, 최소한 큰
개선 효과를 내지 못할 수도 있다. 미국이 1950년대와 1960년대에 경험한
바와 같이, Brown 사건 및 기타 판례들에서의 사법적 인종차별 철폐 노력에
열광했음에도 불구하고, 궁극적으로 필요한 것은 강력한 입법 개입---즉,
민권법(Civil Rights Act) 형태의 개입---이었다. 그리고 결정적인 차이를
만든 것은 사법부(courts)와 입법부(legislatures) 사이의 구별 자체가
아니라, 연방(federal) 기관과 주(state) 기관 사이의 구분이었다. 결국에는
연방 입법부가 결정적 역할을 수행했다.

전반적으로, 우리는 Carolene Products 판결의 각주나 유사한 교의를
일반적인 사법심사 관행을 \emph{레버리지(leverage)} 하기 위한 수단으로
해석해서는 \textbf{(p.~1406)} 안 된다.\footnote{사법심사를 옹호하기 위한
  논거로 이와 같은 논지의 유용성에 한계가 있다는 점에 대해서는, Tushnet,
  supra note 11, at 158--63을 참조하라. 이 부분에 대한 일반적인 논의가
  잘 정리되어 있다.} \emph{discrete and insular minorities}(분리되고
고립된 소수자들)의 문제는 사법심사의 일종의 '트로이 목마'로 간주되어서도
안 되며, 사법심사에 반대하는 논증들을 곤혹스럽게 만들기 위한 논거로
사용되어서도 안 된다. 이러한 사례를 고려하는 목적은 사법심사를
옹호하려는 것이 아니라, 해당 소수자의 \emph{권리(rights)}를 가장
효과적으로 확보할 수 있는 방안을 찾는 데 있다. 우리는 바로 그 목표를
지향해야 하며, 그 과정에서 이러한 사례들이 사법심사를 옹호하는 일반
논증의 일종의 이념적 전위(ideological vanguard)로 간주될 수 없다는 점을
명확히 인식해야 한다.

\subsection{결론(CONCLUSION)}\label{uxacb0uxb860conclusion}

나는 입법에 대한 사법심사(judicial review)의 실천이 모든 상황에서
부적절하다고 보이려는 것이 아니었다. 대신, 나는 \emph{권리(right)}에
기반한 사법심사가 다음과 같은 사회에서는 부적절하다는 점을 보이고자
했다: 입법 제도가 기능 장애(dysfunction)을 일으키는 것이 주요 문제가
아니라, 그 구성원들이 \emph{권리(rights)}에 대해
\emph{의견불일치(disagreement)}를 보이고 있는, 비교적 민주적인 사회.

\emph{권리(rights)}에 대한 \emph{의견불일치(disagreement)}는 비합리적인
것이 아니다. 사람들은 서로의 \emph{권리(rights)}를 진지하게
받아들이면서도 그에 관하여 \emph{의견불일치}할 수 있다. 이러한
상황에서는, 당사자들이 수백만 명에 이르는 구성원들의
\emph{권리(rights)}가 논쟁의 대상이 되는 가운데 그들의
\emph{의견(opinions)}과 \emph{목소리(voices)}를 존중하고, 그들을 절차
내에서 동등한 존재로 대우하는 방식으로
\emph{의견불일치(disagreements)}를 해결하기 위한 절차를 채택해야 한다.
동시에, 그 절차들은 그러한 \emph{권리에 대한
의견불일치(rights-disagreements)}가 제기하는 어렵고 복잡한 문제들을
책임감 있고 숙고적인(deliberative) 방식으로 다루어야 한다. 나는 이러한
절차적 요구를 일반적인 입법 절차가 충족할 수 있다는 점을 논증해왔다.
그리고 그 위에 법원이 최종 결정을 내릴 수 있도록 덧붙이는 사법심사의
층위는, 절차에 실질적인 기여를 하기는커녕, 오히려 경멸적인 방식으로
시민들의 \emph{발언(a say)}을 박탈하고, 우리가 \emph{권리(rights)}에
관해 벌이는 \emph{의견불일치(disagreements)} 속에 내재한 도덕적 쟁점을
법률적 언어로 모호하게 만드는 역할만을 할 뿐이다.

물론, 특이한 병리들, 기능 장애가 있는 입법 제도, 부패한 정치문화,
인종차별과 같은 구조화된 편견의 유산과 같은 상황들에서는 이러한
모호화(obfuscation)와 배제(disenfranchisement)의 대가를 일시적으로
감수할 가치가 있을지도 모른다. 그러나 사법심사를 옹호하는 이들이라면,
이러한 실천을 정당화하려는 주장을 할 때, 그 기반이 되는 상황을 솔직하게
인정하고, 그 상황에 대해 일정한 겸허함(humility)과 부끄러움(shame)을
갖고 논증을 전개해야 할 것이다. 그것을 마치 \emph{권리(rights)}에 대한
존중의 전형(epitome)인 양, 현대 헌정 민주주의에서 보편적이고 정당한
제도인 양 외부에 전파하려 해서는 안 된다.

\clearpage
\backmatter

\end{document}